\chapter{Tropicalization of analytic spaces}
\index{tropicalization!analytic space}\index{moment map}
\label{chap:analytic-tropicalization}

% The global operator \trop denotes the tropicalization map; the canonical
% section below is chapter-local.  Upright operator names avoid confusion
% with the polyhedral image, which is written \Trop.
\newcommand{\sq}{\operatorname{sq}}

We now explain how the real calculus enters a Berkovich space. Four notions
must remain distinct. A moment map supplies analytic functions; its
tropicalization supplies real coordinates; a tropical chart specifies where
those coordinates are used; and a calibration records the multiplicities
needed for integration.

Throughout, $K$ is a complete non-Archimedean field and $X$ is a good,
boundaryless $K$-analytic space of pure dimension $n$, unless explicitly
stated otherwise.

It is worth saying at the outset what problem this chapter solves.
\Cref{chap:superforms,chap:tropical-spaces,chap:tropical-integration} built a
complete differential calculus on the support of a balanced weighted
polyhedral complex: superforms, the operators $\dprime$ and $\ddprime$,
calibrated integration, and Stokes' formula. None of it applies to a
Berkovich space directly, because a Berkovich space is not a real manifold
and carries no atlas of charts modelled on $\mathbf R^m$. What it does carry
is an enormous supply of invertible analytic functions, and the map
\[
 x\longmapsto\bigl(-\log|f_1(x)|,\ldots,-\log|f_r(x)|\bigr)
\]
is continuous into $\mathbf R^r$ for any such family $f_1,\ldots,f_r$. The
image is locally rational polyhedral; on a compact analytic domain it is a
finite rational polyhedral set, not necessarily a single polyhedron.  So
Berkovich spaces do not have real charts, but they have real \emph{shadows},
and there are enough shadows to see everything: the whole difficulty is that
a shadow is not injective, that different families of functions cast
different shadows, and that a
$\mathbf R$-linear functional pulled back along a shadow remembers only part
of the analytic geometry above it.

The resolution has three layers, and this chapter is organized around them.

First, one identifies, inside $X$, the locus on which a given moment map
$f$ has \emph{maximal} tropical rank. This is the local skeleton, or
characteristic polyhedron, $\Sigma_f$. On it, and only on it, the
tropicalization map is a piecewise immersion of polyhedral spaces, so it is
the correct domain of integration. Its intrinsic characterization
(\Cref{thm:trop-sigma-characterization}) is purely graded-residue-theoretic:
$x\in\Sigma_f$ if and only if $n$ of the initial forms
$\widetilde{f_i(x)}$ are algebraically independent over $\widetilde K$.

Second, one attaches to each facet of $\Sigma_f$ a positive number, the
local analytic degree of a suitable finite flat projection. This is the
calibration. Its construction involves an auxiliary torus projection $q$,
and the whole content of \Cref{thm:trop-calibration-canonical} is that the
answer does not depend on $q$: the degree and the lattice normalization of
the affine section move by reciprocal factors.

Third, one shows that the resulting calibration on a local skeleton is
independent of the moment map used to see the skeleton
(\Cref{thm:trop-calibration-chart-independent}). This is the technical heart
of the chapter. Two moment maps have no common refinement in any naive
sense; what they have is a \emph{product}, and the argument is that both
calibrations agree with the product calibration cell by cell.

Once these three layers are in place, the chapter's final structural
statement (\Cref{cor:trop-calibration-glue}) says that every polyhedral
subspace of a pure $n$-dimensional analytic space carries a canonical
calibration, compatible across tropical charts.  The real calculus of
\Cref{chap:superforms,chap:tropical-spaces,chap:tropical-integration} can
therefore be transported chartwise to the analytic space. That is the input to
\Cref{chap:cld-forms,chap:cld-currents}.

\paragraph{Roadmap.}
\Cref{sec:trop-dictionary} fixes the sign conventions relating the max-plus
terminology of abstract tropical geometry to the valuation conventions of
this book; nothing later works without it. \Cref{sec:trop-moment} defines
moment maps and proves their elementary functoriality, together with the
polyhedrality theorem. \Cref{sec:trop-torus} computes the fundamental
example completely: the canonical skeleton $S(\mathbf G_m^n)$ of a torus, its
intrinsic characterizations, and its behaviour under monomial maps. Every
later multiplicity is measured against this example.
\Cref{sec:trop-charts} introduces tropical charts, refinement, and the
product construction. \Cref{sec:trop-local-skeleta} constructs local skeleta
$\Sigma_f$ and proves their basic properties;
\Cref{sec:trop-dimension} proves the dimension formula in the chain
\[
 \dim_{\mathrm{trop},x}(f)\le d_x(f)\le d_K(x)\le\dim_xX,
\]
identifying each inequality and each case of equality.
\Cref{sec:trop-calibration} constructs the canonical calibration and proves
it canonical; \Cref{sec:trop-refinement} proves refinement independence and
glues. \Cref{sec:trop-multiplicities} derives balancing, the multiplicity
formula for a tropical polyhedron, and independence of integration.
\Cref{sec:trop-graded} isolates the graded reduction as the source of the
multiplicities: a local degree factors as a length, a graded residue degree,
and a defect, and the last is trivial in residue characteristic zero.
\Cref{sec:trop-supports} treats weak analytic structures, boundaries, and
partitions of unity at the level that
\Cref{chap:cld-forms,chap:cld-currents,chap:cld-bedford-taylor} import.
\Cref{sec:trop-worked} carries out five complete computations.

\paragraph{Dependencies.}
The reader needs, from \Cref{chap:valuations}, graded reduction
(\Cref{lem:graded-residue-field}) and the Abhyankar inequality
(\Cref{thm:abhyankar-relative}); from \Cref{chap:berkovich}, analytification,
affinoid domains, the completed residue field $\mathcal H(x)$, the boundary
$\partial X$, and the classification of points on the analytic line
(\Cref{thm:berk-a1-classification}); from \Cref{chap:skeleta}, skeleta of SNC
models, integral affine structures (\Cref{def:integral-affine}), and
retractions; and from
\Cref{chap:tropical-spaces,chap:tropical-integration}, weighted balanced
complexes, calibrations, integer-affine pushforward
(\Cref{prop:tropical-pushforward-balanced}), and subdivision independence of
integrals (\Cref{prop:integral-subdivision-independent}). The polyhedral and
graded conventions are collected in \Cref{app:tropical-algebra}, whose
\Cref{sec:appG-analytic-multiplicity} states the multiplicity formula in the
form used here.

Conversely, this chapter is used in \Cref{chap:cld-forms} to define the
sheaves $A^{p,q}_X$, in \Cref{chap:cld-currents} to define currents and
cycle classes, in \Cref{chap:cld-poincare-lelong,chap:cld-bedford-taylor} to
localize potential theory, in \Cref{chap:delta-forms} as the base on which
delta-forms and harmonic tropicalizations are built, and in
\Cref{chap:compare-ma} for the comparison with model metrics.

\paragraph{Standing hypotheses.}
Unless a statement says otherwise, $K$ is a complete non-Archimedean field,
not assumed algebraically closed, not assumed nontrivially valued, and not
assumed of residue characteristic zero. Analytic spaces are $K$-analytic in
the sense of \Cref{chap:berkovich}, Hausdorff, and of pure dimension $n$;
following the terminology of \cite[\S12.4]{CLD25} we say that $X$ is
\emph{of level $n$} when $\dim X\le n$, since all constructions below refer
to a fixed integer $n$ and simply return nothing on components of smaller
dimension. Goodness --- the existence of affinoid neighbourhoods --- is
assumed wherever local invertible functions must be produced, and is
recalled explicitly in those statements. Boundarylessness is a genuine
hypothesis: several equalities below degrade to inequalities on
$\partial X$, and we mark each occurrence.

Five inputs are imported rather than proved, each with an exact reference
and a description of its mechanism: the polyhedrality of a compact
tropicalization (\Cref{thm:analytic-tropical-polyhedrality}); the local
structure of tropicalization germs (\Cref{thm:trop-local-germ}); the
polyhedrality of local skeleta (\Cref{thm:trop-sigma-polyhedral}); the
dimension, flatness and semicontinuity properties of analytic morphisms
collected in \Cref{prop:trop-analytic-inputs}; and the invariance of local
degrees under a tropical translation
(\Cref{thm:trop-degree-translation}). Everything else --- in particular
every independence, gluing and multiplicity statement --- is proved here.
The only place where an alternative to the curve-slope proof of balancing is
needed is recorded in \Cref{rem:trop-balancing-two-routes}.

\section{The analytic--tropical dictionary}
\label{sec:trop-dictionary}

The constructions in this chapter have two convention layers.  Abstract
tropical geometry follows Mikhalkin--Rau: max-plus arithmetic,
integer-affine maps, facets as top-dimensional cells, and tropical cycles as
equivalence classes of balanced weighted complexes
\cite[\S\S1.2--1.3, 4.2]{MR18Tropical}.  Analytic geometry follows the
valuation convention
\[
                         u(f)=-\log|f|.
\]
Thus an abstract max-plus coordinate is $x=\log|f|=-u(f)$.  If
\[
 F_{\min}(u)=\min_m\{v(a_m)+\langle m,u\rangle\},
\]
then the corresponding Mikhalkin--Rau polynomial is
\begin{equation}
\label{eq:analytic-max-min-firewall}
 F_{\max}(x)=-F_{\min}(-x)
 =\max_m\{\log|a_m|+\langle m,x\rangle\}.
\end{equation}
Coordinate reflection and value negation are both present.  Reflection
alone leaves a minimum a minimum, so we never describe it as a change of
semiring convention.

By the end of the chapter the reader should be able to distinguish:
\begin{enumerate}
\item the analytic morphism furnishing invertible functions;
\item its continuous tropicalization map;
\item its rational polyhedral image and, in full analytic rank, the
  weighted $n$-cycle carried by its $n$-dimensional part;
\item a tropical chart used to represent a differential form;
\item the calibration that records analytic multiplicity.
\end{enumerate}
These distinctions are needed even on an annulus: the map, its interval
image, the skeletal section, and the weight-one calibration are four
different objects.

We record the translation once and for all, in the form in which it is used
below. Let $N=\mathbf Z^r$ and $M=N^\vee$ with the pairing
$\langle\cdot,\cdot\rangle$. A Laurent monomial $\chi^m$ on
$\mathbf G_{m,K}^r$, $m\in M$, has
\begin{equation}
\label{eq:trop-monomial-affine}
 -\log|\chi^m(x)|=\langle m,u\rangle,
 \qquad u=\bigl(-\log|T_1(x)|,\ldots,-\log|T_r(x)|\bigr),
\end{equation}
so monomials become integral linear forms, and a Laurent polynomial
$f=\sum_ma_m\chi^m$ satisfies the ultrametric bound
\begin{equation}
\label{eq:trop-ultrametric-bound}
 -\log|f(x)|\ \ge\ \min_m\bigl\{v_K(a_m)+\langle m,u\rangle\bigr\}
 =\operatorname{trop}(f)(u),
\end{equation}
in the notation of \eqref{eq:appG-tropical-polynomial}. Three things must be
kept apart in \eqref{eq:trop-ultrametric-bound}. The inequality is always
true and is the ultrametric inequality. Equality for \emph{all} $f$
characterizes the points of the canonical skeleton
(\Cref{prop:trop-torus-skeleton}). Failure of equality for some $f$ is
exactly cancellation of leading terms, and the locus of $u$ where
cancellation is possible is the corner locus of
\Cref{def:appG-corner-locus}. The passage from the last to the
Mikhalkin--Rau max-plus corner locus is \eqref{eq:analytic-max-min-firewall}
and nothing else.

\begin{warning}{Two reflections, not one}
The change from $\min$ to $\max$ and the change from $u$ to $x=-u$ are
independent operations. Performing only one of them turns a convex function
into a concave one but leaves a minimum a minimum; performing both returns a
maximum, which is what Mikhalkin--Rau write. In particular the balanced
cycle underlying a tropical hypersurface is carried between the two
conventions by the linear involution $u\mapsto-u$ of $N_{\mathbf R}$, which
preserves lattices, primitive vectors, and lattice indices. Hence every
weight computed below is the same in both conventions; only the pictures are
reflected.
\end{warning}

A last orientation remark. In this book \emph{tropicalization} always names
a map out of an analytic space, never an abstract combinatorial object. The
abstract objects are the ones of \Cref{chap:tropical-spaces}: polyhedral
complexes, weights, balancing, calibrations. The content of the present
chapter is a functor from analytic geometry to that category, and the
theorems are about the functor being well defined.

\section{Moment maps}
\label{sec:trop-moment}

\begin{definition}
\label{def:trop-moment-map}
Let $U\subset X$ be an analytic domain. A \emph{moment map} is an analytic
morphism
\[
                         \varphi:U\longrightarrow\mathbf G_{m,K}^r.
\]
Writing $\varphi=(f_1,\ldots,f_r)$ with $f_i\in\mathcal O(U)^\times$, its
\emph{tropicalization} is
\[
\varphi_{\mathrm{trop}}:U\longrightarrow\mathbf R^r,\qquad
x\longmapsto(-\log|f_1(x)|,\ldots,-\log|f_r(x)|).
\]
Its image is denoted $\operatorname{Trop}_\varphi(U)$.
\end{definition}

The minus sign agrees with $v=-\log|\cdot|$. If
$g=c f_1^{a_1}\cdots f_r^{a_r}$, then
\[
                         -\log|g(x)|
=\sum_i a_i(-\log|f_i(x)|)-\log|c|,
\]
so monomial changes of moment map become integer-affine maps.

In the terminology fixed in \Cref{chap:tropical-spaces}, these are
integer-affine maps.  If one instead uses the Mikhalkin--Rau coordinate
$x_i=\log|f_i|$, the same formula has the max-plus sign convention.

It is convenient, and for the invariance statements essential, to phrase the
definition without a choice of coordinates on the target.

\begin{definition}[Invariant form]
\label{def:trop-moment-invariant}
Let $T$ be a $K$-split torus with character lattice $M=X^*(T)$ and
cocharacter lattice $N=M^\vee$. Set
\[
 T_{\mathrm{trop}}=\operatorname{Hom}(M,\mathbf R)=N_{\mathbf R}.
\]
A \emph{moment map} on $X$ is a morphism of $K$-analytic spaces
$f:X\to T$, and its \emph{tropicalization} is
\[
 f_{\mathrm{trop}}:X\longrightarrow T_{\mathrm{trop}},
 \qquad
 \langle m,f_{\mathrm{trop}}(x)\rangle=-\log\bigl|\chi^m(f(x))\bigr|
 \quad(m\in M).
\]
A trivialization $T\simeq\mathbf G_{m,K}^r$ identifies
$T_{\mathrm{trop}}$ with $\mathbf R^r$ and $f$ with an $r$-tuple of
invertible analytic functions, recovering the coordinate description above.
\end{definition}

That $f_{\mathrm{trop}}(x)$ is a well-defined element of
$\operatorname{Hom}(M,\mathbf R)$, and not merely a function on $M$, is
multiplicativity of the seminorm $|\cdot|_x$: for $m,m'\in M$ one has
$\chi^{m+m'}=\chi^m\chi^{m'}$, hence
$-\log|\chi^{m+m'}(f(x))|=-\log|\chi^m(f(x))|-\log|\chi^{m'}(f(x))|$.
This is the first and most elementary point at which the multiplicativity of
Berkovich seminorms --- as opposed to submultiplicativity --- is used; it is
what makes tropicalization linear rather than merely subadditive.

\begin{lemma}
\label{lem:trop-moment-basic}
Let $f:X\to T$ be a moment map.
\begin{enumerate}
\item\label{itm:trop-moment:cont} $f_{\mathrm{trop}}$ is continuous.
\item\label{itm:trop-moment:funct} If $g:Y\to X$ is a morphism of
  $K$-analytic spaces, then $f\circ g$ is a moment map on $Y$ and
  $(f\circ g)_{\mathrm{trop}}=f_{\mathrm{trop}}\circ g$.
\item\label{itm:trop-moment:torus} If $q:T\to T'$ is a homomorphism of tori,
  then $q\circ f$ is a moment map and
  $(q\circ f)_{\mathrm{trop}}=q_{\mathrm{trop}}\circ f_{\mathrm{trop}}$,
  where $q_{\mathrm{trop}}=\operatorname{Hom}(q^*,\mathbf R):N_{\mathbf R}\to
  N'_{\mathbf R}$ is the integer-linear map induced by
  $q^*:M'\to M$. A translation of $q$ by a point of $T'(K)$ changes
  $q_{\mathrm{trop}}$ by the constant vector $\Trop(t)$ of that point, so an
  arbitrary morphism of tori induces an integer-\emph{affine} map with
  integral linear part.
\item\label{itm:trop-moment:prod} If $f:X\to T$ and $f':X\to T'$ are moment
  maps, then $(f,f'):X\to T\times T'$ is a moment map with
  $(f,f')_{\mathrm{trop}}=(f_{\mathrm{trop}},f'_{\mathrm{trop}})$, and the
  two projections induce $q_{\mathrm{trop}}$, $q'_{\mathrm{trop}}$ recovering
  $f_{\mathrm{trop}}$ and $f'_{\mathrm{trop}}$.
\item\label{itm:trop-moment:base} If $L/K$ is a complete valued extension and
  $p:X_L\to X$ the base-change morphism, then
  $(f_L)_{\mathrm{trop}}=f_{\mathrm{trop}}\circ p$.
\end{enumerate}
\end{lemma}

\begin{proof}
\cref{itm:trop-moment:cont} In a trivialization, $f_{\mathrm{trop}}$ has
components $-\log|f_i|$ with $f_i\in\mathcal O(X)^\times$. The function
$x\mapsto|f_i(x)|$ is continuous by the definition of the topology on a
Berkovich space, and takes values in $(0,+\infty)$ because $f_i$ is
invertible, on which $-\log$ is continuous.

\cref{itm:trop-moment:funct} The composition of morphisms of analytic
spaces is a morphism, and $|f_i(g(y))|=|(f_i\circ g)(y)|$ by definition of
the seminorm at $g(y)$.

\cref{itm:trop-moment:torus} For $m'\in M'$ one has
$\chi^{m'}\circ q=\chi^{q^*m'}$, hence
\[
 \langle m',(q\circ f)_{\mathrm{trop}}(x)\rangle
 =-\log\bigl|\chi^{q^*m'}(f(x))\bigr|
 =\langle q^*m',f_{\mathrm{trop}}(x)\rangle
 =\langle m',q_{\mathrm{trop}}(f_{\mathrm{trop}}(x))\rangle .
\]
Since $q^*$ is a homomorphism of lattices, $q_{\mathrm{trop}}$ has integral
matrix in any pair of bases. If $q$ is composed with translation by
$t\in T'(K)$, then $\chi^{m'}$ acquires the constant factor $\chi^{m'}(t)$,
whose $-\log|\cdot|$ is $\langle m',\Trop(t)\rangle$; this adds the constant
$\Trop(t)$.

\cref{itm:trop-moment:prod} $X^*(T\times T')=M\oplus M'$, and the
displayed formula for $(f,f')_{\mathrm{trop}}$ is the definition applied to
each summand. The projections induce the two inclusions on characters,
hence the two coordinate projections on tropicalizations.

\cref{itm:trop-moment:base} The completed residue field of $x_L\in X_L$
receives $\mathcal H(p(x_L))$ isometrically, so
$|f_i(x_L)|=|f_i(p(x_L))|$.
\end{proof}

Assertion \cref{itm:trop-moment:torus} is the reason that lattice indices
appear later: an integral linear map between lattices of the same rank need
not be an isomorphism of lattices even when it is an isomorphism of real
vector spaces, and the discrepancy $[N':q(N)]$ is precisely the
normalization factor \eqref{eq:appG-lattice-index}. Assertion
\cref{itm:trop-moment:prod} is the only construction available for
comparing two unrelated moment maps; there is no notion of a common
\emph{sub}moment map, and this asymmetry shapes every refinement argument in
the chapter.

\begin{lemma}
\label{lem:trop-moment-local-existence}
Let $X$ be a $K$-analytic space, $x\in X$, and suppose $X$ is good at $x$.
Let $a_1,\ldots,a_d$ be homogeneous elements of the graded residue field
$\widetilde{\mathcal H(x)}$ that are algebraically independent over
$\widetilde K$. Then there are a compact analytic neighbourhood $V$ of $x$
and invertible analytic functions $f_1,\ldots,f_d$ on $V$ with
$\widetilde{f_i(x)}=a_i$ for every $i$. In particular $x$ lies in the local
skeleton of the moment map $f=(f_1,\ldots,f_d)$ once $d=n$.
\end{lemma}

\begin{proof}
Fix $i$ and let $r_i$ be the degree of $a_i$, so $r_i>0$. Choose
$\alpha_i\in\mathcal H(x)$ with $|\alpha_i|=r_i$ representing $a_i$.
Goodness at $x$ gives an affinoid neighbourhood $W=\mathcal M(A)$ of $x$.
By construction $\mathcal H(x)$ is the completion of
$\kappa(x)=\operatorname{Frac}(A/\mathfrak p_x)$, so $\kappa(x)$ is dense in
$\mathcal H(x)$ and there is an element of $\kappa(x)$ within distance
$<r_i$ of $\alpha_i$; writing it as a fraction of two elements of $A$ whose
denominator does not vanish at $x$, we obtain an analytic function $g_i$ on
an open neighbourhood of $x$ in $W$ with
\[
 |g_i(x)-\alpha_i|<r_i .
\]
Then $|g_i(x)|=r_i$ by the ultrametric equality case
(\Cref{lem:strong-triangle-equality}), and $\widetilde{g_i(x)}=a_i$, since
$g_i(x)$ and $\alpha_i$ differ by an element of strictly smaller absolute
value, which dies in the graded piece of degree $r_i$ by
\Cref{prop:appG-initial-calculus}.

Now fix $r_i'$ with $0<r_i'<r_i$ and set
\[
 V=\{y\in W'\mid |g_i(y)|\ge r_i'\text{ for }1\le i\le d\},
\]
where $W'\subset W$ is an affinoid neighbourhood of $x$ on which all the
$g_i$ are defined. Then $V$ is a Laurent domain in $W'$
(\Cref{ex:berk-weierstrass-laurent}), hence compact and analytic; it
contains the open set $\{|g_i|>r_i'\ \forall i\}\ni x$, hence is a
neighbourhood of $x$; and every $g_i$ is invertible on $V$, since it does
not vanish there. Set $f_i=g_i|_V$. The last assertion is
\Cref{thm:trop-sigma-characterization} below.
\end{proof}

The germ of a tropicalization near a point is more rigid than its global
image, and this local rigidity is what makes the notion of tropical
dimension well behaved. We import the statement.

\begin{theorem}
\label{thm:trop-local-germ}
Let $X$ be a $K$-analytic space, $f:X\to T$ a moment map, $x\in X$, and
$\xi=f_{\mathrm{trop}}(x)$. Let $d_x(f)$ denote the transcendence degree
over $\widetilde K$ of the graded subfield of $\widetilde{\mathcal H(x)}$
generated by $\widetilde K$ and the initial forms
$\widetilde{\chi^m(f(x))}$, $m\in M$. Then there is a compact analytic
neighbourhood $U$ of $x$ such that for every compact analytic
neighbourhood $V\subset U$ of $x$:
\begin{enumerate}
\item\label{itm:trop-germ:germ} the germs at $\xi$ of
  $f_{\mathrm{trop}}(U)$ and of $f_{\mathrm{trop}}(V)$ coincide; we write
  $f_{\mathrm{trop}}(X,x)$ for this germ;
\item\label{itm:trop-germ:bound}
  $\dim_\xi f_{\mathrm{trop}}(V)\le d_x(f)$;
\item\label{itm:trop-germ:interior} if $x\notin\partial X$, then
  $f_{\mathrm{trop}}(V)$ is purely of dimension $d_x(f)$ at $\xi$; if
  moreover $d_x(f)=\dim T$, then $f_{\mathrm{trop}}(V)$ is a neighbourhood of
  $\xi$ in $T_{\mathrm{trop}}$.
\end{enumerate}
\end{theorem}

\begin{proof}[Foundational local tropicalization input]
This is the local germ theorem of Ducros, quoted in the present form as
\cite[Proposition~12.2.8 and Remark~12.2.9]{CLD25}; the proof there reduces
the assertion to Temkin's graded reduction of the germ $(X,x)$, in which the
germ of the tropicalization is the rational cone cut out by the
monomial inequalities defining the graded reduction, and then invokes the
dimension theory of graded fields recalled in \Cref{sec:appG-graded-reduction}.
We use \cref{itm:trop-germ:germ}--\cref{itm:trop-germ:interior} as
stated and prove everything else.
\end{proof}

The number $d_x(f)$ of \Cref{thm:trop-local-germ} is the \emph{residual
transcendence degree of $f$ at $x$}. It is defined without reference to any
trivialization: a change of basis of $M$ replaces the family
$(\widetilde{f_i(x)})$ by a family of invertible monomials in it, which
generates the same graded subfield. We record this and the comparison with
the invariant $d_K(x)$ of \Cref{chap:berkovich}.

\begin{definition}
\label{def:trop-residual-degree}
For a moment map $f:X\to T$ and $x\in X$, write $F_x(f)\subset
\widetilde{\mathcal H(x)}$ for the graded subfield generated by
$\widetilde K$ and the elements $\widetilde{\chi^m(f(x))}$, $m\in M$, and
\[
 d_x(f)=\operatorname{trdeg}_{\widetilde K}F_x(f),
 \qquad
 d_K(x)=\operatorname{trdeg}_{\widetilde K}\widetilde{\mathcal H(x)} .
\]
Here transcendence degrees are taken in the graded sense of
\Cref{sec:appG-graded-reduction}; for $X$ a curve, $d_K(x)=s(x)+t(x)$ in the
notation of \Cref{prop:berk-type-invariants}.
\end{definition}

\begin{lemma}
\label{lem:trop-residual-degree-bounds}
For every moment map $f:X\to T$ and every $x\in X$,
\[
 d_x(f)\ \le\ \min\{\dim T,\ d_K(x)\},
 \qquad
 d_K(x)\ \le\ \dim_xX .
\]
Moreover $d_x(f)$ is unchanged when $f$ is replaced by $f\cdot t$ for a
$K$-point $t$ of $T$, and $d_x(q\circ f)\le d_x(f)$ for every homomorphism
of tori $q$, with equality whenever $q$ is an isogeny.
\end{lemma}

\begin{proof}
The subfield $F_x(f)$ is generated over $\widetilde K$ by the $\dim T$
elements $\widetilde{f_i(x)}$ in any trivialization, whence
$d_x(f)\le\dim T$; and it is a subfield of $\widetilde{\mathcal H(x)}$,
whence $d_x(f)\le d_K(x)$. The inequality $d_K(x)\le\dim_xX$ is the
Abhyankar inequality for the graded residue field of an analytic point, in
the graded formulation of \Cref{sec:appG-graded-reduction}; it is
\Cref{thm:abhyankar-relative} applied to an irreducible component through
$x$.

For the last two assertions: multiplying $f$ by a $K$-point multiplies each
$\chi^m(f(x))$ by a nonzero scalar of $K$, whose initial form is a unit of
$\widetilde K$, so $F_x(f)$ is unchanged. Composing with $q$ replaces the
generating family by the subfamily of monomials $\chi^{q^*m'}$, so
$F_x(q\circ f)\subset F_x(f)$ and $d_x(q\circ f)\le d_x(f)$. If $q$ is
an isogeny then $q^*$ is injective with finite cokernel, so $F_x(q\circ f)$
and $F_x(f)$ have the same algebraic closure in
$\widetilde{\mathcal H(x)}$ --- every generator of the larger has a power
lying in the smaller --- and the transcendence degrees agree.
\end{proof}

\begin{example}[The annulus]
\label{ex:trop-annulus}
Assume here that $v_K(\pi)=1$.  On
$U=\{|\pi|\leq|T|\leq1\}$ the moment map $T:U\to\mathbf G_m$ has tropical
image $[0,1]$. Its canonical section is the family of Gauss points $v_r$ with
$v_r(T)=r$. Thus tropicalization retracts the annulus onto the interval
skeleton studied in \Cref{chap:four-entrances}.
\end{example}

\begin{example}[A Laurent polynomial and the max/min firewall]
\label{ex:trop-max-min}
Let $f=a_0+a_1T+a_2S$ on the two-dimensional torus and put
$u=(-\log|T|,-\log|S|)$.  Away from cancellation, the ultrametric inequality
gives
\[
 -\log|f|
 =\min\{-\log|a_0|,\,-\log|a_1|+u_1,\,
                       -\log|a_2|+u_2\}.
\]
The tropical hypersurface is the locus where the minimum occurs at least
twice.  After $x=-u$ and multiplication by $-1$, this is the corner locus of
the max-plus polynomial used by Mikhalkin--Rau.  The underlying balanced
cycle is carried to the other convention by the reflection $u\mapsto-u$.
\end{example}

\begin{theorem}\label{thm:analytic-tropical-polyhedrality}
For a compact analytic domain $U$ and a moment map $\varphi$, the tropical
image is a finite rational polyhedral set of dimension at most $\dim U$.
\end{theorem}
\begin{proof}
This is the compact polyhedrality theorem
\cite[Theorem~12.2.4]{CLD25}.  Its substantive input is the reduction theory
of analytic germs: after passing to affinoid pieces, weighted graded
reductions describe the possible vectors
$(-\log|f_1|,\ldots,-\log|f_r|)$ by finitely many rational affine
inequalities.  The dimension bound follows from the graded Abhyankar
inequality.  Compactness turns the resulting locally finite polyhedral
structure into a finite one.  Multiplicities are deliberately not included
in this imported statement; they are constructed intrinsically from local
analytic degrees in \Cref{sec:trop-calibration,sec:trop-multiplicities}.
\end{proof}

\section{The canonical skeleton of a torus}
\label{sec:trop-torus}

Every multiplicity in this chapter is measured against one example, which we
therefore compute completely. Let $T=\mathbf G_{m,K}^n$ with coordinates
$T_1,\ldots,T_n$, character lattice $M=\mathbf Z^n$, and
$T_{\mathrm{trop}}=\mathbf R^n$. Write $u=(u_1,\ldots,u_n)$ for the
tropical coordinate and $r_i=\rme^{-u_i}$, so that $r^m=\rme^{-\langle
m,u\rangle}$ for $m\in M$.

\begin{definition}[Gauss points of a torus]
\label{def:trop-gauss-point}
\index{Gauss point!of a torus}
For $u\in\mathbf R^n$ define $\eta_u$ on the Laurent polynomial ring
$K[T_1^{\pm1},\ldots,T_n^{\pm1}]$ by
\begin{equation}
\label{eq:trop-gauss-seminorm}
 \Bigl|\sum_ma_mT^m\Bigr|_{\eta_u}
 =\max_m|a_m|\,\rme^{-\langle m,u\rangle}.
\end{equation}
\end{definition}

\begin{lemma}
\label{lem:trop-gauss-multiplicative}
Formula \eqref{eq:trop-gauss-seminorm} defines a multiplicative norm on
$K[T^{\pm1}]$ extending the absolute value of $K$; hence, by
\Cref{prop:berk-pair-description}, a point
$\eta_u\in T^{\An}=\mathbf G_{m,K}^{n,\An}$. It satisfies
$\trop(\eta_u)=u$, where $\trop=\mathrm{id}_{\mathrm{trop}}$.
\end{lemma}

\begin{proof}
The right-hand side of \eqref{eq:trop-gauss-seminorm} is a maximum of
finitely many nonnegative numbers, vanishing only if all $a_m$ vanish;
it is $\le$-subadditive for the maximum because
$|a_m+b_m|\le\max\{|a_m|,|b_m|\}$, and restricts to $|\cdot|$ on $K$ since
only $m=0$ occurs. It remains to prove multiplicativity.

Let $f=\sum_ma_mT^m$ and $g=\sum_mb_mT^m$ be nonzero, and put
$\rho=|f|_{\eta_u}$, $\sigma=|g|_{\eta_u}$. Every coefficient of $fg$ is a
finite sum $\sum_{m+m'=\ell}a_mb_{m'}$, so
$|fg|_{\eta_u}\le\rho\sigma$ by the ultrametric inequality. For the reverse
inequality, let
\[
 S=\{m:|a_m|r^m=\rho\},\qquad S'=\{m':|b_{m'}|r^{m'}=\sigma\};
\]
both are nonempty and finite. Fix the lexicographic total order on
$M=\mathbf Z^n$; it is compatible with addition. Let $m_0=\max_{\mathrm{lex}}S$
and $m_0'=\max_{\mathrm{lex}}S'$, and set $\ell=m_0+m_0'$. The coefficient
of $T^\ell$ in $fg$ is
\[
 c_\ell=a_{m_0}b_{m_0'}+\sum_{\substack{m+m'=\ell\\ (m,m')\ne(m_0,m_0')}}
 a_mb_{m'} .
\]
Consider a pair $(m,m')\ne(m_0,m_0')$ with $m+m'=\ell$. If $m\notin S$ or
$m'\notin S'$, then $|a_mb_{m'}|r^\ell<\rho\sigma$. If $m\in S$ and
$m'\in S'$, then $m\le_{\mathrm{lex}}m_0$ and $m'\le_{\mathrm{lex}}m_0'$
by maximality; since $m+m'=m_0+m_0'$ and the order is translation
invariant, both inequalities must be equalities, contradicting
$(m,m')\ne(m_0,m_0')$. Hence exactly one term of $c_\ell$ has absolute value
$\rho\sigma r^{-\ell}$ and the others are strictly smaller, so
$|c_\ell|r^\ell=\rho\sigma$ by the ultrametric equality case
(\Cref{lem:strong-triangle-equality}), and $|fg|_{\eta_u}\ge\rho\sigma$.

Finally $|T_i|_{\eta_u}=\rme^{-u_i}$, so $\trop(\eta_u)=u$.
\end{proof}

\begin{definition}
\label{def:trop-torus-skeleton}
\index{skeleton!canonical, of a torus}
The map $\sq:T_{\mathrm{trop}}\to T^{\An}$, $u\mapsto\eta_u$, is the
\emph{canonical section} of $\trop$. Its image
\[
 S(T)=\sq(T_{\mathrm{trop}})\subset T^{\An}
\]
is the \emph{canonical skeleton} of $T$.
\end{definition}

The following proposition collects the three characterizations of $S(T)$
used below: an equality of norms, an extremality property, and a
graded-reduction condition. The third is the one that generalizes, and it
is the reason graded reduction is unavoidable in this theory.

\begin{proposition}
\label{prop:trop-torus-skeleton}
Let $x\in T^{\An}=\mathbf G_{m,K}^{n,\An}$ and put
$u=\trop(x)$. The following are equivalent.
\begin{enumerate}
\item\label{itm:trop-sk:eq} For every Laurent polynomial
  $f=\sum_ma_mT^m$ one has
  \[
   |f(x)|=\max_m|a_m|\rme^{-\langle m,u\rangle};
  \]
  that is, $x=\eta_u$.
\item\label{itm:trop-sk:max} $x$ is maximal for the partial order on
  $T^{\An}$ defined by $y\le z$ if and only if $|f(y)|\le|f(z)|$ for every
  analytic function $f$ on $T^{\An}$.
\item\label{itm:trop-sk:graded} The initial forms
  $\widetilde{T_1(x)},\ldots,\widetilde{T_n(x)}$ are algebraically
  independent over $\widetilde K$ in the graded residue field
  $\widetilde{\mathcal H(x)}$; equivalently $d_x(\mathrm{id}_T)=n$.
\end{enumerate}
Moreover $\sq$ is continuous, $S(T)$ is closed in $T^{\An}$, and $\sq$ is a
homeomorphism of $T_{\mathrm{trop}}$ onto $S(T)$ inverse to
$\trop|_{S(T)}$. When $x\in S(T)$ one has $d_K(x)=n$.
\end{proposition}

\begin{proof}
\emph{\cref{itm:trop-sk:eq} $\Rightarrow$ \cref{itm:trop-sk:max}.} Suppose
$x=\eta_u$ and $\eta_u\le z$. Put $s_i=|T_i(z)|$ and $s=(s_1,\ldots,s_n)$,
$w=-\log s$. Applying $\eta_u\le z$ to $T_i$ and to $T_i^{-1}$ gives
$r_i\le s_i$ and $r_i^{-1}\le s_i^{-1}$, hence $s=r$ and $w=u$. On the other
hand, for any $y\in T^{\An}$ with $\trop(y)=u$ the ultrametric inequality
gives $|f(y)|\le\max_m|a_m|r^m=|f(\eta_u)|$ for every Laurent polynomial
$f$, and hence for every analytic function by density; so $z\le\eta_u$.
Therefore $z=\eta_u$: any $z$ dominating $\eta_u$ equals it.

\emph{\cref{itm:trop-sk:max} $\Rightarrow$ \cref{itm:trop-sk:eq}.} Let
$x$ be arbitrary and $u=\trop(x)$. The computation just made shows
$x\le\eta_u$. If $x$ is maximal, then $x=\eta_u$.

\emph{\cref{itm:trop-sk:eq} $\Rightarrow$ \cref{itm:trop-sk:graded}.} Let
$P=\sum_ma_mS^m$ be a nonzero graded polynomial relation over
$\widetilde K$ among the $\widetilde{T_i(x)}$; here the monomial $S^m$ has
degree $r^m$, and we may assume $P$ homogeneous of some degree $\rho$,
since the graded components of a relation are again relations. Homogeneity
of degree $\rho$ means that only exponents $m$ with $|a_m|r^m=\rho$ occur,
where $a_m\in K$ represents $\widetilde{a_m}$. Lifting, $f=\sum a_mT^m$
satisfies $|f|_{\eta_u}=\rho$, whereas $P(\widetilde{T(x)})=0$ says
precisely $|f(x)|<\rho$. This contradicts \cref{itm:trop-sk:eq}. Hence no
nontrivial relation exists.

\emph{\cref{itm:trop-sk:graded} $\Rightarrow$ \cref{itm:trop-sk:eq}.}
Let $f=\sum_ma_mT^m$ and let $\rho=\max_m|a_m|r^m$. Only the exponents $m$
with $|a_m|r^m=\rho$ contribute to the degree-$\rho$ part; by
\Cref{prop:appG-initial-calculus} the class of $f(x)$ in degree $\rho$ is
$\sum_{|a_m|r^m=\rho}\widetilde{a_m}\,\widetilde{T(x)}^m$, a nonzero graded
polynomial expression in algebraically independent elements, hence nonzero.
So $|f(x)|=\rho$.

\emph{Topology.} Continuity of $\sq$: for a Laurent polynomial $f$, the
function $u\mapsto|f(\eta_u)|=\max_m|a_m|\rme^{-\langle m,u\rangle}$ is
continuous on $\mathbf R^n$, and these functions generate the topology of
$T^{\An}$. Since $\trop\circ\sq=\mathrm{id}$, the map $\sq$ is injective and
$\trop|_{S(T)}$ is a continuous inverse; therefore $\sq$ is a homeomorphism
onto its image. Closedness: if $x\notin S(T)$, then by
\cref{itm:trop-sk:eq} there is a Laurent polynomial $f$ with
$|f(x)|<|f(\eta_{\trop(x)})|$; the set of $y$ with
$|f(y)|<|f(\eta_{\trop(y)})|$ is open, because $y\mapsto|f(y)|$ is
continuous and $y\mapsto|f(\eta_{\trop(y)})|$ is continuous by the previous
sentence, and it avoids $S(T)$.

\emph{The invariant $d_K(x)$.} If $x\in S(T)$ then
$d_K(x)\ge d_x(\mathrm{id}_T)=n$ by \cref{itm:trop-sk:graded}, while
$d_K(x)\le\dim_xT^{\An}=n$ by \Cref{lem:trop-residual-degree-bounds}.
\end{proof}

\begin{corollary}
\label{cor:trop-skeleton-intrinsic}
Let $T$ be any $K$-split torus of dimension $n$. The subset
$S(T)\subset T^{\An}$ obtained from \Cref{def:trop-torus-skeleton} through
one trivialization $T\simeq\mathbf G_{m,K}^n$ is independent of the
trivialization, and $\trop$ restricts to a homeomorphism
$S(T)\to T_{\mathrm{trop}}$. Transporting the group law of
$T_{\mathrm{trop}}$ makes $S(T)$ a topological group, and the composite
$T(K)\to T_{\mathrm{trop}}\xrightarrow{\ \sq\ }S(T)$ is a group
homomorphism.
\end{corollary}

\begin{proof}
Characterizations \cref{itm:trop-sk:max} and \cref{itm:trop-sk:graded} of
\Cref{prop:trop-torus-skeleton} do not mention the coordinates: the former
is stated with all analytic functions, and the latter says
$d_x(\mathrm{id}_T)=n$, which by
\Cref{lem:trop-residual-degree-bounds} is unchanged under a monomial change
of basis of $M$, since such a change replaces the generating family
$(\widetilde{T_i(x)})$ by a family of Laurent monomials in it generating the
same graded subfield. The remaining statements are immediate from the
homeomorphism $\sq$ and the identity $\trop(st)=\trop(s)+\trop(t)$, which is
multiplicativity of the absolute values of characters.
\end{proof}

The next lemma is the model computation for every tropical multiplicity in
this chapter. Note that both the analytic degree and the lattice index occur,
and that they occur on opposite sides.

\begin{lemma}
\label{lem:trop-monomial-degree}
Let $A\in M_n(\mathbf Z)$ with $\delta=|\det A|\ne0$, let
$c=(c_1,\ldots,c_n)\in(K^\times)^n$, and let
$\varphi:\mathbf G_{m,K}^n\to\mathbf G_{m,K}^n$ be the homomorphism
followed by translation given by $\varphi^*(T_i)=c_iT^{A_i}$, where $A_i$
is the $i$-th row of $A$. Then:
\begin{enumerate}
\item\label{itm:trop-mono:affine} $\varphi_{\mathrm{trop}}(u)=Au+v_K(c)$ is
  integer affine with linear part $A$, and
  $[\mathbf Z^n:A\mathbf Z^n]=\delta$;
\item\label{itm:trop-mono:finflat} $\varphi$ is finite and flat of degree
  $\delta$;
\item\label{itm:trop-mono:preimage}
  $\varphi^{-1}(S(\mathbf G_m^n))=S(\mathbf G_m^n)$, and $\varphi$ induces a
  homeomorphism of $S(\mathbf G_m^n)$ onto itself;
\item\label{itm:trop-mono:degree} for $x\in S(\mathbf G_m^n)$ the fibre
  $\varphi^{-1}(\varphi(x))$ is the single reduced point
  $\mathcal M(\mathcal H(x))$, the graded residue extension
  $\widetilde{\mathcal H(x)}/\widetilde{\mathcal H(\varphi(x))}$ has degree
  $\delta$, and $\deg_x(\varphi)=\delta$.
\end{enumerate}
\end{lemma}

\begin{proof}
\cref{itm:trop-mono:affine} is \Cref{lem:trop-moment-basic}
\cref{itm:trop-moment:torus} together with the elementary identity
$[\mathbf Z^n:A\mathbf Z^n]=|\det A|$ for an injective integral matrix,
which follows from the Smith normal form.

\cref{itm:trop-mono:finflat} After translating we may take $c=1$. Writing
$A=UDV$ with $U,V\in\mathrm{GL}_n(\mathbf Z)$ and $D$ diagonal with entries
$d_1,\ldots,d_n$, the map $\varphi$ becomes the composition of two
automorphisms of $\mathbf G_m^n$ with the map
$(t_i)\mapsto(t_i^{d_i})$; the latter corresponds to the ring extension
$K[T^{\pm1}]\subset K[T^{\pm1}]$, $T_i\mapsto T_i^{d_i}$, which is free of
rank $\prod|d_i|=\delta$. Freeness gives finiteness and flatness, and the
degree is $\delta$.

\cref{itm:trop-mono:preimage} and \cref{itm:trop-mono:degree}. Let
$x\in\mathbf G_m^{n,\An}$ and $t=\varphi(x)$. The initial forms
$\widetilde{T_i(t)}=\widetilde{c_i}\,\widetilde{T(x)}^{A_i}$ are Laurent
monomials in the $\widetilde{T_j(x)}$ with exponent matrix $A$. Since $A$
is invertible over $\mathbf Q$, some power of each $\widetilde{T_j(x)}$
is a monomial in the $\widetilde{T_i(t)}$ up to a unit of $\widetilde K$;
hence the two families are simultaneously algebraically independent over
$\widetilde K$. By \Cref{prop:trop-torus-skeleton}
\cref{itm:trop-sk:graded}, $x\in S$ if and only if $t\in S$. Since
$\varphi_{\mathrm{trop}}=A$ is bijective on $\mathbf R^n$, the induced map
on skeleta is a homeomorphism.

Suppose now $x\in S$ and let $E\subset\widetilde{\mathcal H(x)}$ be the
graded subfield generated by $\widetilde K$ and the $\widetilde{T_j(x)}$,
and $F\subset E$ the graded subfield generated by $\widetilde K$ and the
$\widetilde{T_i(t)}$. Both are graded Laurent polynomial fields in $n$
independent variables, and $E/F$ is the extension of graded fraction fields
determined by the lattice inclusion $A\mathbf Z^n\subset\mathbf Z^n$, so
$[E:F]=[\mathbf Z^n:A\mathbf Z^n]=\delta$. Now $E=\widetilde{\mathcal H(x)}$
and $F=\widetilde{\mathcal H(t)}$: indeed $x=\eta_u$ and $t=\eta_{Au+v_K(c)}$
lie on the skeleton, $K[T^{\pm1}]$ is dense in $\mathcal H(\eta_u)$, and if
$|g-h|<|h|$ then $\widetilde g=\widetilde h$
(\Cref{prop:appG-initial-calculus}), so every initial form in
$\widetilde{\mathcal H(\eta_u)}$ is the initial form of a Laurent
polynomial, that is, a $\widetilde K$-linear combination of monomials in the
$\widetilde{T_j(x)}$. Therefore
$[\widetilde{\mathcal H(x)}:\widetilde{\mathcal H(t)}]=\delta$. Since the
graded residue degree of an extension of complete valued fields is the
product $ef$ of its ramification index and residue degree, the fundamental
inequality (\Cref{prop:fundamental-inequality}) gives
$[\mathcal H(x):\mathcal H(t)]\ge\delta$, once we know that this degree is
finite; and it is, because $\varphi$ is finite. As $\varphi$ is finite
and flat of degree $\delta$ by \cref{itm:trop-mono:finflat}, the fibre
$\varphi^{-1}(t)$ is an $\mathcal H(t)$-algebra of dimension $\delta$
containing $\mathcal H(x)$ as a residue field of one of its local factors;
comparing dimensions forces the fibre to be the single reduced point
$\mathcal M(\mathcal H(x))$ and $\deg_x(\varphi)=\delta$.
\end{proof}

\begin{example}[The tropicalization of a torus, computed]
\label{ex:trop-torus-explicit}
Take $X=T=\mathbf G_{m,K}^n$ and the identity moment map $f=\mathrm{id}$.
Then:
\begin{enumerate}
\item $\Trop_f(T^{\An})=\mathbf R^n$. Indeed $\trop(\eta_u)=u$ for all $u$
  by \Cref{lem:trop-gauss-multiplicative}, so $\trop$ is surjective; and
  $\Trop_f\subset\mathbf R^n$ trivially. Thus the tropicalization of a
  torus is all of $N_{\mathbf R}$, with no polyhedral boundary.
\item The fibre of $\trop$ over $u$ is the set of $x$ with $x\le\eta_u$,
  by the argument in the proof of \Cref{prop:trop-torus-skeleton}; it is a
  compact connected set with a unique maximal point. For $n=1$ and
  $K$ algebraically closed and nontrivially valued it is the set of points
  of the analytic circle $\{|T|=\rme^{-u}\}$, whose unique maximal element
  is the type-II or type-III point $\zeta_{0,\rme^{-u}}$ of
  \Cref{thm:berk-a1-classification}; the description of the fibre as the set
  of points dominated by a Gauss point is \Cref{cor:berk-gauss-sup}.
\item The local skeleton is $\Sigma_{\mathrm{id}}=S(T)$, by
  \Cref{prop:trop-torus-skeleton} and the definition
  \eqref{eq:trop-sigma-def} below, and $\trop|_{S(T)}$ is an isomorphism of
  integer-affine spaces onto $\mathbf R^n$.
\item The canonical calibration of $S(T)$ is
  $\mu=1\cdot|e_1\wedge\cdots\wedge e_n|$: taking $q=\mathrm{id}$ in
  \Cref{def:trop-canonical-calibration}, the map $q\circ f=\mathrm{id}$ is
  finite flat of degree $1$, and the affine section $\sigma_q$ is the
  identity. In particular the tropical multiplicity of a torus is $1$, and
  every multiplicity computed later is a comparison with this one.
\item The map $\rho=\sq\circ\trop:T^{\An}\to S(T)$ is a continuous
  retraction. It is the toric case of the skeletal retractions of
  \Cref{prop:skeleton-retraction}; the comparison with the retraction
  attached to a toric model, and with the support-function calculus, is the
  content of \Cref{sec:toric-convex}, whose coordinate
  \eqref{eq:toric-trop} agrees with $\trop$ here.
\end{enumerate}
\end{example}

\begin{example}[Why naive fibre counting fails]
\label{ex:trop-fibre-counting}
Let $K=\mathbf Q_p$, $n=1$, and $\varphi(T)=T^p$. By
\Cref{lem:trop-monomial-degree}, $\varphi$ is finite flat of degree $p$
and, for every $u$, the fibre $\varphi^{-1}(\eta_{pu})$ inside
$\mathbf G_m^{\An}$ meets $S(\mathbf G_m)$ in the single point $\eta_u$,
with $\deg_{\eta_u}(\varphi)=p$. Counting points of the fibre would give
$1$; the correct answer $p$ is the degree of the graded residue extension
$\widetilde{\mathcal H(\eta_u)}/\widetilde{\mathcal H(\eta_{pu})}$, which is
generated by an element $\widetilde{T}$ with
$\widetilde T^{\,p}=\widetilde{T^p}$ in the appropriate degree. Over an
algebraically closed field the same map has, over a classical point,
$p$ distinct preimages when $\operatorname{char}K\ne p$ --- which is the case
here, since $\mu_p\subset\widehat{\overline{\mathbf Q_p}}$ --- and one
preimage of multiplicity $p$ only in equal characteristic $p$; note that the
criterion is the characteristic of $K$, not of $\widetilde K$.  The analytic
degree is $p$ in every case. This is why \Cref{def:trop-canonical-calibration} is phrased with the
local degree of a finite flat morphism and never with a cardinality.
\end{example}

\section{Tropical charts}
\label{sec:trop-charts}

\begin{definition}
\label{def:trop-chart}
A \emph{tropical chart} on $X$ is a pair $(V,\varphi)$ consisting of a moment
map on an analytic domain $U$ and an open subset
$\Omega\subset\operatorname{Trop}_\varphi(U)$ such that
\[
                         V=\varphi_{\mathrm{trop}}^{-1}(\Omega).
\]
A superform on $V$ represented in this chart is the pullback
$\varphi_{\mathrm{trop}}^*\alpha$ of a superform $\alpha$ on $\Omega$.
\end{definition}

The notation is symbolic at this stage: $\varphi_{\mathrm{trop}}$ is only
continuous, so pullback means the equivalence class represented by the chart,
not classical differentiation of a map between manifolds.

Two structural facts about charts are used constantly and are best isolated
now: the domains of charts form a basis of the topology near any good point,
and the intersection of two chart domains is again covered by chart domains.

\begin{lemma}
\label{lem:trop-charts-basis}
Let $X$ be a $K$-analytic space, let $x\in X$ be a point at which $X$ is
good, and let $W\ni x$ be open. Then there is a tropical chart
$(V,\varphi)$ with $x\in V\subset W$. If moreover $d_K(x)=n$, the moment map
$\varphi$ may be chosen with $d_x(\varphi)=n$.
\end{lemma}

\begin{proof}
Goodness gives an affinoid neighbourhood $\mathcal M(A)\subset W$ of $x$.
Choose finitely many $g_1,\ldots,g_r\in A$ generating $A$ topologically,
and let $J=\{i:|g_i(x)|>0\}$. On the Laurent domain
$U=\{y:|g_i(y)|\ge\varepsilon_i,\ i\in J\}$ with
$0<\varepsilon_i<|g_i(x)|$, the functions $g_i$, $i\in J$, are invertible,
so $\varphi=(g_i)_{i\in J}$ is a moment map on $U$. The set
\[
 \Omega=\{w\in\Trop_\varphi(U):w_i<-\log\varepsilon_i\ \text{for }i\in J\}
\]
is open in $\Trop_\varphi(U)$ and contains $\varphi_{\mathrm{trop}}(x)$, and
$V=\varphi_{\mathrm{trop}}^{-1}(\Omega)$ is the open subset
$\{y\in U:|g_i(y)|>\varepsilon_i\ \forall i\in J\}$ of $U$, hence contained
in $\mathcal M(A)\subset W$. This is a tropical chart containing $x$.

For the refinement, suppose $d_K(x)=n$ and choose a graded transcendence
basis $a_1,\ldots,a_n$ of $\widetilde{\mathcal H(x)}$ over $\widetilde K$.
\Cref{lem:trop-moment-local-existence} produces invertible
$f_1,\ldots,f_n$ near $x$ with $\widetilde{f_j(x)}=a_j$. Replacing
$\varphi$ by $(\varphi,f_1,\ldots,f_n)$ --- a moment map on the
intersection of the two domains, by \Cref{lem:trop-moment-basic}
\cref{itm:trop-moment:prod} --- gives $d_x(\varphi)\ge n$, hence $=n$ by
\Cref{lem:trop-residual-degree-bounds}.
\end{proof}

\begin{definition}
\label{def:trop-chart-refinement}
A chart $(V',\psi)$ \emph{refines} $(V,\varphi)$ if $V'\subset V$ and there is
an integer-affine map $F$ between tropical targets with
\[
                         \varphi_{\mathrm{trop}}|_{V'}
                         =F\circ\psi_{\mathrm{trop}}.
\]
Then a representative $\alpha$ in the first chart is represented by
$F^*\alpha$ in the refinement.
\end{definition}

\begin{lemma}\label{lem:product-moment-map}
If $(V_1,\varphi)$ and $(V_2,\psi)$ are two charts, then on every suitable
analytic domain contained in $V_1\cap V_2$ the product
\[
                     (\varphi,\psi):U\longrightarrow\mathbf G_m^{r+s}
\]
is a moment map refining both charts through the coordinate projections.
\end{lemma}
\begin{proof}
The union of the invertible analytic coordinate functions of $\varphi$ and
$\psi$ defines the product moment map. Tropicalization of the product is the
ordered pair of the two tropicalizations. The coordinate projections are
integer linear and give the required factorization.
\end{proof}

We spell out the two properties of refinement that the independence proofs
use, together with the warning that refinement is not a symmetric relation
and that a product refinement genuinely enlarges the target dimension.

\begin{lemma}
\label{lem:trop-refinement-calculus}
\begin{enumerate}
\item\label{itm:trop-ref:trans} Refinement is transitive: if
  $(V'',\chi)$ refines $(V',\psi)$ through $G$ and $(V',\psi)$ refines
  $(V,\varphi)$ through $F$, then $(V'',\chi)$ refines $(V,\varphi)$ through
  $F\circ G$.
\item\label{itm:trop-ref:prod} For moment maps $\varphi:U\to T$ and
  $\psi:U\to T'$ on the same domain, the product $(\varphi,\psi)$ refines
  both through the projections $p:T\times T'\to T$, $p':T\times T'\to T'$,
  and $p_{\mathrm{trop}}$, $p'_{\mathrm{trop}}$ are surjective integral
  linear maps.
\item\label{itm:trop-ref:sigma} If $(V',\psi)$ refines $(V,\varphi)$ through
  $F$, then $d_x(\varphi)\le d_x(\psi)$ for every $x\in V'$; in particular
  $\Sigma^{(n)}_{\varphi}\cap V'\subset\Sigma^{(n)}_{\psi}$ in the notation
  of \Cref{sec:trop-local-skeleta}.
\item\label{itm:trop-ref:trans2} If $\varphi_{\mathrm{trop}}$ and
  $\varphi'_{\mathrm{trop}}$ differ by a constant on $V$, then each of the
  two charts refines the other through a translation.
\end{enumerate}
\end{lemma}

\begin{proof}
\cref{itm:trop-ref:trans} A composition of integer-affine maps is integer
affine, and
$\varphi_{\mathrm{trop}}|_{V''}=F\circ\psi_{\mathrm{trop}}|_{V''}
=F\circ G\circ\chi_{\mathrm{trop}}$.

\cref{itm:trop-ref:prod} is \Cref{lem:trop-moment-basic}
\cref{itm:trop-moment:prod}; the projections induce the inclusions
$M\hookrightarrow M\oplus M'$ on characters, whose duals are the coordinate
projections, which are surjective with $\mathbf Z$-saturated image.

\cref{itm:trop-ref:sigma} Composition with a homomorphism of tori shrinks the
generated graded subfield, so $F_x(\varphi)\subset F_x(\psi)$ by
\Cref{lem:trop-residual-degree-bounds}; hence
$d_x(\varphi)\le d_x(\psi)$.

\cref{itm:trop-ref:trans2} A translation is integer affine with linear part
the identity, and is invertible with integer-affine inverse.
\end{proof}

\begin{warning}{Refinement is directed, not symmetric}
If $(V',\psi)$ refines $(V,\varphi)$, the coordinates of $\varphi$ are
recovered from those of $\psi$ but not conversely: information is lost, not
gained, along $F$. Two arbitrary charts therefore have no common
\emph{coarsening}; the only universal construction is the product, which
refines both. Consequently every comparison of two charts below proceeds by
passing \emph{up} to the product, never down, and the corresponding
statements about weights are pushforward statements, carrying the lattice
index \eqref{eq:appG-lattice-index}.
\end{warning}

\begin{example}[A common refinement with a nontrivial index]
\label{ex:trop-common-refinement-index}
Let $A=\{r\leq |T|\leq1\}$ be a closed annulus and fix an integer
$m\geq1$.  The two moment maps $T$ and $T^m$ have tropical coordinates
$u=-\log|T|$ and $u'=mu$.  Their product has image
\[
             \{(u,mu):0\leq u\leq-\log r\}\subset\mathbf R^2,
\]
whose primitive tangent vector is $(1,m)$.  Projection to the first factor
has lattice index $1$, whereas projection to the second has lattice index
$m$.  Thus the product chart has weight $1$, the $T$-chart has weight $1$,
and the $T^m$-chart has weight $m$.  Formula
\eqref{eq:analytic-tropical-pushforward} therefore reproduces the analytic
degree of $T\mapsto T^m$ exactly once.  This one-dimensional calculation is
a useful normalization test for every chart-refinement formula below.
\end{example}

\begin{proposition}\label{prop:tropical-form-refinement}
Suppose representatives $\alpha$ and $\beta$ in two tropical charts have
equal pullbacks to a common product refinement. Then equality is preserved
under further refinement, wedge products, $\dprime$, $\ddprime$, and pullback by an
analytic morphism. Hence compatible representatives define a sheaf
$A_X^{p,q}$.
\end{proposition}
\begin{proof}
Further pullback preserves equality. The compatibility with wedge and the
differentials was proved for affine pullback in \Cref{chap:superforms}. An
analytic morphism pulls a moment map back by composing its invertible
coordinates; product refinements commute with this operation. The sheaf
axioms follow by taking common refinements on overlaps and using locality of
smooth forms on polyhedral supports.
\end{proof}

\section{Local skeleta}
\label{sec:trop-local-skeleta}

Tropicalization is not usually injective, and its image is not usually of the
right dimension. Inside $X$ there is, however, a canonical closed subset on
which the moment coordinates attain maximal rank. It is a polyhedral space,
tropicalization restricts to a piecewise immersion on its top-dimensional
part, and it is the correct domain of integration. Following
\cite[\S12.6]{CLD25} we call it the characteristic polyhedron of the moment
map; because it plays the role that skeleta of models play in
\Cref{chap:skeleta}, and because it is defined only relative to a moment map
on a domain, we also call it a \emph{local skeleton}.

\begin{definition}[Local skeleton]
\label{def:trop-local-skeleton}
\index{skeleton!local, of a moment map}\index{characteristic polyhedron}
Let $X$ be a $K$-analytic space of level $n$ and let $f:V\to T$ be a moment
map on an analytic domain $V\subset X$. The \emph{local skeleton} of $f$ is
\begin{equation}
\label{eq:trop-sigma-def}
 \Sigma_f=\bigcup_q\,(q\circ f)^{-1}\bigl(S(\mathbf G_{m}^n)\bigr)
 \subset V,
\end{equation}
the union being over all homomorphisms of tori $q:T\to\mathbf G_{m}^n$. We
write $\Sigma_f^{(n)}$ for the union of the local components of $\Sigma_f$ of
dimension $n$, that is, the set of $x\in\Sigma_f$ with
$\dim_x\Sigma_f=n$.
\end{definition}

Two remarks before the theorem. First, \eqref{eq:trop-sigma-def} is a union
over an infinite set, so closedness does not follow formally; the theorem shows it is
a finite union. Second, the number $n$ in \eqref{eq:trop-sigma-def} is the
level, not $\dim T$: one always projects to a torus of dimension exactly $n$,
and $\Sigma_f=\varnothing$ as soon as $\dim T<n$.

\begin{theorem}
\label{thm:trop-sigma-characterization}
Let $X$ be of level $n$, let $f:V\to T$ be a moment map, and fix a
trivialization $T\simeq\mathbf G_{m}^m$, so that $f=(f_1,\ldots,f_m)$ with
$f_i\in\mathcal O(V)^\times$. For $x\in V$ the following are equivalent.
\begin{enumerate}
\item\label{itm:trop-sigmachar:def} $x\in\Sigma_f$.
\item\label{itm:trop-sigmachar:deg} $d_x(f)=n$.
\item\label{itm:trop-sigmachar:coord} There is a subset
  $I\subset\{1,\ldots,m\}$ with $\#I=n$ such that the initial forms
  $\bigl(\widetilde{f_i(x)}\bigr)_{i\in I}$ are algebraically independent
  over $\widetilde K$.
\item\label{itm:trop-sigmachar:proj} There is $I\subset\{1,\ldots,m\}$ with
  $\#I=n$ such that $x\in(q_I\circ f)^{-1}(S(\mathbf G_m^n))$, where
  $q_I:\mathbf G_m^m\to\mathbf G_m^n$ is the projection onto the coordinates
  indexed by $I$.
\end{enumerate}
Consequently
\begin{equation}
\label{eq:trop-sigma-finite-union}
 \Sigma_f=\bigcup_{\#I=n}(q_I\circ f)^{-1}\bigl(S(\mathbf G_m^n)\bigr),
\end{equation}
a finite union of closed subsets; in particular $\Sigma_f$ is closed in $V$,
and $\Sigma_f=\varnothing$ if $\dim T<n$.
\end{theorem}

\begin{proof}
\cref{itm:trop-sigmachar:def} $\Rightarrow$ \cref{itm:trop-sigmachar:deg}.
Suppose $q\circ f(x)\in S(\mathbf G_m^n)$ for some $q$. By
\Cref{prop:trop-torus-skeleton} \cref{itm:trop-sk:graded} applied to
$\mathbf G_m^n$ at the point $q(f(x))$ we get $d_x(q\circ f)=n$, and
$d_x(q\circ f)\le d_x(f)$ by \Cref{lem:trop-residual-degree-bounds}. On the
other hand $d_x(f)\le d_K(x)\le\dim_xX\le n$ by the same lemma and the
level hypothesis. Hence $d_x(f)=n$.

\cref{itm:trop-sigmachar:deg} $\Rightarrow$ \cref{itm:trop-sigmachar:coord}.
The graded field $F_x(f)$ is generated over $\widetilde K$ by the family
$\bigl(\widetilde{f_i(x)}\bigr)_{1\le i\le m}$, since the initial form of
$f_i^{-1}$ is the inverse of that of $f_i$
(\Cref{prop:appG-initial-calculus}) and every $\chi^m\circ f$ is a monomial
in the $f_i^{\pm1}$. A generating family of a field extension of
transcendence degree $n$ contains a transcendence basis; the graded case is
identical, because a graded field extension is algebraic exactly when its
homogeneous generators are, and the exchange lemma holds verbatim for
homogeneous elements. Choose $I$ indexing such a basis.

\cref{itm:trop-sigmachar:coord} $\Rightarrow$ \cref{itm:trop-sigmachar:proj}.
The subfield $F_x(q_I\circ f)$ is generated by the
$\widetilde{f_i(x)}$, $i\in I$, so $d_x(q_I\circ f)=n$, and
\Cref{prop:trop-torus-skeleton} \cref{itm:trop-sk:graded} applied at
$q_I(f(x))\in\mathbf G_m^{n,\An}$ gives $q_I(f(x))\in S(\mathbf G_m^n)$.

\cref{itm:trop-sigmachar:proj} $\Rightarrow$ \cref{itm:trop-sigmachar:def} is
trivial, the $q_I$ being particular homomorphisms of tori. Equation
\eqref{eq:trop-sigma-finite-union} restates the equivalence
\cref{itm:trop-sigmachar:def} $\Leftrightarrow$
\cref{itm:trop-sigmachar:proj}. Each set in the union is closed, since
$S(\mathbf G_m^n)$ is closed by \Cref{prop:trop-torus-skeleton} and
$q_I\circ f$ is continuous; the union is finite. If $\dim T=m<n$ there is
no admissible $I$, so $\Sigma_f=\varnothing$.
\end{proof}

The characterization is worth restating in words: \emph{a point lies in the
local skeleton of a moment map exactly when $n$ of the moment coordinates
have algebraically independent initial forms there}. No cardinality, no
fibre, no model is involved. This is the precise sense in which graded
reduction, and not any counting procedure, controls the theory.

\begin{corollary}
\label{cor:trop-sigma-abhyankar}
Let $x\in\Sigma_f$. Then
\[
 d_x(f)=d_K(x)=\dim_xX=n .
\]
Moreover $x$ belongs to no Zariski-closed subset of dimension $<n$ of any
analytic domain of $X$ containing it. In particular $x$ lies on exactly one
irreducible component of $X$, in which it is Zariski dense, and
$\Sigma_f\subset X^{(n)}$, the union of the $n$-dimensional irreducible
components.
\end{corollary}

\begin{proof}
The chain $n=d_x(f)\le d_K(x)\le\dim_xX\le n$ of
\Cref{thm:trop-sigma-characterization} and
\Cref{lem:trop-residual-degree-bounds} forces equality throughout. Suppose
$x\in Y$ with $Y$ Zariski closed in an analytic domain $W\ni x$ and
$\dim_xY<n$. The completed residue field $\mathcal H(x)$ computed in $Y$ and
in $W$ is the same, so $d_K(x)$ is the same; applying
$d_K(x)\le\dim_xY$ (\Cref{lem:trop-residual-degree-bounds} on $Y$) gives
$n\le\dim_xY<n$, absurd. If $x$ lay on two irreducible components, it would
lie on their intersection, a Zariski-closed subset of dimension $<n$;
similarly the closure of $\{x\}$ inside its component must be everything.
\end{proof}

The polyhedral structure of $\Sigma_f$ is the one substantial import of this
section. We state it once, in the form used below, and reference it
explicitly at each use.

\begin{theorem}
\label{thm:trop-sigma-polyhedral}
Let $X$ be a topologically separated $K$-analytic space of level $n$ and let
$f:V\to T$ be a moment map on an analytic domain. Then:
\begin{enumerate}
\item\label{itm:trop-poly:para} $\Sigma_f$ is a closed polyhedral subspace of
  $V$ of dimension at most $n$, in the sense of the tropical spaces of
  \Cref{chap:tropical-spaces};
\item\label{itm:trop-poly:top} $\Sigma_f^{(n)}$ is a closed polyhedral
  subspace, purely of dimension $n$, and $f_{\mathrm{trop}}$ restricts on it
  to a piecewise immersion of polyhedral spaces;
\item\label{itm:trop-poly:union} local skeleta are stable under
  restriction to analytic domains, $\Sigma_{f|_W}=\Sigma_f\cap W$, and under
  preimage, $\Sigma_{f\circ p}=p^{-1}(\Sigma_f)$ for a morphism $p$ of
  spaces of level $n$; a locally finite union of local skeleta is again
  $G$-locally contained in one.
\end{enumerate}
\end{theorem}

\begin{proof}
Assertions \cref{itm:trop-poly:para} and \cref{itm:trop-poly:top} are proved
in \cite[\S\S12.6 and 12.9]{CLD25}; the argument
there proceeds by a toric Noether normalization reducing $\dim T$ to $n$
and then by induction on $n$ using
\Cref{thm:trop-local-germ}. The two stability assertions of
\cref{itm:trop-poly:union} are elementary and we prove them. Write
$S=S(\mathbf G_m^n)$. For a domain $W\subset V$, a morphism $p:Y\to V$, and
any $q$, one has
\[
 (q\circ f|_W)^{-1}(S)=(q\circ f)^{-1}(S)\cap W,
 \qquad
 (q\circ f\circ p)^{-1}(S)=p^{-1}\bigl((q\circ f)^{-1}(S)\bigr),
\]
and the union over $q$ commutes with intersection and with preimage. The
last assertion of
\cref{itm:trop-poly:union} is \cite[Proposition~12.6.7]{CLD25}.
\end{proof}

\begin{example}[Local skeleta in dimension one]
\label{ex:trop-sigma-dim-one}
Let $K$ be nontrivially valued and let $X=\mathbf G_{m,K}^{\An}$, of level
$1$.
\begin{enumerate}
\item For $f=T$ one has $\Sigma_f=S(\mathbf G_m)=\{\eta_u:u\in\mathbf R\}$,
  a closed real line inside $X$; this is the skeleton of the annulus
  computation of \Cref{ex:berk-standard-annulus} and of
  \Cref{ex:retraction-annulus}, extended to the whole punctured line.
\item For $f=T^m$ with $m\ne0$, \Cref{lem:trop-monomial-degree}
  \cref{itm:trop-mono:preimage} gives $\Sigma_{T^m}=\Sigma_T$: the local
  skeleton does not see the exponent. What does see it is the calibration,
  computed in \Cref{ex:trop-two-charts}.
\item For $f=(T,1+T)$ on the domain $V=\{|T|<1\}\cap\mathbf G_m^{\An}$, one
  has $|1+T|=1$ on $V$, so $\widetilde{(1+T)(x)}=1$ and the pair
  $\bigl(\widetilde{T(x)},\widetilde{(1+T)(x)}\bigr)$ has the same
  graded span as $\widetilde{T(x)}$ alone. By
  \Cref{thm:trop-sigma-characterization},
  $\Sigma_{(T,1+T)}=\Sigma_T\cap V$.
\item Take $X=\mathcal M(K)$ or, more generally, $X$ of dimension $0$ and
  level $1$. Then $\dim_xX=0<1$, so $\Sigma_f=\varnothing$ for every moment
  map, by \Cref{cor:trop-sigma-abhyankar}. Nothing is integrated over a
  space of the wrong dimension, and no convention is needed to say so.
\end{enumerate}
\end{example}

\section{Dimension formulae}
\label{sec:trop-dimension}

The tropical dimension of a moment map at a point measures how much of the
local geometry survives tropicalization. It is defined as an infimum over
shrinking neighbourhoods precisely because a global tropical image can be
larger than any germ suggests.

\begin{definition}
\label{def:trop-tropical-dimension}
\index{tropical dimension}
Let $f:X\to T$ be a moment map and $x\in X$. The \emph{tropical dimension}
of $f$ at $x$ is
\[
 \dim_{\mathrm{trop},x}(f)
 =\inf_V\ \dim_{f_{\mathrm{trop}}(x)}f_{\mathrm{trop}}(V),
\]
the infimum over compact analytic neighbourhoods $V$ of $x$ in $X$. By
\Cref{thm:trop-local-germ} \cref{itm:trop-germ:germ} the infimum is attained
on every sufficiently small $V$.
\end{definition}

\begin{theorem}
\label{thm:trop-dimension-formula}
Let $X$ be a $K$-analytic space, $f:X\to T$ a moment map, and $x\in X$.
Then
\begin{equation}
\label{eq:trop-dimension-chain}
 \dim_{\mathrm{trop},x}(f)\ \le\ d_x(f)\ \le\ d_K(x)\ \le\ \dim_xX .
\end{equation}
The first inequality is an equality if $x\notin\partial X$. The second is an
equality if and only if the graded subfield $F_x(f)$ generated by the moment
coordinates is a graded transcendence basis of
$\widetilde{\mathcal H(x)}$ up to algebraic extension. The third is an
equality if and only if $x$ is an Abhyankar point of $X$, and then
$\mathcal H(x)$ is an Abhyankar extension of $K$ in the sense of
\Cref{sec:trop-graded}. If $X$ is compact,
\begin{equation}
\label{eq:trop-dimension-sup}
 \dim\Trop_f(X)=\sup_{x\in X}\dim_{\mathrm{trop},x}(f).
\end{equation}
\end{theorem}

\begin{proof}
The middle and right inequalities of \eqref{eq:trop-dimension-chain} are
\Cref{lem:trop-residual-degree-bounds}, together with their equality cases,
which are the definitions of transcendence basis and of an Abhyankar point.
The left inequality is \Cref{thm:trop-local-germ}
\cref{itm:trop-germ:bound}, and the equality for interior $x$ is
\Cref{thm:trop-local-germ} \cref{itm:trop-germ:interior}: there
$f_{\mathrm{trop}}(V)$ is purely of dimension $d_x(f)$ at
$f_{\mathrm{trop}}(x)$ for every small $V$, so the infimum equals
$d_x(f)$.

For \eqref{eq:trop-dimension-sup}: the inequality $\ge$ is clear because
$f_{\mathrm{trop}}(V)\subset\Trop_f(X)$ for every $V$. For $\le$, let
$\xi\in\Trop_f(X)$ be a point at which the finite polyhedral set
$\Trop_f(X)$ --- polyhedral by
\Cref{thm:analytic-tropical-polyhedrality} --- has its
maximal dimension $\delta$, and choose a $\delta$-dimensional cell $C$ of it
containing $\xi$ in its relative interior. Since $X$ is compact,
$f_{\mathrm{trop}}^{-1}(C)$ is a compact analytic domain whose tropical image
is $C$; covering it by finitely many small compact neighbourhoods $V_j$ and
using that a finite union of sets of dimension $<\delta$ cannot equal $C$,
some $V_j$ has $\dim f_{\mathrm{trop}}(V_j)=\delta$, and then some point of
$V_j$ has tropical dimension $\delta$ by the same argument applied inside
$V_j$ and \Cref{thm:trop-local-germ} \cref{itm:trop-germ:germ}.
\end{proof}

We now connect the tropical dimension to the local skeleton. The next lemma
is the geometric heart: a top-dimensional piece of the tropical image is
always \emph{hit} by the canonical section, so top-dimensional tropical
behaviour is genuinely present in the analytic space.

\begin{lemma}
\label{lem:trop-skeleton-attained}
Let $X$ be a compact $K$-analytic space and $f:X\to T$ a moment map with
$\dim T=n$. Then
\[
 f(X)\ \supset\ \sq\bigl(\Trop_f(X)^{(n)}\bigr),
\]
where $\Trop_f(X)^{(n)}$ denotes the union of the $n$-dimensional local
components of the polyhedron $\Trop_f(X)$.
\end{lemma}

\begin{proof}
First let $C\subset\Trop_f(X)$ be a compact $n$-dimensional polyhedron; we
show that $f^{-1}(\sq(C))\ne\varnothing$. Writing $C$ as a finite
intersection of half-spaces $\langle m_i,u\rangle\le c_i$ with
$m_i\in M$, the set
$W=f_{\mathrm{trop}}^{-1}(C)=\{x:|\chi^{m_i}(f(x))|\ge\rme^{-c_i}\}$ is a
compact analytic domain of $X$ with $f_{\mathrm{trop}}(W)=C$. Since
$\dim f_{\mathrm{trop}}(W)=n$ and $W$ is compact,
\eqref{eq:trop-dimension-sup} gives a point $x\in W$ with
$\dim_{\mathrm{trop},x}(f|_W)=n$. By \eqref{eq:trop-dimension-chain},
$d_x(f)=n=\dim T$, so $f(x)\in S(T)$ by
\Cref{prop:trop-torus-skeleton} \cref{itm:trop-sk:graded}. As
$\trop(f(x))\in C$, we get $f(x)\in\sq(C)$.

Now let $s\in\Trop_f(X)^{(n)}$ and let $(C_\alpha)$ be a basis of compact
neighbourhoods of $s$ in $\Trop_f(X)$ consisting of $n$-dimensional
polyhedra. By the previous paragraph $f(X)$ meets each $\sq(C_\alpha)$, and
$(\sq(C_\alpha))$ is a basis of neighbourhoods of $\sq(s)$ in $S(T)$ because
$\sq$ is a homeomorphism onto $S(T)$
(\Cref{prop:trop-torus-skeleton}). Hence $\sq(s)$ lies in the closure of
$f(X)$; and $f(X)$ is compact, hence closed.
\end{proof}

\begin{theorem}
\label{thm:trop-dimension-sigma}
Let $X$ be of level $n$ and $f:X\to T$ a moment map.
\begin{enumerate}
\item\label{itm:trop-dimsig:a} For every $x\in X$ one has
  $\dim_{\mathrm{trop},x}(f)\le n$, with equality if and only if
  $x\in\Sigma_f^{(n)}$.
\item\label{itm:trop-dimsig:b} If $x\in\Int(X)$, the three conditions
  $x\in\Sigma_f$, $\dim_{\mathrm{trop},x}(f)=n$, and
  $x\in\Sigma_f^{(n)}$ are equivalent.
\item\label{itm:trop-dimsig:c} If $X$ is compact, or more generally if
  $f_{\mathrm{trop}}$ is topologically proper, then
  $\Trop_f(X)^{(n)}=f_{\mathrm{trop}}\bigl(\Sigma_f^{(n)}\bigr)$.
\item\label{itm:trop-dimsig:d} If $X$ is boundaryless, then
  $\Sigma_f=\Sigma_f^{(n)}$, and $\Sigma_f$ depends only on the map
  $f_{\mathrm{trop}}$, not on $f$.
\end{enumerate}
\end{theorem}

\begin{proof}
\cref{itm:trop-dimsig:a} The bound $\dim_{\mathrm{trop},x}(f)\le d_x(f)\le n$
is \eqref{eq:trop-dimension-chain} and the level hypothesis.

Suppose $\dim_{\mathrm{trop},x}(f)=n$. Then $d_x(f)=n$, so $x\in\Sigma_f$
by \Cref{thm:trop-sigma-characterization}. We must show
$\dim_x\Sigma_f=n$. Let $V$ be any compact analytic neighbourhood of $x$;
then $\dim f_{\mathrm{trop}}(V)=n$. Choose an $n$-dimensional cell
$C\subset f_{\mathrm{trop}}(V)$ and a homomorphism $q:T\to\mathbf G_m^n$
with $q_{\mathrm{trop}}$ injective on the affine span $\langle C\rangle$;
such a $q$ exists because $\langle C\rangle$ is a $\mathbf Q$-rational
affine subspace of $N_{\mathbf R}$, so a generic projection $M\to\mathbf Z^n$
restricts to an injection on its direction lattice. Then
$\dim(q\circ f)_{\mathrm{trop}}(V)=n$, and
\Cref{lem:trop-skeleton-attained} applied to $q\circ f$ on $V$ gives
\[
 \sq\bigl((q\circ f)_{\mathrm{trop}}(V)^{(n)}\bigr)\subset(q\circ f)(V).
\]
Every point of $V$ mapping into $S(\mathbf G_m^n)$ under $q\circ f$ lies in
$\Sigma_f$ by \eqref{eq:trop-sigma-def}; hence $(q\circ f)(\Sigma_f\cap V)$
contains an $n$-dimensional subset of $S(\mathbf G_m^n)$, so
$\dim(\Sigma_f\cap V)\ge n$. Combined with
\Cref{thm:trop-sigma-polyhedral} \cref{itm:trop-poly:para} this gives
$\dim(\Sigma_f\cap V)=n$ for every $V$, that is $\dim_x\Sigma_f=n$.

Conversely let $x\in\Sigma_f^{(n)}$ and let $V$ be a compact neighbourhood
of $x$. Then $V\cap\Sigma_f^{(n)}$ is a neighbourhood of $x$ in
$\Sigma_f^{(n)}$, hence contains an $n$-dimensional cell; since
$f_{\mathrm{trop}}$ restricted to $\Sigma_f^{(n)}$ is a piecewise immersion
(\Cref{thm:trop-sigma-polyhedral} \cref{itm:trop-poly:top}), the image
$f_{\mathrm{trop}}(V)$ has dimension at least $n$, hence exactly $n$.
Therefore $\dim_{\mathrm{trop},x}(f)=n$.

\cref{itm:trop-dimsig:b} Let $x\in\Int(X)$ and $x\in\Sigma_f$. Then
$d_x(f)=n$ by \Cref{thm:trop-sigma-characterization}, and since
$x\notin\partial X$ the first inequality of \eqref{eq:trop-dimension-chain}
is an equality, so $\dim_{\mathrm{trop},x}(f)=n$. By
\cref{itm:trop-dimsig:a} this gives $x\in\Sigma_f^{(n)}$. The implications
$\Sigma_f^{(n)}\subset\Sigma_f$ and \cref{itm:trop-dimsig:a} close the
circle.

\cref{itm:trop-dimsig:c} Since $f_{\mathrm{trop}}|_{\Sigma_f^{(n)}}$ is a
piecewise immersion from a space purely of dimension $n$, its image is
purely of dimension $n$ and is closed when $f_{\mathrm{trop}}$ is proper;
hence it is contained in $\Trop_f(X)^{(n)}$. Conversely let
$\xi\in\Trop_f(X)$ not lie in $f_{\mathrm{trop}}(\Sigma_f^{(n)})$ and let
$F=f_{\mathrm{trop}}^{-1}(\xi)$. Every $x\in F$ lies outside
$\Sigma_f^{(n)}$, so \cref{itm:trop-dimsig:a} and the definition of local
tropical dimension give a compact analytic neighbourhood $V_x$ of $x$ with
$\dim f_{\mathrm{trop}}(V_x)\le n-1$. Properness makes $F$ compact; choose
$x_1,\ldots,x_r$ with $F\subset V:=\bigcup_i\mathring V_{x_i}$. The closed
set $X\setminus V$ has closed image under the proper map
$f_{\mathrm{trop}}$, and that image does not contain $\xi$. Hence some
polyhedral neighbourhood $U$ of $\xi$ in $\Trop_f(X)$ satisfies
$f_{\mathrm{trop}}^{-1}(U)\subset V$. Therefore
\[
 U\subset\bigcup_{i=1}^r f_{\mathrm{trop}}(V_{x_i}),
\]
a finite union of polyhedral sets of dimension at most $n-1$. Thus
$\dim U\le n-1$ and $\xi\notin\Trop_f(X)^{(n)}$.

\cref{itm:trop-dimsig:d} If $X$ is boundaryless then $X=\Int(X)$ and
\cref{itm:trop-dimsig:b} gives $\Sigma_f=\Sigma_f^{(n)}$. By
\cref{itm:trop-dimsig:a}, $\Sigma_f^{(n)}$ is the locus where
$\dim_{\mathrm{trop},x}(f)=n$, and $\dim_{\mathrm{trop},x}$ is computed
from the images $f_{\mathrm{trop}}(V)$ alone.
\end{proof}

\begin{warning}{$\Sigma_f$ itself is not a tropical invariant}
Without the boundarylessness hypothesis, $\Sigma_f$ genuinely depends on $f$
and not only on $f_{\mathrm{trop}}$. Let $X=\{x\in\mathbf G_m^{\An}:|T(x)|=1\}$,
of level $1$, and let $f=T$. Then $f_{\mathrm{trop}}\equiv0$ is the same map
as the tropicalization of the constant moment map $1$, yet
$\Sigma_f=\{\eta_0\}$ is the Gauss point while $\Sigma_1=\varnothing$. The
discrepancy lives entirely on $\partial X$: here $X$ is affinoid and
$\eta_0$ is its Shilov boundary point. This is why every subsequent
statement about calibrations is made on $\Sigma_f^{(n)}$, or on interiors of
cells, and never on $\Sigma_f$ as a whole.
\end{warning}

\begin{corollary}
\label{cor:trop-sigma-translation}
Let $X$ be of level $n$ and let $f,f':X\to\mathbf G_m^n$ be moment maps such
that $f'_{\mathrm{trop}}-f_{\mathrm{trop}}$ is constant on $X$. Then
\[
 \Sigma_f\cap\Int(X)=\Sigma_{f'}\cap\Int(X)
 =\Sigma_f^{(n)}\cap\Int(X)=\Sigma_{f'}^{(n)}\cap\Int(X).
\]
\end{corollary}

\begin{proof}
A translation of $\mathbf R^n$ preserves dimensions of germs, so
$\dim_{\mathrm{trop},x}(f)=\dim_{\mathrm{trop},x}(f')$ for every $x$. Now
apply \Cref{thm:trop-dimension-sigma} \cref{itm:trop-dimsig:a} and
\cref{itm:trop-dimsig:b}.
\end{proof}

We can now place the original dimension statement of this chapter, which
concerns cycles rather than spaces, in its context.

\begin{proposition}
\label{prop:analytic-tropical-dimension}
Let $Z$ be an irreducible analytic cycle of dimension $d$ and let $f$ be a
moment map defined near its support. Every tropical image of $Z$ has
dimension at most $d$. For an interior point $x$ of the reduced support,
\[
 \dim_{\mathrm{trop},x}(f)=d
 \quad\Longleftrightarrow\quad
 d_x(f)=d.
\]
Equivalently, the local tropical image has full dimension at $x$ precisely
when the initial forms of the chosen moment functions generate a graded
subfield of $\widetilde{\mathcal H(x)}$ of transcendence degree $d$ over
$\widetilde K$.
\end{proposition}
\begin{proof}
Apply \eqref{eq:trop-dimension-chain} to the reduced support of $Z$:
$\dim_{\mathrm{trop},x}(f)\le d_x(f)\le d_K(x)\le d$.  Together with
\eqref{eq:trop-dimension-sup} on compact analytic domains, this bounds every
cell of the tropical image by $d$. If $x$ is interior, the first inequality
is an equality by \Cref{thm:trop-local-germ}; hence the local image has full
dimension exactly when $d_x(f)=d$. The last formulation is the definition of
$d_x(f)$ in \Cref{def:trop-residual-degree}.
\end{proof}

\begin{remark}
\label{rem:trop-dimension-two-halves}
The full-dimensionality criterion must use graded transcendence degree, not
the rational rank of the value group alone. The identity separating its two
contributions is the graded
Abhyankar equality
\[
 d_K(x)=\operatorname{trdeg}_{\widetilde K}\widetilde{\mathcal H(x)}
 =\operatorname{trdeg}_{\widetilde K_0}\widetilde{\mathcal H(x)}_0
 +\dim_{\mathbf Q}\Bigl(\bigl(|\mathcal H(x)^\times|/|K^\times|\bigr)
 \otimes_{\mathbf Z}\mathbf Q\Bigr),
\]
where the subscript zero denotes the degree-zero part, that is the ordinary
residue field.  The left-hand transcendence degree is the graded one of
\Cref{sec:appG-graded-reduction}, which is what
\Cref{def:trop-residual-degree} uses in defining $d_K(x)$. For curves this is
the equality $d_K(x)=s(x)+t(x)$ of
\Cref{prop:berk-type-invariants}. A type-II point contributes through residue
transcendence while a type-III point contributes through value-group rank,
yet both can support a one-dimensional tropical germ. Neither summand alone
detects tropical dimension. A type-IV point has $d_K(x)=0$ and never lies on a
local skeleton.
\end{remark}

\section{The canonical calibration}
\label{sec:trop-calibration}

A local skeleton is a polyhedral space; to integrate over it we must
calibrate it, that is, attach to each of its facets an $n$-vector in the
sense of \Cref{def:appG-calibration}. The recipe is short to state: project
the moment map to an $n$-dimensional torus so that the projection is
injective on the facet, take the local degree of the resulting finite flat
morphism, and multiply it by the volume vector transported through the
inverse of the projection. Everything in this section is the verification
that the recipe is well posed.

\subsection{Cells and the local degree}

\begin{definition}[Cells of a local skeleton]
\label{def:trop-f-cell}
Let $X$ be of level $n$ and $f:X\to T$ a moment map. An \emph{$f$-cell} is a
compact polyhedral subspace $C\subset\Sigma_f$ such that
$f_{\mathrm{trop}}|_C$ is an isomorphism of polyhedral spaces onto a cell of
$T_{\mathrm{trop}}$. An $f$-cell $C$ of dimension $n$ is \emph{convenient}
if its relative interior $\mathring C$ is open in $\Sigma_f^{(n)}$, is
connected, and does not meet $\partial X$.
\end{definition}

The existence of decompositions of $\Sigma_f^{(n)}$ into convenient
$f$-cells is part of \Cref{thm:trop-sigma-polyhedral}
\cref{itm:trop-poly:top}: $f_{\mathrm{trop}}$ is a piecewise immersion there,
so a sufficiently fine polyhedral decomposition has cells on which it is
injective, and $f_{\mathrm{trop}}(\partial X\cap\Sigma_f^{(n)})$ is
contained in the $(n-1)$-skeleton of a suitable decomposition, by
\Cref{prop:trop-boundary-thin} below. We refer to
\cite[Lemma~14.1.5 and Proposition~14.1.12]{CLD25} for the combinatorial
bookkeeping.

\begin{definition}[Local degree]
\label{def:trop-local-degree}
\index{degree!local, of an analytic morphism}
Let $g:Y\to Z$ be a morphism of $K$-analytic spaces, $y\in Y$, and
$z=g(y)$. If $\dim_yg^{-1}(z)>0$ we set $\deg_y(g)=+\infty$. Otherwise $y$
is isolated in $g^{-1}(z)$, the local ring $\mathcal O_{g^{-1}(z),y}$ is a
finite-dimensional $\mathcal H(z)$-algebra, and we set
\[
 \deg_y(g)=\dim_{\mathcal H(z)}\mathcal O_{g^{-1}(z),y},
 \qquad
 \deg^z(g)=\sum_{y\in g^{-1}(z)}\deg_y(g).
\]
When $g$ is finite and flat these agree with the usual degrees, and
$\deg^z(g)$ is then locally constant in $z$.
\end{definition}

The local degree is insensitive to shrinking: if $V\ni y$ and $U\supset g(V)$
are analytic domains and $g'=g|_V:V\to U$, then $\deg_y(g)=\deg_y(g')$,
because the local rings of $g^{-1}(z)$ and of $(g')^{-1}(z)$ at $y$
coincide. This is what allows the definition to be made cell by cell.

We isolate the analytic inputs, each with an exact reference.

\begin{proposition}
\label{prop:trop-analytic-inputs}
Let $X$ be a $K$-analytic space of level $n$.
\begin{enumerate}
\item\label{itm:trop-input:qf} \emph{(Quasi-finiteness.)} If
  $g:X\to Y$ is a morphism, $x\in X$ and $d_K(g(x))=n=\dim_xX$, then $g$ is
  quasi-finite at $x$: $x$ is isolated in $g^{-1}(g(x))$.
\item\label{itm:trop-input:flat} \emph{(Finiteness and flatness.)} If in
  addition $Y$ is reduced and $x\in\Int(X/Y)$, then $g$ is finite and flat
  at $x$; in particular $\deg_x(g)$ is a positive integer and $g$ induces a
  finite flat morphism of a compact neighbourhood of $x$ onto a
  neighbourhood of $g(x)$, of that degree.
\item\label{itm:trop-input:sc} \emph{(Semicontinuity.)} If $g$ is compact
  and separated, then $y\mapsto\deg^y(g)$ is upper semicontinuous.
\item\label{itm:trop-input:mult} \emph{(Multiplicativity.)} If
  $h:Z\to Y$ and $g:Y\to X$ are morphisms, $z\in Z$ is isolated in its
  fibre for $g\circ h$, and $g^{-1}(g(h(z)))$ is reduced at $h(z)$, then
  $\deg_z(g\circ h)=\deg_{h(z)}(g)\deg_z(h)$.
\end{enumerate}
\end{proposition}

\begin{proof}
These are \cite[Lemma~11.5.1, Lemma~11.6.5, Proposition~13.4.2 and
Proposition~13.3.10]{CLD25}. Assertion \cref{itm:trop-input:qf} is the
fibre-dimension inequality
$\dim_xg^{-1}(g(x))\le\dim_xX-d_K(g(x))$ of Ducros; \cref{itm:trop-input:flat}
uses that the target local ring at a point with $d_K(g(x))=n$ is artinian and,
when reduced, a field; \cref{itm:trop-input:sc} is proved by
Nakayama's lemma after a local factorization of $g$ into a finite morphism
followed by a quasi-\'etale one; \cref{itm:trop-input:mult} is the
composition rule for calibrated degrees, discussed again in
\Cref{sec:trop-graded}.
\end{proof}

\begin{lemma}
\label{lem:trop-quasifinite-flat}
Let $X$ be of level $n$ and $f:X\to\mathbf G_m^n$ a moment map. Then $f$ is
quasi-finite at every point of $\Sigma_f$, and finite and flat at every
point of $\Sigma_f\setminus\partial X$.
\end{lemma}

\begin{proof}
Let $x\in\Sigma_f$. Since the target is already $n$-dimensional, the
equivalence \cref{itm:trop-sigmachar:def} $\Leftrightarrow$
\cref{itm:trop-sigmachar:proj} of
\Cref{thm:trop-sigma-characterization} applies with $I=\{1,\ldots,n\}$ and
gives $f(x)\in S(\mathbf G_m^n)$, so
$d_K(f(x))=n$ by \Cref{prop:trop-torus-skeleton}, and $\dim_xX=n$ by
\Cref{cor:trop-sigma-abhyankar}. \Cref{prop:trop-analytic-inputs}
\cref{itm:trop-input:qf} gives quasi-finiteness. If moreover
$x\notin\partial X$, then $x\in\Int(X)\subset\Int(X/\mathbf G_m^n)$, and
$\mathbf G_m^n$ is reduced, so \cref{itm:trop-input:flat} applies.
\end{proof}

\begin{proposition}
\label{prop:trop-degree-constant}
Let $X$ be of level $n$, let $f:X\to T$ be a moment map, let $C$ be a
convenient $f$-cell of dimension $n$, and let $q:T\to\mathbf G_m^n$ be a
homomorphism of tori such that $q_{\mathrm{trop}}$ is injective on
$f_{\mathrm{trop}}(C)$. Then $q\circ f$ has relative dimension zero at every
point of $C$, and there is a unique integer $d_C(q)>0$ such that $q\circ f$
is finite and flat of degree $d_C(q)$ at every point of $\mathring C$.
\end{proposition}

\begin{proof}
\emph{Step 1: $C\subset\Sigma_{q\circ f}$.} Because $C$ is an $f$-cell of
dimension $n$ and $q_{\mathrm{trop}}$ is injective on
$f_{\mathrm{trop}}(C)$, the composite
$(q\circ f)_{\mathrm{trop}}|_C=q_{\mathrm{trop}}\circ f_{\mathrm{trop}}|_C$
is injective, so $(q\circ f)_{\mathrm{trop}}(C)$ is $n$-dimensional and
$\dim_{\mathrm{trop},x}(q\circ f)=n$ for $x\in C$. Hence
$d_x(q\circ f)=n$ by \eqref{eq:trop-dimension-chain}, so
$(q\circ f)(x)\in S(\mathbf G_m^n)$ by
\Cref{prop:trop-torus-skeleton}, that is $x\in\Sigma_{q\circ f}$.

\emph{Step 2: finiteness and flatness on $\mathring C$.} By
\Cref{lem:trop-quasifinite-flat}, $q\circ f$ is quasi-finite at every point
of $C$ --- this is the assertion on relative dimension --- and, since
$\mathring C\cap\partial X=\varnothing$, finite and flat at every point of
$\mathring C$. Write $\delta(x)=\deg_x(q\circ f)$, a positive integer for
$x\in\mathring C$.

\emph{Step 3: local constancy.} Let $D=\sq\bigl(q_{\mathrm{trop}}
(f_{\mathrm{trop}}(C))\bigr)\subset\mathbf G_m^n$, an $n$-cell of the
canonical skeleton. If $y\in\Int(X)$ satisfies $(q\circ f)(y)\in\mathring D$,
then $q\circ f$ has tropical dimension $n$ at $y$, hence so does $f$, so
$y\in\Sigma_f^{(n)}$ by \Cref{thm:trop-dimension-sigma}
\cref{itm:trop-dimsig:a}. Thus
\[
 \mathring C\ \subset\ (q\circ f)^{-1}(\mathring D)\cap\Int(X)
 \ \subset\ \Sigma_f^{(n)} .
\]
Fix $x\in\mathring C$. By hypothesis $\mathring C$ is open in
$\Sigma_f^{(n)}$, hence open in $(q\circ f)^{-1}(\mathring D)\cap\Int(X)$.
As $q\circ f$ is finite and flat at $x$, there are an open neighbourhood
$U\ni x$ in $\Int(X)$ and an open $W\ni(q\circ f)(x)$ such that $q\circ f$
induces a finite flat morphism $U\to W$ of degree $\delta(x)$ and
$(q\circ f)^{-1}(\mathring D)\cap U\subset\mathring C$. Any
$y\in\mathring C\cap U$ is then the unique preimage in $U$ of
$(q\circ f)(y)$, whence $\delta(y)=\delta(x)$. So $\delta$ is locally
constant on $\mathring C$.

\emph{Step 4: constancy.} $\mathring C$ is connected by the definition of a
convenient cell, so $\delta$ is constant; call the value $d_C(q)$.
\end{proof}

\subsection{Independence of the auxiliary projection}

We now prove that the product
\[
 d_C(q)\cdot\sigma_{q,*}(e)
\]
does not depend on $q$. Here $e=|e_1\wedge\cdots\wedge e_n|$ is the standard
volume $n$-vector of $\mathbf R^n$ and $\sigma_q$ is defined as follows.

\begin{definition}
\label{def:trop-affine-section}
Let $C$ be an $n$-dimensional $f$-cell, put
$C_f=f_{\mathrm{trop}}(C)\subset T_{\mathrm{trop}}$, and let
$\langle C_f\rangle$ be its affine span, an $n$-dimensional
$\mathbf Q$-rational affine subspace of $N_{\mathbf R}$. For a homomorphism
of tori $q:T\to\mathbf G_m^n$ with $q_{\mathrm{trop}}$ injective on $C_f$,
the restriction $q_{\mathrm{trop}}|_{\langle C_f\rangle}$ is an affine
isomorphism onto $\mathbf R^n$; its inverse
\[
 \sigma_q:\mathbf R^n\longrightarrow\langle C_f\rangle\subset N_{\mathbf R}
\]
is the unique affine section of $q_{\mathrm{trop}}$ with image
$\langle C_f\rangle$. We write $\sigma_{q,*}(e)$ for the image of $e$ under
$\bigwedge^n$ of the linear part of $\sigma_q$; it is an $n$-vector on
$\langle C_f\rangle$.
\end{definition}

The rationality of $\langle C_f\rangle$ is what makes the argument work, and
it holds because $C_f$ is a rational polyhedron: its affine span is cut out
by integral linear equations.

The following theorem is the reason the word \emph{canonical} is used. Its
proof is a two-case reduction that we carry out in full, because the pattern
--- separate a monomial change of projection from a translation, then
interpolate by a rational linear map --- recurs in
\Cref{sec:trop-refinement}.

\begin{theorem}
\label{thm:trop-calibration-canonical}
Let $X$ be of level $n$, let $f:X\to T$ be a moment map, let $C$ be a
convenient $f$-cell of dimension $n$, and let $q,q':T\to\mathbf G_m^n$ be
homomorphisms of tori such that $q_{\mathrm{trop}}$ and
$q'_{\mathrm{trop}}$ are both injective on $C_f=f_{\mathrm{trop}}(C)$. Then
\begin{equation}
\label{eq:trop-calibration-canonical}
 d_C(q)\,\sigma_{q,*}(e)=d_C(q')\,\sigma_{q',*}(e)
\end{equation}
as $n$-vectors on $\langle C_f\rangle$.
\end{theorem}

\begin{proof}
\emph{Case 1: $q'=\phi\circ q$ with $\phi:\mathbf G_m^n\to\mathbf G_m^n$ a
surjective homomorphism of tori.} Let $A\in M_n(\mathbf Z)$ be the matrix of
$\phi$, so $\delta:=|\det A|=\deg\phi$ by
\Cref{lem:trop-monomial-degree} \cref{itm:trop-mono:finflat}. For
$x\in\mathring C$ we have $(q\circ f)(x)\in S(\mathbf G_m^n)$ by Step 1 of
the proof of \Cref{prop:trop-degree-constant}, so
\Cref{lem:trop-monomial-degree} \cref{itm:trop-mono:degree} gives
$\deg_{(q\circ f)(x)}(\phi)=\delta$ and tells us that the fibre of $\phi$
through that point is reduced. \Cref{prop:trop-analytic-inputs}
\cref{itm:trop-input:mult} therefore applies and yields
\[
 d_C(q')=\deg_x(\phi\circ q\circ f)
 =\deg_{(q\circ f)(x)}(\phi)\cdot\deg_x(q\circ f)
 =\delta\,d_C(q).
\]
On the tropical side, $\phi_{\mathrm{trop}}=A$ is a linear automorphism of
$\mathbf R^n$ with $|\det A|=\delta$, and
$q'_{\mathrm{trop}}=A\circ q_{\mathrm{trop}}$, so the affine section with
image $\langle C_f\rangle$ satisfies
$\sigma_{q'}=\sigma_q\circ A^{-1}$. Hence
\[
 \sigma_{q',*}(e)=\sigma_{q,*}\bigl((A^{-1})_*e\bigr)
 =\delta^{-1}\,\sigma_{q,*}(e),
\]
since $\bigwedge^n$ of $A^{-1}$ multiplies the volume vector by
$|\det A^{-1}|=\delta^{-1}$. The two factors are reciprocal, so
\eqref{eq:trop-calibration-canonical} holds.

\emph{Case 2: $q'_{\mathrm{trop}}-q_{\mathrm{trop}}$ is constant on $C_f$.}
Then $\sigma_q$ and $\sigma_{q'}$ differ by a translation of
$\mathbf R^n$, so they have the same linear part and
$\sigma_{q,*}(e)=\sigma_{q',*}(e)$. It remains to see that
$d_C(q)=d_C(q')$. Both $q\circ f$ and $q'\circ f$ are moment maps
$V\to\mathbf G_m^n$ on the analytic domain
$V=f_{\mathrm{trop}}^{-1}(C_f)$, whose tropicalizations differ by a
constant on $V$, and $\mathring C\subset\Int(V)$. Now apply
\Cref{thm:trop-degree-translation} below.

\emph{General case.} Let $\psi$ be the unique linear self-map of
$\mathbf R^n$ such that $\psi\circ q_{\mathrm{trop}}$ and
$q'_{\mathrm{trop}}$ differ by a translation on $C_f$; it exists and is
invertible because both $q_{\mathrm{trop}}$ and $q'_{\mathrm{trop}}$ restrict
to affine isomorphisms $\langle C_f\rangle\to\mathbf R^n$, and one may take
$\psi$ to be the linear part of
$q'_{\mathrm{trop}}|_{\langle C_f\rangle}\circ
 \bigl(q_{\mathrm{trop}}|_{\langle C_f\rangle}\bigr)^{-1}$.
Since $\langle C_f\rangle$ is $\mathbf Q$-rational and both restrictions are
defined by integral matrices, $\psi\in\mathrm{GL}_n(\mathbf Q)$. Choose an
integer $\nu\ge1$ clearing the denominators of $\psi$, and let
$\phi,\phi':\mathbf G_m^n\to\mathbf G_m^n$ be the surjective homomorphisms
with $\phi_{\mathrm{trop}}=\nu\psi$ and
$\phi'_{\mathrm{trop}}=\nu\,\mathrm{id}$; both matrices are integral, and
$\phi'_{\mathrm{trop}}\circ\psi=\phi_{\mathrm{trop}}$.

Consider the four projections $q$, $\phi\circ q$, $q'$, $\phi'\circ q'$. By
construction
\[
 (\phi'\circ q')_{\mathrm{trop}}
 =\nu\,q'_{\mathrm{trop}}
 =\nu\psi\circ q_{\mathrm{trop}}+\text{const}
 =(\phi\circ q)_{\mathrm{trop}}+\text{const}
\]
on $C_f$. Case 1 gives
\[
 \begin{aligned}
  d_C(\phi\circ q)\,\sigma_{\phi\circ q,*}(e)
   &=d_C(q)\,\sigma_{q,*}(e),\\
  d_C(\phi'\circ q')\,\sigma_{\phi'\circ q',*}(e)
   &=d_C(q')\,\sigma_{q',*}(e),
 \end{aligned}
\]
and Case 2, applied to the pair $\phi\circ q$ and $\phi'\circ q'$, gives
\[
 d_C(\phi\circ q)\,\sigma_{\phi\circ q,*}(e)
 =d_C(\phi'\circ q')\,\sigma_{\phi'\circ q',*}(e).
\]
Chaining the three equalities proves
\eqref{eq:trop-calibration-canonical}.
\end{proof}

The one ingredient of the proof that we did not establish is the invariance
of the local degree under a translation of the tropicalization. It is a
genuine theorem, and it is where the analytic geometry, rather than the
combinatorics, does the work.

\begin{theorem}
\label{thm:trop-degree-translation}
Let $X$ be a $K$-analytic space of level $n$ and let
$f,f':X\to\mathbf G_m^n$ be moment maps such that
$f'_{\mathrm{trop}}-f_{\mathrm{trop}}$ is constant on $X$. Then $f$ and
$f'$ are finite and flat at every point $x$ of
$\Sigma_f\cap\Int(X)=\Sigma_{f'}\cap\Int(X)$, and
\[
 \deg_x(f)=\deg_x(f')\qquad\text{for every such }x.
\]
\end{theorem}

\begin{proof}
The identification of the two skeleta is
\Cref{cor:trop-sigma-translation}, and finiteness and flatness are
\Cref{lem:trop-quasifinite-flat}. The equality of degrees is
\cite[Theorem~14.2.7]{CLD25}. We indicate the mechanism, which is
instructive. One first reduces to $K$ algebraically closed, using that a
point of $S(\mathbf G_m^n)$ has a unique preimage in
$S(\mathbf G_{m,\widehat{K^{\mathrm{alg}}}}^n)$ --- visible on the defining
formula \eqref{eq:trop-gauss-seminorm}, which is invariant under
$\mathrm{Aut}(\widehat{K^{\mathrm{alg}}}/K)$ --- together with the base-change
additivity of degrees. Over an algebraically closed field, the local
constant-norm theorem \cite[Proposition~12.2.12]{CLD25}, applied to
$f_i'/f_i$, gives
$f_i'=(1+g_i)\lambda_if_i$ with $\lambda_i\in K^\times$ and $|g_i|<1$; after
translating $f$ one may take $\lambda_i=1$. One then interpolates by the
family
\[
 h_t=\bigl((1+tg_1)f_1,\ldots,(1+tg_n)f_n\bigr),\qquad t\in D(0,R),
\]
with $R>1$ chosen so that $R|g_i|<1$, on a Cohen--Macaulay affinoid
neighbourhood of $x$. Since $R|g_i|<1$, the tropicalization of $h_t$ is
independent of $t$; one checks that $h$ is finite and flat over an open
polydisc-shaped piece of the canonical skeleton of
$\mathbf G_m^n\times D(0,R)$, and the degree of a finite flat morphism is
constant over a connected base. Specializing to $t=0$ and $t=1$ gives
$\deg_x(f)=\deg_x(f')$.
\end{proof}

\begin{warning}{The two factors move in opposite directions}
Case 1 of \Cref{thm:trop-calibration-canonical} is the whole normalization
story of this book in one line. Enlarging the projection by a monomial map of
degree $\delta$ multiplies the analytic degree by $\delta$ and divides the
lattice normalization of the section by $\delta$. A construction that keeps
only one of the two factors --- for instance one that defines the weight as
the analytic degree of a fixed coordinate projection, or one that defines it
as a lattice index and forgets the analytic geometry --- is not invariant and
will contradict the projection formula. \Cref{ex:appG-index-test} and
\Cref{ex:trop-two-charts} are the two smallest tests of this.
\end{warning}

\subsection{The canonical calibration of a local skeleton}

\begin{definition}[Canonical calibration]
\label{def:trop-canonical-calibration}
\index{calibration!canonical}
Let $X$ be of level $n$, let $f:V\to T$ be a moment map on an analytic
domain, and let $\Sigma\subset\Sigma_f$ be a paracompact polyhedral
subspace. Choose a decomposition $\mathscr C$ of $\Sigma$ into convenient
$f$-cells. For an $n$-dimensional $C\in\mathscr C$ choose $q$ as in
\Cref{prop:trop-degree-constant} and set
\begin{equation}
\label{eq:trop-canonical-calibration}
 \mu_C=d_C(q)\,\bigl(f_{\mathrm{trop}}|_C\bigr)^*\sigma_{q,*}(e),
\end{equation}
an $n$-vector on $C$; by \Cref{thm:trop-calibration-canonical} it does not
depend on $q$. The family $(\mu_C)$ is the \emph{canonical calibration}
$\mu_{\Sigma,f,\mathscr C}$ of $\Sigma$ relative to the data
$(V,f,\mathscr C)$.
\end{definition}

Three checks remain: that $\mu_{\Sigma,f,\mathscr C}$ is compatible with
subdivision of $\mathscr C$, that it does not depend on $(V,f)$, and that
the resulting local objects glue. The first is immediate and we do it now;
the other two occupy \Cref{sec:trop-refinement}.

\begin{lemma}
\label{lem:trop-calibration-subdivision}
If $\mathscr C'$ subdivides $\mathscr C$, then
$\mu_{\Sigma,f,\mathscr C'}=\mu_{\Sigma,f,\mathscr C}$.
\end{lemma}

\begin{proof}
Let $C'\in\mathscr C'$ be $n$-dimensional and let $C\in\mathscr C$ be the
cell with $\mathring{C'}\subset\mathring C$; such a $C$ exists and is unique
because the relative interiors of the $n$-cells of $\mathscr C$ partition
those of $\mathscr C'$ up to lower-dimensional sets. The affine spans agree,
$\langle f_{\mathrm{trop}}(C')\rangle=\langle f_{\mathrm{trop}}(C)\rangle$,
so a projection $q$ admissible for $C$ is admissible for $C'$ and gives the
same $\sigma_q$. Moreover $d_{C'}(q)=d_C(q)$ because both are the value of
$x\mapsto\deg_x(q\circ f)$ on the nonempty open subset
$\mathring{C'}\subset\mathring C$, on which that function is constant by
\Cref{prop:trop-degree-constant}. Hence $\mu_{C'}=\mu_C|_{C'}$.
\end{proof}

\section{Independence under refinement of tropical charts}
\label{sec:trop-refinement}

This is the technical heart of the chapter. We must show that the
calibration of \Cref{def:trop-canonical-calibration} depends only on the
polyhedral subspace $\Sigma\subset X$, not on the moment map used to see it.
The obstacle is that two moment maps have nothing in common except their
product, and that the product lives on a torus of larger dimension, where
the projections $q$ available are different. The proof is a cell-by-cell
comparison with the product, using the freedom in
\Cref{thm:trop-calibration-canonical} to choose $q$ conveniently.

\begin{theorem}
\label{thm:trop-calibration-chart-independent}
Let $X$ be a topologically separated $K$-analytic space of level $n$, and
let $\Sigma\subset X$ be a paracompact polyhedral subspace. Suppose given
two sets of data
\[
 (V,f:V\to T,\mathscr C),\qquad (V',f':V'\to T',\mathscr C'),
\]
where $V,V'$ are analytic domains containing $\Sigma$, $f,f'$ are moment
maps with $\Sigma\subset\Sigma_f$ and $\Sigma\subset\Sigma_{f'}$, and
$\mathscr C,\mathscr C'$ are decompositions of $\Sigma$ into convenient
$f$-, respectively $f'$-cells. Then
\[
 \mu_{\Sigma,f,\mathscr C}=\mu_{\Sigma,f',\mathscr C'}.
\]
\end{theorem}

\begin{proof}
Put $W=V\cap V'$, an analytic domain containing $\Sigma$, and let
\[
 g=(f|_W,f'|_W):W\longrightarrow T\times T'
\]
be the product moment map (\Cref{lem:product-moment-map}). By
\Cref{lem:trop-refinement-calculus} \cref{itm:trop-ref:sigma},
$d_x(g)\ge\max\{d_x(f),d_x(f')\}=n$ for $x\in\Sigma$, so
$\Sigma\subset\Sigma_g$ by \Cref{thm:trop-sigma-characterization}. Choose a
decomposition $\mathscr D$ of $\Sigma$ into convenient $g$-cells refining
both $\mathscr C$ and $\mathscr C'$; this is possible because the three
decompositions are polyhedral and admit a common refinement
(\Cref{prop:tropical-common-refinement}), and a subdivision of a convenient
cell into cells with connected interiors again consists of convenient cells.
By \Cref{lem:trop-calibration-subdivision} it suffices to prove
\begin{equation}
\label{eq:trop-chart-indep-goal}
 \mu_{\Sigma,f,\mathscr D}=\mu_{\Sigma,g,\mathscr D},
\end{equation}
since the same argument applied to $f'$ then gives
$\mu_{\Sigma,f',\mathscr D}=\mu_{\Sigma,g,\mathscr D}$, and both sides of
the theorem equal $\mu_{\Sigma,g,\mathscr D}$ after subdivision.

Fix an $n$-dimensional cell $D\in\mathscr D$ and let $C\in\mathscr C$ be the
cell containing $D$. Let $p:T\times T'\to T$ be the first projection, so
that $p\circ g=f|_W$ and $p_{\mathrm{trop}}\circ g_{\mathrm{trop}}
=f_{\mathrm{trop}}|_W$.

\emph{Step 1: an admissible projection for both.} Choose
$q:T\to\mathbf G_m^n$ with $q_{\mathrm{trop}}$ injective on
$f_{\mathrm{trop}}(D)$; this is possible because
$\langle f_{\mathrm{trop}}(D)\rangle$ is a $\mathbf Q$-rational
$n$-dimensional affine subspace of $N_{\mathbf R}$. Then
$(q\circ p)_{\mathrm{trop}}=q_{\mathrm{trop}}\circ p_{\mathrm{trop}}$ is
injective on $g_{\mathrm{trop}}(D)$, because
$g_{\mathrm{trop}}|_D$ is injective with
$p_{\mathrm{trop}}(g_{\mathrm{trop}}(D))=f_{\mathrm{trop}}(D)$ and the
composite $(q\circ p)_{\mathrm{trop}}\circ g_{\mathrm{trop}}|_D
=q_{\mathrm{trop}}\circ f_{\mathrm{trop}}|_D$ is injective. So $q\circ p$ is
admissible for $D$ as a $g$-cell, and $q$ is admissible for $D$ viewed
inside the $f$-cell $C$.

\emph{Step 2: the two degrees coincide.} By construction
$(q\circ p)\circ g=q\circ f|_W$ as morphisms $W\to\mathbf G_m^n$. Now
$\mathring D$ does not meet $\partial W$, because $D$ is a convenient
$g$-cell; hence $\mathring D\subset\Int(W/V)$ and, for $x\in\mathring D$,
the local degree of $q\circ f$ at $x$ computed on $W$ equals the one
computed on $V$ (the local rings of the fibres coincide, as recorded after
\Cref{def:trop-local-degree}). \Cref{prop:trop-degree-constant} applied to
the $g$-cell $D$ and the projection $q\circ p$ gives an integer $d_D>0$ with
$\deg_x\bigl((q\circ p)\circ g\bigr)=d_D$ on $\mathring D$; applied to the
$f$-cell $C$ and the projection $q$ it gives
$\deg_x(q\circ f)=d_C(q)$ on $\mathring C\supset\mathring D$. Since the two
morphisms are equal, $d_D=d_C(q)$.

\emph{Step 3: the two sections coincide.} The affine section
$\sigma_{q\circ p}$ of $(q\circ p)_{\mathrm{trop}}$ with image
$\langle g_{\mathrm{trop}}(D)\rangle$ satisfies
$p_{\mathrm{trop}}\circ\sigma_{q\circ p}=\sigma_q$, because both sides are
affine sections of $q_{\mathrm{trop}}$ with image
$\langle f_{\mathrm{trop}}(D)\rangle=\langle f_{\mathrm{trop}}(C)\rangle$,
and such a section is unique. Pulling back along the isomorphisms
$g_{\mathrm{trop}}|_D$ and $f_{\mathrm{trop}}|_D$ and using
$f_{\mathrm{trop}}|_D=p_{\mathrm{trop}}\circ g_{\mathrm{trop}}|_D$ therefore
gives
\[
 \bigl(g_{\mathrm{trop}}|_D\bigr)^*\sigma_{q\circ p,*}(e)
 =\bigl(f_{\mathrm{trop}}|_D\bigr)^*\sigma_{q,*}(e)
\]
as $n$-vectors on $D$. Indeed, both are the unique $n$-vector on $D$ whose
image under $\bigwedge^n$ of $(q\circ f)_{\mathrm{trop}}|_D$ is $e$.

\emph{Conclusion.} Combining Steps 2 and 3,
\[
 \mu^{g}_D=d_D\,\bigl(g_{\mathrm{trop}}|_D\bigr)^*\sigma_{q\circ p,*}(e)
 =d_C(q)\,\bigl(f_{\mathrm{trop}}|_D\bigr)^*\sigma_{q,*}(e)
 =\mu^{f}_C\bigl|_D,
\]
which is \eqref{eq:trop-chart-indep-goal}.
\end{proof}

\begin{remark}
\label{rem:trop-chart-indep-anatomy}
The proof contains exactly one nontrivial idea, and it is worth naming.
The product moment map $g$ has a larger target, so a priori it admits more
projections $q$; the point is that among them one may choose
$q\circ p$, which \emph{factors through the projection to the smaller
torus}. For that particular choice, the composite $(q\circ p)\circ g$ is
literally the same analytic morphism as $q\circ f$, so no comparison of
degrees is needed at all --- Step 2 is a tautology once the domains are
matched --- and the entire content is transferred to
\Cref{thm:trop-calibration-canonical}, which frees us to make that choice.
This is the precise sense in which the canonicity theorem is logically prior
to the refinement theorem, and why the two must not be proved in the other
order.
\end{remark}

\begin{corollary}
\label{cor:trop-calibration-restriction}
With the notation of \Cref{def:trop-canonical-calibration}, if
$\Sigma'\subset\Sigma$ is a polyhedral subspace, then
$\mu_\Sigma|_{\Sigma'}=\mu_{\Sigma'}$. If $\Sigma$ is contained in an
analytic domain $V\subset X$, the canonical calibration of $\Sigma$ computed
in $V$ agrees with the one computed in $X$.
\end{corollary}

\begin{proof}
Choose compatible decompositions of $\Sigma$ and $\Sigma'$ and pass to a
common subdivision so that every cell of $\Sigma'$ is the trace of a cell of
$\Sigma$. Then
\Cref{def:trop-canonical-calibration} computes $\mu_{\Sigma'}$ from the
cells contained in $\Sigma'$, with the same $q$ and the same degrees, and
\Cref{lem:trop-calibration-subdivision} removes the dependence on the
decomposition. For the second assertion, note that
$\Sigma_{f|_V}=\Sigma_f\cap V$ (\Cref{thm:trop-sigma-polyhedral}
\cref{itm:trop-poly:union}) and that local degrees are computed in local
rings, hence unchanged by restriction to a domain.
\end{proof}

\begin{corollary}
\label{cor:trop-calibration-glue}
Let $X$ be a topologically separated $K$-analytic space of level $n$ and let
$\Sigma\subset X$ be a polyhedral subspace, that is, a subset which is
$G$-locally a closed polyhedral subspace of a local skeleton. Then there is a
unique calibration $\mu_\Sigma$ of $\Sigma$ such that for every analytic
domain $V$, every moment map $f$ on $V$, and every paracompact polyhedral
subspace $\Sigma'\subset\Sigma\cap\Sigma_f$, one has
$\mu_\Sigma|_{\Sigma'}=\mu_{\Sigma'}$ in the sense of
\Cref{def:trop-canonical-calibration}. In particular every $K$-analytic
space of pure dimension $n$ has compatible calibrated tropical charts: each
$n$-dimensional polyhedral subspace, and in particular each local skeleton,
carries its canonical calibration. The calculus of
\Cref{chap:tropical-spaces,chap:tropical-integration} therefore applies
chartwise; the analytic space itself is not being identified with a
polyhedral tropical space.
\end{corollary}

\begin{proof}
By definition of a polyhedral subspace, $\Sigma$ is $G$-covered by
subspaces $\Sigma'$ of the stated form. Calibrations form a $G$-sheaf
(\Cref{def:appG-calibration}), so there is at most one such $\mu_\Sigma$.
For existence one must check the cocycle condition on overlaps: if
$\Sigma'$ and $\Sigma''$ are two such subspaces, then
$\Sigma'\cap\Sigma''$ is again one, and
\Cref{cor:trop-calibration-restriction} gives
\[
 \mu_{\Sigma'}|_{\Sigma'\cap\Sigma''}
 =\mu_{\Sigma'\cap\Sigma''}
 =\mu_{\Sigma''}|_{\Sigma'\cap\Sigma''},
\]
where the outer equalities use \Cref{cor:trop-calibration-restriction} and
the inner one uses \Cref{thm:trop-calibration-chart-independent} to identify
the two computations of $\mu_{\Sigma'\cap\Sigma''}$, made with the moment
maps of $\Sigma'$ and of $\Sigma''$ respectively.
\end{proof}

\begin{proposition}
\label{prop:trop-calibration-base-change}
Let $L/K$ be a complete valued extension, $p:X_L\to X$ the base-change
morphism, $\Sigma\subset X$ a polyhedral subspace, and $\Sigma_L\subset X_L$
its preimage. Then $\Sigma_L$ is a polyhedral subspace of $X_L$ and
$p_*\mu_{\Sigma_L}=\mu_\Sigma$.
\end{proposition}

\begin{proof}
We may assume $X$ compact and $\Sigma=\Sigma_f$ for a moment map
$f:X\to\mathbf G_m^n$; then $\Sigma_L=\Sigma_{f_L}$ by
\Cref{lem:trop-moment-basic} \cref{itm:trop-moment:base} and
\Cref{thm:trop-sigma-polyhedral} \cref{itm:trop-poly:union}. Choose adapted
decompositions $\mathscr C$ of $\Sigma$ and $\mathscr D$ of $\Sigma_L$ into
convenient cells so that $p$ maps cells to cells. Let $C$ be an $n$-cell of
$\mathscr C$, let $q$ be admissible for it, let $x\in\mathring C$ and
$t=(q\circ f)(x)\in S(\mathbf G_m^n)$. By
\Cref{prop:trop-torus-skeleton} the point $t$ has a unique preimage
$t_L\in S(\mathbf G_{m,L}^n)$, since both are described by the same formula
\eqref{eq:trop-gauss-seminorm}. Base-change additivity of degrees
(\Cref{prop:trop-analytic-inputs}, and \cite[Lemma~13.3.3]{CLD25}) gives
\[
 \deg_x(q\circ f)=\sum_{y\in p^{-1}(x)\cap(q\circ f_L)^{-1}(t_L)}
 \deg_y(q\circ f_L)
 =\sum_{y\in p^{-1}(x)\cap\Sigma_{f_L}}\deg_y(q\circ f_L),
\]
the second equality because any $y\in\Sigma_{f_L}$ over $x$ has
$(q\circ f_L)_{\mathrm{trop}}(y)=(q\circ f)_{\mathrm{trop}}(x)=\trop(t)$ and
$(q\circ f_L)(y)\in S(\mathbf G_{m,L}^n)$, hence $(q\circ f_L)(y)=t_L$ by
uniqueness. Since $q_{\mathrm{trop}}$ and the sections $\sigma_q$ are
unchanged by base change, summing the cellwise identities gives
$p_*\mu_{\Sigma_L}=\mu_\Sigma$.
\end{proof}

\section{Calibrations and tropical multiplicities}
\label{sec:trop-multiplicities}

The calibration of \Cref{sec:trop-calibration} lives upstairs, on the local
skeleton inside $X$. What tropical geometry manipulates is its image
downstairs, in $N_{\mathbf R}$. Pushing forward is where the several cells
lying over one polyhedron get added, and where the multiplicities of
\Cref{chap:tropical-spaces} appear.

\begin{definition}[Tropical polyhedron of a moment map]
\label{def:trop-tropical-polyhedron}
\index{tropical polyhedron}
A moment map $f:X\to T$ on a space of level $n$ is \emph{skeletally proper}
if $f_{\mathrm{trop}}$ restricted to $\Sigma_f^{(n)}$ is topologically
proper; this holds automatically when $X$ is compact. For such an $f$, the
\emph{tropical polyhedron} of $f$ is
\[
 \Pi_f=f_{\mathrm{trop}}\bigl(\Sigma_f^{(n)}\bigr)\subset T_{\mathrm{trop}},
\]
equipped with the pushforward $f_{\mathrm{trop},*}\mu_{\Sigma_f^{(n)}}$ of
the canonical calibration.
\end{definition}

Pushforward is legitimate here: $f_{\mathrm{trop}}$ restricted to
$\Sigma_f^{(n)}$ is a piecewise immersion by
\Cref{thm:trop-sigma-polyhedral} \cref{itm:trop-poly:top}, and it is proper
by hypothesis, so the pushforward calibration of
\Cref{prop:appG-calibration-tests} is defined. When $X$ is compact,
$\Pi_f=\Trop_f(X)^{(n)}$ by \Cref{thm:trop-dimension-sigma}
\cref{itm:trop-dimsig:c}: the tropical polyhedron is exactly the
top-dimensional part of the naive tropical image.

\begin{theorem}
\label{thm:trop-multiplicity-formula}
Let $X$ be of level $n$ and $f:X\to T$ a skeletally proper moment map. Let
$\mathscr C$ be a decomposition of $\Sigma_f^{(n)}$ into convenient
$f$-cells such that $\mathscr D=\{f_{\mathrm{trop}}(C):C\in\mathscr C\}$ is a
polyhedral decomposition of $\Pi_f$. Let $D\in\mathscr D$ be
$n$-dimensional, let $q:T\to\mathbf G_m^n$ have $q_{\mathrm{trop}}|_D$
injective, and let $\sigma_D:\mathbf R^n\to\langle D\rangle$ be the affine
section of $q_{\mathrm{trop}}$ with image $\langle D\rangle$. Then
\begin{equation}
\label{eq:trop-multiplicity-formula}
 \mu_{\Pi_f,D}=d_D\,\sigma_{D,*}(e),
 \qquad
 d_D=\sum_{\substack{C\in\mathscr C\\ f_{\mathrm{trop}}(C)=D}}d_C(q),
\end{equation}
where $d_C(q)$ is the constant value of $x\mapsto\deg_x(q\circ f)$ on
$\mathring C$. Writing $e_D$ for the primitive integral volume vector of the
tangent lattice $N_D$, the tropical weight of $D$ is the positive real
number $m_D$ with $\mu_{\Pi_f,D}=m_De_D$.
\end{theorem}

\begin{proof}
By \Cref{def:trop-canonical-calibration} the calibration of a cell
$C\in\mathscr C$ with $f_{\mathrm{trop}}(C)=D$ is
$\mu_C=d_C(q)(f_{\mathrm{trop}}|_C)^*\sigma_{D,*}(e)$, the projection $q$
being admissible for $C$ precisely because it is admissible for
$D=f_{\mathrm{trop}}(C)$. The pushforward of a calibration under a piecewise
immersion is, by definition, the sum over the cells lying above a given
target cell of the images of their calibrations
(\Cref{prop:appG-calibration-tests}); since $f_{\mathrm{trop}}|_C$ is an
isomorphism onto $D$, the image of $\mu_C$ is $d_C(q)\sigma_{D,*}(e)$.
Summing gives \eqref{eq:trop-multiplicity-formula}. Positivity of $m_D$ is
positivity of each $d_C(q)$, which is \Cref{prop:trop-degree-constant}.
\end{proof}

\begin{proposition}
\label{prop:trop-multiplicity-local}
Let $X$ be of level $n$, let $f:X\to T$ be a skeletally proper moment map,
and let $x\in\Pi_f$ have an open neighbourhood $V$ in $\Pi_f$ which is an
$n$-cell of $T_{\mathrm{trop}}$. Let $q:T\to\mathbf G_m^n$ have
$q_{\mathrm{trop}}|_V$ injective, and suppose that the restriction of
$q\circ f$ to $f_{\mathrm{trop}}^{-1}(V)$ is finite and flat over
$\sq(q_{\mathrm{trop}}(x))$, of degree $d_x$. Then the canonical calibration
of $\Pi_f$ is constant equal to $d_x\sigma_{q,*}(e)$ near $x$.
\end{proposition}

\begin{proof}
Let $D\in\mathscr D$ be an $n$-cell containing $x$; then
$\langle D\rangle=\langle V\rangle$, and $q_{\mathrm{trop}}$ is injective on
$D$. By \Cref{thm:trop-multiplicity-formula} it suffices to prove
$d_D=d_x$.

Shrink to an open $W\ni\sq(q_{\mathrm{trop}}(x))$ over which
$q\circ f|_{f_{\mathrm{trop}}^{-1}(V)}$ is finite and flat of degree $d_x$,
and choose an $n$-cell $\Delta\subset D$ containing $x$ with
$\sq(q_{\mathrm{trop}}(\Delta))\subset W$. Fix
$y\in\mathring\Delta$. We claim that the preimages of
$\sq(q_{\mathrm{trop}}(y))$ under $q\circ f|_{f_{\mathrm{trop}}^{-1}(V)}$
are exactly the preimages of $y$ in $\Sigma_f^{(n)}$ under
$f_{\mathrm{trop}}$.

If $z\in f_{\mathrm{trop}}^{-1}(V)$ has
$(q\circ f)(z)=\sq(q_{\mathrm{trop}}(y))$, then applying $\trop$ and using
injectivity of $\sq\circ q_{\mathrm{trop}}$ on $V$ gives
$f_{\mathrm{trop}}(z)=y$; moreover $z\in\Sigma_f$ by
\eqref{eq:trop-sigma-def}, and $q\circ f$ is open near $z$ because it is
finite and flat there, so the image of any neighbourhood of $z$ contains a
neighbourhood of $\sq(q_{\mathrm{trop}}(y))$, whence
$\dim_{\mathrm{trop},z}(q\circ f)=n$ and a fortiori
$\dim_{\mathrm{trop},z}(f)=n$; thus $z\in\Sigma_f^{(n)}$ by
\Cref{thm:trop-dimension-sigma} \cref{itm:trop-dimsig:a}. Conversely, if
$z\in\Sigma_f^{(n)}$ has $f_{\mathrm{trop}}(z)=y$, let $C\in\mathscr C$ be
the cell with $z\in\mathring C$; then $f_{\mathrm{trop}}$ maps $C$
isomorphically onto $D$ and $q\circ f$ has tropical dimension $n$ at $z$, so
$(q\circ f)(z)\in S(\mathbf G_m^n)$; as its tropicalization is
$q_{\mathrm{trop}}(y)$, it equals $\sq(q_{\mathrm{trop}}(y))$.

Summing local degrees over the common set of preimages and using
\Cref{prop:trop-degree-constant} on each cell gives
$d_x=\sum_Cd_C(q)=d_D$.
\end{proof}

\begin{corollary}
\label{cor:trop-weights-integral}
In the situation of \Cref{thm:trop-multiplicity-formula}, let $N_D$ be the
tangent lattice of $D$ in $N$ and let $e_D$ be its primitive volume vector.
Then $q_{\mathrm{trop}}(N_D)\subset\mathbf Z^n$ has finite index and
\begin{equation}
\label{eq:trop-weight-index}
 m_D=\frac{d_D}{\bigl[\mathbf Z^n:q_{\mathrm{trop}}(N_D)\bigr]} .
\end{equation}
Since $N_D$ is saturated in $N$, there exist admissible $q$ with
$q_{\mathrm{trop}}(N_D)=\mathbf Z^n$; for any such choice
$m_D=d_D\in\mathbf Z_{>0}$.
\end{corollary}

\begin{proof}
Because $\sigma_q$ is a section of $q_{\mathrm{trop}}$ and
$q_{\mathrm{trop}}|_{N_{D,\mathbf R}}$ is injective, the lattice
$\Lambda=\sigma_q(\mathbf Z^n)$ contains $N_D$, with
$[\Lambda:N_D]=[\mathbf Z^n:q_{\mathrm{trop}}(N_D)]=:j$. Taking
$\bigwedge^n$ of an inclusion of lattices of index $j$ multiplies the
primitive volume vector by $j$, so $e_D=j\,\sigma_{q,*}(e)$, that is
$\sigma_{q,*}(e)=j^{-1}e_D$. Substituting in
\eqref{eq:trop-multiplicity-formula} gives
\eqref{eq:trop-weight-index}.

For the last assertion, $N_D=N\cap\operatorname{span}_{\mathbf R}(D-D)$ is
saturated, so $N/N_D$ is free and the exact sequence
$0\to N_D\to N\to N/N_D\to0$ splits; the resulting projection $N\to N_D$,
composed with an isomorphism $N_D\simeq\mathbf Z^n$, is an integral linear
map, hence the tropicalization of a homomorphism of tori
$q:T\to\mathbf G_m^n$, and it is injective on $N_{D,\mathbf R}$, hence on
$\langle D\rangle$. For this $q$ one has $j=1$. Positivity and integrality
of $d_D$ are \Cref{prop:trop-degree-constant}.
\end{proof}

\begin{warning}{The scalar weight is not the invariant}
In the integral setting of \Cref{cor:trop-weights-integral}, the tangent
lattice is the saturated lattice
$N_D=N\cap\operatorname{span}_{\mathbf R}(D-D)$ and $m_D$ is a positive
integer. What is invariant under changes of chart and subdivision is the
calibrated volume vector $\mu_D=m_De_D$, not a coefficient written relative
to an arbitrary auxiliary lattice. Re-expressing the same vector relative
to a non-saturated auxiliary lattice can produce a rational coefficient,
but that coefficient is not the canonical integral weight. Likewise a
nonreduced source contributes its integral generic length to the fundamental
cycle; the rational calibrated degree of the reduction map in
\Cref{prop:trop-calibrated-degree} \cref{itm:trop-dc:red} compensates for
that length and does not make the canonical space weight fractional.
Genuinely rational weights enter only after rational coefficients are
allowed on analytic cycles, as in \Cref{chap:cld-currents}. This is why
\Cref{def:appG-calibration} records $\mu_D$ rather than a bare number.
\end{warning}

We can now restate the multiplicity in the language of cycles, which is the
form in which \Cref{chap:cld-currents} uses it. Let $Z$ be an irreducible
$n$-dimensional analytic cycle in $X$, and let $f$ be a moment map defined
near its support. On the relative interior of an
$n$-dimensional cell $\sigma$ of $\operatorname{Trop}_f(Z)$, the
tropicalization is generically finite in the graded-reduction sense, by
\Cref{prop:trop-degree-constant} applied to the components of $Z$.

\begin{definition}
\label{def:trop-cycle-multiplicity}
Let $e_\sigma$ be the primitive volume vector of the saturated tangent
lattice of $\sigma$. The \emph{tropical multiplicity}
$m_\sigma(Z,f)$ is the unique positive scalar for which the local
pushforward of the canonical calibration of $Z$ is
$m_\sigma(Z,f)e_\sigma$ on $\mathring\sigma$. Equivalently, for an
admissible projection $q$ it is the sum of the local analytic degrees above
$\sigma$, divided by the lattice index in
\eqref{eq:trop-weight-index}. By \Cref{thm:trop-degree-graded}, each analytic
degree already contains the length, graded residue degree, and possible
defect; none of these factors is to be inserted a second time.
\end{definition}

This definition is intentionally conceptual. In non-strict or non-algebraically
closed settings, naive counting of inverse images is wrong; graded reduction
and residue degrees are indispensable.

\begin{theorem}\label{thm:analytic-tropical-balancing}
The weighted $n$-dimensional part of the tropicalization of an
$n$-dimensional analytic cycle is balanced at every codimension-one cell not
meeting the image of the analytic boundary; it is zero if no
$n$-dimensional cell occurs.  In particular it is balanced everywhere when
the cycle is boundaryless.  The boundary hypothesis is not decorative: the
proof below runs the slope formula on a transversal fibre and uses that the
divisor of a character on it has degree zero, which fails in the presence of
an analytic boundary, and a harmonic tropicalization of a closed disc gives
an unbalanced image (\Cref{chap:delta-forms}).  Suppose that a moment map $\psi$ refines a moment map $\varphi$
through an integer-affine map $F$, so that
$\varphi_{\mathrm{trop}}=F\circ\psi_{\mathrm{trop}}$, and that $F$ is
generically finite on each relevant source facet.  Then the weight on a
target facet is
\begin{equation}\label{eq:analytic-tropical-pushforward}
 m_{\sigma'}\bigl(F_*\operatorname{Trop}_{\psi}(Z)\bigr)
 =\sum_{F(\sigma)=\sigma'}
   [N_{\sigma'}:F(N_\sigma)]\,m_\sigma.
\end{equation}
If an analytic morphism $g:Z\to Z'$ of generic degree $\delta$ induces this
projection and $Z'$ is given its reduced cycle structure, the left-hand side
of \eqref{eq:analytic-tropical-pushforward} is
$\delta m_{\sigma'}(\operatorname{Trop}_{\varphi}(Z'))$.
\end{theorem}
\begin{proof}
Let $\tau$ be a codimension-one cell and let
$\sigma_1,\ldots,\sigma_s$ be the facets containing it.  Balancing is an
identity in $N/N_\tau$, so it suffices to pair it with every integral
covector $m$ annihilating $N_\tau$.  Choose a point in the relative interior
of $\tau$ and a sufficiently small transversal polyhedral neighborhood.
After fixing $n-1$ generic monomial coordinates parallel to $\tau$, the
corresponding analytic fiber is one-dimensional.  Its tropicalization has
outgoing branches indexed by the $\sigma_i$; the degree of the branch is
exactly the initial-degeneration multiplicity $m_{\sigma_i}$.

The character $\chi^m$ restricts to an invertible analytic function on this
fiber.  Along the branch indexed by $\sigma_i$, its outgoing slope is
$\langle m,u_{\sigma_i/\tau}\rangle$.  The non-Archimedean slope formula on
the compactified analytic curve (\Cref{thm:ch08-slope-formula}, in the form
\eqref{eq:ch08-slope-formula}) says that the sum of the outgoing slopes,
weighted by the branch degrees, equals the degree of the divisor of
$\chi^m$.  The divisor has degree zero.  Hence
\[
 \sum_i m_{\sigma_i}
       \langle m,u_{\sigma_i/\tau}\rangle=0.
\]
Since the integral covectors annihilating $N_\tau$ span the dual of
$N/N_\tau$, this is the balancing equation.

Now let $F:N\to N'$ be the linear part of the integer-affine projection.
Refine so that $F$ has constant rank on every facet.  On a facet on which it
is generically finite, the induced map of initial degenerations is finite.
Localize at a generic point of the target initial component.  The source
coordinate algebra is then a finite module, and additivity of length gives
\begin{equation*}
 m_{\sigma'}\bigl(F_*\operatorname{Trop}_{\psi}(Z)\bigr)
 =\sum_{F(\sigma)=\sigma'}
   [N_{\sigma'}:F(N_\sigma)]\,m_\sigma.
\end{equation*}
The lattice index is the degree contributed by the monomial torus map.
The coefficient of the cycle-theoretic pushforward of the tropical cycle
already contains the degree of the finite analytic map; inserting a second
analytic-degree factor would count that degree twice.  Equivalently, for the
analytic morphism in the statement one has $g_*[Z]=\delta[Z']$ over the
generic locus.  Taking initial degenerations commutes with this finite
pushforward, and additivity of length gives the second formulation.
This is precisely the pushforward rule for weighted tropical cycles.
Conservation of rank in a tower of finite modules proves functoriality under
composition.

The two analytic ingredients used in this reduction are the slope formula
and conservation of degree for graded reductions; their invariant
construction, including non-strict fields and residue-degree factors, is
proved in \cite[\S\S12--15]{CLD25}.
\end{proof}

\begin{remark}[Two routes to balancing]
\label{rem:trop-balancing-two-routes}
The proof just given is the curve-slope route: one cuts down to a
one-dimensional analytic fibre transverse to a codimension-one cell and
applies the slope formula there. It is the shortest route and the one that
shows why balancing is a residue theorem in disguise. There is a second,
invariant route, which does not choose transversal coordinates: one proves
that the interior of an analytic space of level $n$, equipped with its
canonical calibration, is \emph{harmonious}, meaning that the calibrated
discordance vector at every codimension-one face vanishes; pairing that
vector with integral covectors annihilating $N_\tau$ gives back the scalar
identities above. This is \cite[Theorem~15.6.1]{CLD25}, quoted in this book
as \Cref{thm:appG-analytic-calibration}. The two routes give the same
weights, because both compute the calibration of
\Cref{def:trop-canonical-calibration}, whose uniqueness is
\Cref{cor:trop-calibration-glue}. The harmony route is the one that extends
to the boundary statements and to fibre integration in
\Cref{chap:cld-currents}; the curve route is the one that generalizes to the
curve-test criteria of \Cref{chap:delta-forms}.
\end{remark}

The projection formula is exactly what makes calibrations compatible under
the product refinements of \Cref{lem:product-moment-map}.

\begin{corollary}\label{cor:analytic-integration-independent}
Let $\alpha$ be a compactly supported $(n,n)$-form on $X$. Compute its
integral in any finite family of tropical charts using the canonical
calibrations and a partition of unity. The result is independent of all
charts, common refinements, subdivisions, and the partition of unity.
\end{corollary}
\begin{proof}
On a common product refinement, the projection formula of
\Cref{thm:analytic-tropical-balancing} identifies the two calibrated
integrals. Subdivision independence is
\Cref{prop:integral-subdivision-independent}. Additivity and the
identity $\sum\rho_i=1$ give independence of the partition. Compact support
reduces the comparison to finitely many charts.
\end{proof}

The last compatibility needed downstream concerns not a change of chart but
a change of space: an analytic morphism with finite fibres over a local
skeleton. This is the analytic counterpart of the integer-affine pushforward
of \Cref{prop:tropical-pushforward-balanced}, and it is the statement that
enters the projection formula for currents in \Cref{chap:cld-currents}.

\begin{theorem}
\label{thm:trop-finite-pushforward}
Let $p:Y\to X$ be a morphism of $K$-analytic spaces of level $n$, let
$\Sigma\subset X$ be a polyhedral subspace with canonical calibration
$\mu_\Sigma$, and let $\mathrm T=p^{-1}(\Sigma)$, a polyhedral subspace of
$Y$ with canonical calibration $\mu_{\mathrm T}$. Assume that $p$ has finite
fibres over $\Sigma$ and that $p_\Sigma:\mathrm T\to\Sigma$ is
topologically compact. Then
\begin{equation}
\label{eq:trop-finite-pushforward}
 p_*\mu_{\mathrm T}=\mathrm{dc}(p)\,\mu_\Sigma,
\end{equation}
where $\mathrm{dc}_x(p)=\sum_{y\in p^{-1}(x)}\mathrm{dc}_y(p)$ is the
calibrated degree of \Cref{sec:trop-graded}, a function which is locally
constant on $\Sigma$ off a polyhedral subset of dimension $<n$.
\end{theorem}

\begin{proof}
That $\mathrm T$ is polyhedral is
\Cref{thm:trop-sigma-polyhedral} \cref{itm:trop-poly:union}. Both sides of
\eqref{eq:trop-finite-pushforward} are calibrations, and calibrations form a
$G$-sheaf, so we may work $G$-locally: replace $X$ by a compact analytic
domain $V$ carrying a moment map $f:V\to\mathbf G_m^n$ with
$\Sigma\subset\Sigma_f$, and $Y$ by a compact analytic neighbourhood of
$p^{-1}(\Sigma)$; this is legitimate because canonical calibrations and
local degrees commute with restriction to analytic domains
(\Cref{cor:trop-calibration-restriction} and the remark after
\Cref{def:trop-local-degree}).

Both $f$ and $f\circ p$ are then skeletally proper, and
$\Sigma_{f\circ p}=p^{-1}(\Sigma_f)$,
$\Sigma_{f\circ p}^{(n)}=p^{-1}(\Sigma_f^{(n)})$. Choose adapted
decompositions: $\mathscr C$ of $\Sigma_f^{(n)}$ into convenient $f$-cells,
$\mathscr D$ of $\Sigma_{f\circ p}^{(n)}$ into convenient
$(f\circ p)$-cells, such that $p$ carries cells to cells.

Let $C\in\mathscr C$ be an $n$-cell and let $D_1,\ldots,D_r$ be the
$n$-cells of $\mathscr D$ mapping onto $C$; over $\mathring C$ the map
$p_\Sigma$ is a finite topological covering with sheets
$\mathring D_1,\ldots,\mathring D_r$. Fix $x\in\mathring C$ and let $y_i$ be
its unique preimage in $\mathring D_i$. By convenience of the cells,
$x\in\Int(X)$ and $y_i\in\Int(Y)$, so $f$ is finite and flat at $x$ and
$f\circ p$ is finite and flat at each $y_i$
(\Cref{lem:trop-quasifinite-flat}). For a morphism finite and flat at a
point, the calibrated degree equals the local degree
(\Cref{prop:trop-calibrated-degree} \cref{itm:trop-dc:flat}), so
$\deg_{y_i}(f\circ p)=\mathrm{dc}_{y_i}(f\circ p)$; multiplicativity of
calibrated degrees (\Cref{prop:trop-calibrated-degree}
\cref{itm:trop-dc:mult}) gives
\[
 \deg_{y_i}(f\circ p)=\mathrm{dc}_{y_i}(p)\,\mathrm{dc}_x(f)
 =\mathrm{dc}_{y_i}(p)\,\deg_x(f).
\]
Summing over $i$,
\[
 \sum_{i=1}^r\deg_{y_i}(f\circ p)
 =\Bigl(\sum_{i=1}^r\mathrm{dc}_{y_i}(p)\Bigr)\deg_x(f)
 =\mathrm{dc}_x(p)\,\deg_x(f).
\]
The left-hand side is, by \Cref{thm:trop-multiplicity-formula} applied to
$f\circ p$ and by \Cref{def:trop-canonical-calibration}, the coefficient
appearing in $p_*\mu_{\mathrm T}$ on $C$, while $\deg_x(f)$ is the
coefficient of $\mu_\Sigma$ on $C$; the affine sections are the same because
$(f\circ p)_{\mathrm{trop}}=f_{\mathrm{trop}}\circ p$ and $p$ maps
$\mathring D_i$ isomorphically onto $\mathring C$. This is
\eqref{eq:trop-finite-pushforward} on $\mathring C$. Local constancy of
$\mathrm{dc}(p)$ off a lower-dimensional polyhedral set follows, since the
two sides are cellwise constant on a common decomposition. The full
statement, including the case where $p$ is not flat, is
\cite[Proposition~14.4.2]{CLD25}.
\end{proof}

\begin{corollary}
\label{cor:trop-finite-field-extension}
Let $F/K$ be a finite extension of complete valued fields and let $X$ be an
$F$-analytic space of level $n$; write $X_0$ for $X$ regarded as a
$K$-analytic space. A subset of $X$ is a polyhedral subspace for $F$ if and
only if it is one for $K$, and the two canonical calibrations satisfy
\[
 \mu_0=[F:K]\,\mu .
\]
\end{corollary}

\begin{proof}
The canonical skeleton $S(\mathbf G_{m,F}^n)$ is the preimage of
$S(\mathbf G_{m,K}^n)$ under the projection
$\mathbf G_{m,F}^n\to\mathbf G_{m,K}^n$, and surjects onto it: indeed the
graded transcendence degree of the $\widetilde{T_i(x)}$ over
$\widetilde F$ equals that over $\widetilde K$, since
$\widetilde F/\widetilde K$ is algebraic. Hence the two notions of local
skeleton agree, by \Cref{thm:trop-sigma-characterization}. The projection
$\mathbf G_{m,F}^n\to\mathbf G_{m,K}^n$ is finite and flat of degree
$[F:K]$, so applying \Cref{def:trop-canonical-calibration} with the
$K$-moment map obtained by composition multiplies every local degree by
$[F:K]$ while leaving the affine sections unchanged.
\end{proof}

\section{Graded reduction, cycles, and the source of the multiplicities}
\label{sec:trop-graded}

Everything so far has used the local degree of a finite flat morphism as a
black box. The purpose of this section is to open it. We show that the
local degree factors as
\[
 \deg_x(g)=\underbrace{\ell_x(g)}_{\text{length}}\cdot
 \underbrace{[\widetilde{\mathcal H(x)}:\widetilde{\mathcal H(g(x))}]}
 _{\text{graded residue degree}}\cdot
 \underbrace{\mathrm{def}(\mathcal H(x)/\mathcal H(g(x)))}_{\text{defect}},
\]
that the defect is $1$ in residue characteristic zero, and that the
resulting multiplicities are additive over analytic cycles. This makes
graded reduction, and not any auxiliary combinatorial recipe, the source of
the tropical multiplicities --- which is how the theory is used in
\Cref{chap:cld-currents,chap:cld-metrized-bundles}.

\subsection{Calibrated degrees}

The bookkeeping is cleanest in terms of a rational refinement of the local
degree, which divides out the generic multiplicity of the source.

\begin{definition}[Calibrated degree]
\label{def:trop-calibrated-degree}
\index{degree!calibrated}
For a local artinian ring $A$ and a finite-length $A$-module $N$, put
$\delta_A(N)=\ell_A(N)/\ell_A(A)\in\mathbf Q_{>0}$. Let $g:Y\to X$ be a
morphism of $K$-analytic spaces, $y\in Y$, $x=g(y)$, and suppose $g$ is
quasi-finite at $y$ and that the local ring $\mathcal O_{X_G,x}$ is
artinian. The \emph{calibrated degree} of $g$ at $y$ is
\[
 \mathrm{dc}_y(g)=\delta_{\mathcal O_{X_G,x}}(\mathcal O_{Y_G,y}),
 \qquad
 \mathrm{dc}_x(g)=\sum_{y\in g^{-1}(x)}\mathrm{dc}_y(g).
\]
\end{definition}

\begin{proposition}
\label{prop:trop-calibrated-degree}
With the notation of \Cref{def:trop-calibrated-degree}:
\begin{enumerate}
\item\label{itm:trop-dc:flat} if $g$ is quasi-finite and flat at $y$ and
  $\mathcal O_{X_G,x}$ is artinian, then
  $\mathrm{dc}_y(g)=\deg_y(g)$;
\item\label{itm:trop-dc:mult} calibrated degrees are multiplicative: for
  $h:Z\to Y$ and $z\in h^{-1}(y)$ with $\dim_zh=0$,
  $\mathrm{dc}_z(g\circ h)=\mathrm{dc}_y(g)\,\mathrm{dc}_z(h)$;
\item\label{itm:trop-dc:red} for the closed immersion
  $\iota:X_{\mathrm{red}}\to X$ one has
  $\mathrm{dc}_x(\iota)=1/\ell(\mathcal O_{X_G,x})$, the reciprocal of the
  generic multiplicity of $X$ along the unique irreducible component through
  $x$;
\item\label{itm:trop-dc:factor} if $A\to B$ is a finite local homomorphism
  of artinian local rings with residue fields $\kappa$ and $\lambda$, then
  \[
   \delta_A(B)=\frac{\ell_B(B)}{\ell_A(A)}\,[\lambda:\kappa].
  \]
\end{enumerate}
\end{proposition}

\begin{proof}
\cref{itm:trop-dc:flat} A flat module of finite type over a local ring is
free, of rank $r$ say; then $\ell_A(B)=r\ell_A(A)$ and
$\delta_A(B)=r=\deg_y(g)$.

\cref{itm:trop-dc:factor} Choose a composition series of $B$ as a
$B$-module; it has $\ell_B(B)$ successive quotients, each isomorphic to
$\lambda$. As an $A$-module, $\lambda$ has length $\dim_\kappa\lambda$, so
$\ell_A(B)=\ell_B(B)[\lambda:\kappa]$ and the assertion follows on dividing
by $\ell_A(A)$.

\cref{itm:trop-dc:mult} Apply \cref{itm:trop-dc:factor} twice, to
$A\to B\to C$ with residue fields $\kappa,\lambda,\mu$:
\[
 \delta_A(C)=\frac{\ell_C(C)}{\ell_A(A)}[\mu:\kappa]
 =\frac{\ell_C(C)}{\ell_B(B)}[\mu:\lambda]\cdot
  \frac{\ell_B(B)}{\ell_A(A)}[\lambda:\kappa]
 =\delta_B(C)\,\delta_A(B).
\]

\cref{itm:trop-dc:red} Here $B=A_{\mathrm{red}}$, so $\ell_B(B)=1$ and
$\lambda=\kappa$, giving $\delta_A(B)=1/\ell_A(A)$; the identification of
$\ell(\mathcal O_{X_G,x})$ with the generic multiplicity along the component
through $x$ is the definition of the latter.
\end{proof}

Assertion \cref{itm:trop-dc:factor} already displays the two ingredients
promised in the chapter's title claim: a \emph{length}, which is the
scheme-theoretic multiplicity, and a \emph{residue degree}. On a local
skeleton the residue degree can be computed in the graded world, and that is
the content of the next theorem.

\subsection{Abhyankar points, defect, and the graded residue degree}

\begin{definition}
\label{def:trop-abhyankar-defect}
\index{Abhyankar extension}\index{defect}
A complete valued extension $L/K$ is an \emph{Abhyankar extension} if there
are a $K$-analytic space $X$ and a point $x$ with $d_K(x)=\dim_xX$ and
$\mathcal H(x)\simeq L$ over $K$; equivalently, if
$\operatorname{trdeg}_{\widetilde K}\widetilde L$ attains the Abhyankar
bound. A family $t=(t_1,\ldots,t_d)$ in $L$ is an \emph{Abhyankar basis} if
the $\widetilde{t_i}$ form a graded transcendence basis of $\widetilde L$
over $\widetilde K$; equivalently if
\[
 \Bigl|\sum_Ia_It^I\Bigr|=\max_I|a_I|\,|t|^I
 \qquad(a_I\in K,\ \text{almost all zero}),
\]
so that $\widehat{K(t)}\simeq\mathcal H(\eta_r)$ with
$r=(|t_1|,\ldots,|t_d|)$. The \emph{defect} of $L$ is
\[
 \mathrm{def}_K(L)
 =\frac{[L:\widehat{K(t)}]}
 {\bigl[\widetilde L:\widetilde{\widehat{K(t)}}\bigr]},
\]
which is independent of the Abhyankar basis $t$.
\end{definition}

That $L$ is finite over $\widehat{K(t)}$, and that the quotient above does
not depend on $t$, is \cite[Lemma~13.8.6 and Theorem~13.8.10]{CLD25}. For
$d=0$ the definition specializes to the classical defect of a finite
extension of complete valued fields. We use two standard facts, quoted for
reference: the defect is a positive integer, and it is a power of the
residue characteristic; in particular
\begin{equation}
\label{eq:trop-defect-trivial}
 \mathrm{def}_K(L)=1
 \qquad\text{whenever }\operatorname{char}\widetilde K=0,
\end{equation}
by Ostrowski's lemma for complete rank-one valued fields, see
\cite[Ch.~VI]{ZariskiSamuelII} and \cite[\S13.8]{CLD25}.

\begin{theorem}
\label{thm:trop-degree-graded}
Let $X$ be of level $n$, let $f:X\to\mathbf G_m^n$ be a moment map, and let
$x\in\Sigma_f\setminus\partial X$, so that $f$ is finite and flat at $x$
(\Cref{lem:trop-quasifinite-flat}). Put $t=f(x)\in S(\mathbf G_m^n)$. Then
\begin{equation}
\label{eq:trop-degree-graded}
 \deg_x(f)
 =\ell\bigl(\mathcal O_{f^{-1}(t),x}\bigr)\cdot
  \bigl[\widetilde{\mathcal H(x)}:\widetilde{\mathcal H(t)}\bigr]\cdot
  \mathrm{def}_{\mathcal H(t)}\bigl(\mathcal H(x)\bigr).
\end{equation}
In particular, if $\operatorname{char}\widetilde K=0$ the local degree is the
product of a length and a graded residue degree; and if in addition
$f^{-1}(t)$ is reduced at $x$, it is exactly the graded residue degree
$[\widetilde{\mathcal H(x)}:\widetilde{\mathcal H(t)}]$.
\end{theorem}

\begin{proof}
Since $x$ is isolated in the fibre $f^{-1}(t)$, the local ring
$A=\mathcal O_{f^{-1}(t),x}$ is an artinian local
$\mathcal H(t)$-algebra of finite dimension, and by definition
$\deg_x(f)=\dim_{\mathcal H(t)}A$. Its residue field is a finite extension
of $\mathcal H(t)$, hence complete, hence equal to $\mathcal H(x)$: indeed
$\mathcal H(x)$ is by construction the completion of the residue field of
$A$, and a finite extension of a complete field is complete. Applying
\Cref{prop:trop-calibrated-degree} \cref{itm:trop-dc:factor} with
$A_0=\mathcal H(t)$ --- a field, so $\ell(A_0)=1$ --- gives
\[
 \deg_x(f)=\dim_{\mathcal H(t)}A
 =\ell(A)\cdot[\mathcal H(x):\mathcal H(t)].
\]
It remains to factor the field degree. By
\Cref{cor:trop-sigma-abhyankar}, $d_K(x)=n$, and $d_K(t)=n$ by
\Cref{prop:trop-torus-skeleton}; so $\mathcal H(x)$ and $\mathcal H(t)$ are
Abhyankar extensions of $K$ of the same graded transcendence degree, and
$\mathcal H(x)/\mathcal H(t)$ is a finite extension of Abhyankar type with
graded transcendence degree zero. \Cref{def:trop-abhyankar-defect} applied
with $d=0$ over the base $\mathcal H(t)$ is exactly the classical
multiplicativity
\[
 [\mathcal H(x):\mathcal H(t)]
 =\bigl[\widetilde{\mathcal H(x)}:\widetilde{\mathcal H(t)}\bigr]\cdot
  \mathrm{def}_{\mathcal H(t)}(\mathcal H(x)),
\]
the graded residue degree being the product of the classical residue degree
and the ramification index, by \Cref{lem:graded-residue-field}. Substituting
gives \eqref{eq:trop-degree-graded}, and
\eqref{eq:trop-defect-trivial} gives the residue-characteristic-zero
statement. If $f^{-1}(t)$ is reduced at $x$ then $A$ is a field and
$\ell(A)=1$.
\end{proof}

\begin{remark}
\label{rem:trop-degree-graded-reading}
\Cref{thm:trop-degree-graded} is the promised identification. Reading
\eqref{eq:trop-degree-graded} from right to left: the defect is an
arithmetic pathology absent in residue characteristic zero; the graded
residue degree is the size of the extension of graded reductions, that is,
of initial-form data; and the length is the scheme-theoretic multiplicity of
the fibre. \Cref{prop:trop-calibrated-degree} \cref{itm:trop-dc:red} adds the
fourth ingredient present for nonreduced ambient spaces: the generic
multiplicity of $X$ itself, which is divided out by the calibrated degree.
Thus the tropical weight of a facet is
\[
 m_D=\Bigl[\text{lattice index}\Bigr]^{-1}
 \sum_{C\mapsto D}\ \bigl(\text{length}\bigr)\cdot
 \bigl(\text{graded residue degree}\bigr)\cdot\bigl(\text{defect}\bigr),
\]
each factor coming from a different source, and none of them a cardinality.
\Cref{ex:trop-multiplicity-two} exhibits a weight produced entirely by the
graded residue degree, and \Cref{ex:trop-length-factor} one produced
entirely by the length; \Cref{ex:trop-two-charts} exhibits the lattice
index; and the defect is invisible in residue characteristic zero by
\eqref{eq:trop-defect-trivial}.
\end{remark}

\subsection{Cycles}

Everything above was stated for a space. Analytic cycles enter by
additivity.

\begin{definition}
\label{def:trop-cycle-calibration}
Let $X$ be a $K$-analytic space and let $Z=\sum_j n_j[Z_j]$ be an effective
analytic cycle of pure dimension $n$ on $X$, the $Z_j$ being the reduced
irreducible closed analytic subspaces of dimension $n$ and $n_j>0$. For a
moment map $f$ defined near the support of $Z$, assume that
$f_{\mathrm{trop}}$ is proper on the top-dimensional skeletal support
$\bigcup_j\Sigma_{f|_{Z_j}}^{(n)}$; this holds, for example, when the support
of $Z$ is compact. Put
\[
 \Sigma_f(Z)=\bigcup_j\Sigma_{f|_{Z_j}},
 \qquad
 \mu_{\Sigma_f(Z)}=\sum_jn_j\,\mu_{\Sigma_{f|_{Z_j}}},
\]
the sum being taken after extending each calibration by zero. The
\emph{tropicalization of the cycle} is the pushforward
$\Trop_f(Z)=f_{\mathrm{trop},*}\mu_{\Sigma_f(Z)}$.
\end{definition}

\begin{proposition}
\label{prop:trop-cycle-additive}
Whenever the pushforwards involved are proper on their skeletal supports,
the assignment $Z\mapsto\Trop_f(Z)$ is additive in $Z$, compatible with
restriction to analytic domains and with extension of the base field, and
satisfies
\[
 \Trop_{q\circ f}(Z)=q_{\mathrm{trop},*}\Trop_f(Z)
\]
for every homomorphism of tori $q$ that is generically finite on each facet
of $\Trop_f(Z)$. If $g:Z\to Z'$ is a finite morphism of pure
$n$-dimensional reduced analytic spaces with $g_*[Z]=\delta[Z']$ generically,
and $f'$ is a moment map on $Z'$, then
$\Trop_{f'\circ g}(Z)=\delta\,\Trop_{f'}(Z')$ over the locus where $g$ has
constant calibrated degree $\delta$.
\end{proposition}

\begin{proof}
Additivity is \Cref{def:trop-cycle-calibration}, and the compatibilities are
\Cref{cor:trop-calibration-restriction} and
\Cref{prop:trop-calibration-base-change} applied componentwise. The formula
for $q\circ f$ is \Cref{thm:trop-multiplicity-formula} together with the
transitivity of pushforward of calibrations
(\Cref{prop:appG-calibration-tests}), the lattice index
\eqref{eq:appG-lattice-index} being exactly the discrepancy between
$\sigma_q$ computed upstairs and downstairs. The last assertion is
\Cref{thm:trop-finite-pushforward} applied to $g$, together with
\Cref{prop:trop-calibrated-degree} \cref{itm:trop-dc:mult}, which converts
the calibrated degree of $g$ into the cycle-theoretic coefficient $\delta$.
\end{proof}

\begin{warning}{Do not count the analytic degree twice}
In \Cref{prop:trop-cycle-additive} the coefficient $\delta$ appears exactly
once. The calibration $\mu_{\Sigma_f(Z)}$ upstairs already contains the local
degrees of $Z$ over the torus; pushing forward along $g$ multiplies them by
the calibrated degree of $g$ and by nothing else. Multiplying again by
$\deg g$, as one is tempted to do when thinking of
\eqref{eq:analytic-tropical-pushforward} as a set-theoretic count, produces
$\delta^2$ and contradicts the one-dimensional test of
\Cref{ex:trop-two-charts}. The same warning appears at
\eqref{eq:analytic-tropical-pushforward} itself and in
\Cref{prop:appG-calibration-tests}.
\end{warning}

\section{Compact supports, weak analytic structures, and partitions of unity}
\label{sec:trop-supports}

Three technical facts are exported from this chapter to the rest of Part~II.
The first says that certain subsets which are not analytic spaces --- most
importantly, boundaries --- still have polyhedral tropicalizations of the
expected dimension; without it one cannot even assert that a boundary is
negligible for integration. The second is the existence of partitions of
unity, which is what makes chartwise definitions global. The third is that
compactly supported forms need only finitely many charts.

\subsection{Weak analytic structures}

\begin{definition}[Weakly analytic subset]
\label{def:trop-weakly-analytic}
\index{weakly analytic subset}
Let $X$ be a $K$-analytic space, $A\subset X$ a closed subset, and
$d\in\mathbf N$. We say that $A$ is \emph{weakly analytic of dimension at
most $d$} if for every compact analytic domain $V\subset X$ there are a
finite family of morphisms of $K$-analytic spaces $g_i:V_i\to W_i$, defined
on compact analytic domains $V_i\subset V$, and points $w_i\in W_i$, with
\[
 A\cap V=\bigcup_ig_i^{-1}(w_i),
 \qquad
 \dim_{g_i}(x)\le d\ \text{ for all }x\in A\cap V_i,
\]
where $\dim_{g_i}(x)$ denotes the dimension at $x$ of the fibre of $g_i$
through $x$. This is a $G$-local property, and a locally finite union of
weakly analytic subsets is weakly analytic.
\end{definition}

The point of the definition is that a fibre of a morphism carries enough
structure to be tropicalized, even though it may fail to be an analytic
space over $K$: base change along $\mathcal H(w_i)$ turns it into an
$\mathcal H(w_i)$-analytic space, and a moment map on the ambient space
restricts to one on it. This suffices for the polyhedrality statement.

\begin{proposition}
\label{prop:trop-weakly-analytic-image}
Let $X$ be a compact $K$-analytic space, let $A\subset X$ be weakly analytic
of dimension at most $d$, and let $f:X\to T$ be a moment map. Then
$f_{\mathrm{trop}}(A)$ is a polyhedral subset of $T_{\mathrm{trop}}$ of
dimension at most $d$.
\end{proposition}

\begin{proof}
Since $X$ is compact and the condition is $G$-local and stable under finite
unions, we may assume $A=g^{-1}(w)$ for a morphism $g:X\to W$ of
$K$-analytic spaces and a point $w\in W$, with all fibre dimensions along
$A$ at most $d$. The fibre product $X\times_W\mathcal M(\mathcal H(w))$
realizes $A$ as a compact $\mathcal H(w)$-analytic space of dimension at most
$d$, and $f|_A$ as a moment map on it with values in
$T_{\mathcal H(w)}$. Tropicalization is unchanged by this base change
(\Cref{lem:trop-moment-basic} \cref{itm:trop-moment:base}). Now apply
\Cref{thm:analytic-tropical-polyhedrality} over the field
$\mathcal H(w)$: the image is a polyhedron of dimension at most
$\dim A\le d$.
\end{proof}

\begin{proposition}
\label{prop:trop-boundary-thin}
Let $X$ be a compact $K$-analytic space of dimension $n$ and $f:X\to T$ a
moment map. Then $\partial X$ is weakly analytic of dimension at most
$n-1$, and $f_{\mathrm{trop}}(\partial X)$ is contained in a polyhedral
subset of $T_{\mathrm{trop}}$ of dimension at most $n-1$. If $X$ is
affinoid, $f_{\mathrm{trop}}(\partial X)$ is itself such a subset.
\end{proposition}

\begin{proof}
For an affinoid space, the boundary theorem of Ducros gives directly that
$\partial X$ is weakly analytic of dimension at most $\dim X-1$; this is
\cite[Example~11.10.2]{CLD25}. The nearby
\cite[Proposition~11.10.4]{CLD25} gives a complementary graded-reduction
description of weakly analytic strata, but the Shilov boundary is not to be
confused with the whole analytic boundary. In general, take an affinoid
$G$-cover: the analytic boundary is $G$-locally contained in the
corresponding affinoid boundaries, and weak analyticity is $G$-local and
stable under locally finite unions. Thus $\partial X$ has dimension at most
$n-1$. The assertions about $f_{\mathrm{trop}}(\partial X)$ now follow from
\Cref{prop:trop-weakly-analytic-image}; in the affinoid case that proposition
applies directly to $\partial X$, while in general it gives the asserted
containing polyhedral set.
\end{proof}

\begin{corollary}
\label{cor:trop-boundary-negligible}
Let $X$ be of level $n$ and $f:X\to T$ a skeletally proper moment map. Then
$f_{\mathrm{trop}}(\partial X\cap\Sigma_f^{(n)})$ is a closed polyhedral
subset of $\Pi_f$ of dimension at most $n-1$. Consequently, the integral of
a compactly supported $(n,n)$-form against the canonical calibration is
unchanged if $\partial X$ is removed from the domain of integration, and the
degree functions of \Cref{prop:trop-degree-constant} may be computed on
relative interiors of cells only.
\end{corollary}

\begin{proof}
The first assertion is \Cref{prop:trop-boundary-thin} applied $G$-locally
combined with the local finiteness of a cell decomposition; the second is
\Cref{prop:integral-subdivision-independent}, a subset of dimension at most
$n-1$ carrying no mass for a calibrated $n$-dimensional integral.
\end{proof}

\subsection{Partitions of unity}

\begin{theorem}\label{thm:cld-partitions-unity}
Let $X$ be a good, topologically separated, paracompact Berkovich analytic
space. Every open cover of $X$ admits a subordinate smooth partition of
unity in the CLD sheaf. Compactly supported forms can therefore be integrated
using finitely many tropical charts.
\end{theorem}
\begin{proof}
A good topologically separated analytic space is locally holomorphically
separated, because each point has an affinoid neighbourhood. The partition
theorem is therefore exactly \cite[Proposition~15.2.8]{CLD25}, obtained by
applying the smooth partition theorem for locally differentially separated
tropical spaces. Local finiteness implies that a compact support meets only
finitely many members of the partition; after refining those members by
tropical charts, integration is a finite chartwise sum.
\end{proof}

For spaces not satisfying goodness or paracompactness, we do not silently use
partitions of unity; compact-support statements are made chartwise or after
imposing the needed topology.

\section{Worked examples}
\label{sec:trop-worked}

We now compute five examples in full. The first two are the standard tests
of the multiplicity conventions; the third exhibits refinement independence
concretely; the fourth is the running degeneration of Parts~I and~II; the
fifth produces a weight different from one and traces it back to graded
reduction. Throughout, $K$ is complete and nontrivially valued, and we say
explicitly when it is assumed algebraically closed.

\subsection{A line in a two-dimensional torus}

\begin{example}[Tropicalization of a line in $\mathbf G_m^2$]
\label{ex:trop-line-in-torus}
Assume $K$ algebraically closed. Let
\[
 Z=\{1+T_1+T_2=0\}\subset\mathbf G_{m,K}^2,
\]
a smooth affine curve, and let $X=Z^{\An}$, of level $1$. Take the moment
map $f=(T_1,T_2)|_Z$ into $T=\mathbf G_m^2$, so $M=\mathbf Z^2$ and
$T_{\mathrm{trop}}=\mathbf R^2$. Put $s=T_1|_Z$; then $T_2|_Z=-1-s$ and
\[
 Z\ \simeq\ \mathbf A^1_K\setminus\{0,-1\},
 \qquad t\longmapsto(t,-1-t),
\]
so $X$ is the analytic line with two points removed. We compute all of
$\Trop_f$, $\Sigma_f$, and the calibration.

\emph{The tropical image.} Write $u=(u_1,u_2)=(v(s),v(1+s))$ where
$v=-\log|\cdot|$. If $v(s)>0$ then $|s|<1$, so $|1+s|=1$ and $u=(u_1,0)$ with
$u_1>0$. If $v(s)<0$ then $|s|>1$, so $|1+s|=|s|$ and $u=(u_1,u_1)$ with
$u_1<0$. If $v(s)=0$ then $u=(0,u_2)$ with $u_2\ge0$, and every such $u_2$
occurs, for instance at $s$ close to $-1$. Hence
\[
 \Trop_f(X)=\mathbf R_{\ge0}(1,0)\ \cup\ \mathbf R_{\ge0}(0,1)
 \ \cup\ \mathbf R_{\ge0}(-1,-1),
\]
the tropical line of \Cref{ex:appG-tropical-line} and
\Cref{fig:tropical-line-balancing}, in the valuation convention.

\emph{The local skeleton.} By \Cref{thm:trop-sigma-characterization},
$x\in\Sigma_f$ if and only if $\widetilde{T_1(x)}$ or $\widetilde{T_2(x)}$ is
transcendental over $\widetilde K$. In the coordinate $s$ this says that $x$
is a point of type II or III of $\mathbf A^{1,\An}_K$
(\Cref{prop:berk-type-invariants}) lying on one of the three segments
joining the Gauss point $\zeta_{0,1}$ to the three punctures
$s=0$, $s=-1$, $s=\infty$. Thus $\Sigma_f$ is the open tripod
\[
 \Sigma_f=\{\zeta_{0,\rho}:0<\rho<1\}\cup\{\zeta_{-1,\rho}:0<\rho<1\}
 \cup\{\zeta_{0,\rho}:\rho>1\}\cup\{\zeta_{0,1}\},
\]
and $f_{\mathrm{trop}}$ restricts to a bijection $\Sigma_f\to\Trop_f(X)$
carrying $\zeta_{0,\rho}$ with $\rho<1$ to $(-\log\rho,0)$, carrying
$\zeta_{-1,\rho}$ to $(0,-\log\rho)$, and carrying $\zeta_{0,\rho}$ with
$\rho>1$ to $(-\log\rho,-\log\rho)$. Here tropicalization is injective on
the local skeleton --- a happy accident of this example, not a general
phenomenon.

\emph{The calibration.} We treat the three rays separately.

On $\sigma_1=\mathbf R_{\ge0}(1,0)$ take $q=\mathrm{pr}_1$, so
$q_{\mathrm{trop}}(u)=u_1$ is injective on $\sigma_1$. Then
$q\circ f=s:X\to\mathbf G_m^{\An}$ is the open immersion onto
$\mathbf G_m^{\An}\setminus\{-1\}$; a point of $X$ is the only point of its
fibre and $\mathcal H(x)=\mathcal H(s(x))$, so $\deg_x(q\circ f)=1$. The
affine section with image $\langle\sigma_1\rangle=\mathbf R\times\{0\}$ is
$\sigma_q(a)=(a,0)$, whence $\sigma_{q,*}(e)=|(1,0)|$ and
$\mu_{\sigma_1}=1\cdot|(1,0)|$: weight $1$.

On $\sigma_2=\mathbf R_{\ge0}(0,1)$ take $q=\mathrm{pr}_2$; then
$q\circ f=-1-s$ is again an open immersion, $\sigma_q(a)=(0,a)$, and
$\mu_{\sigma_2}=1\cdot|(0,1)|$: weight $1$.

On $\sigma_3=\mathbf R_{\ge0}(-1,-1)$ take $q=\mathrm{pr}_1$, injective on
$\sigma_3$; the section with image $\langle\sigma_3\rangle$ is
$\sigma_q(a)=(a,a)$, so $\sigma_{q,*}(e)=|(1,1)|$, which is the primitive
volume vector of $N_{\sigma_3}=\mathbf Z(1,1)$. Again $\deg=1$, so the weight
is $1$. Choosing instead $q=\mathrm{pr}_2$ gives the same section and the
same degree, as \Cref{thm:trop-calibration-canonical} predicts.

\emph{Balancing.} The primitive outgoing directions at the origin are
$(1,0)$, $(0,1)$, $(-1,-1)$ with weights $1,1,1$, and
$(1,0)+(0,1)+(-1,-1)=0$; this is the balancing condition
\eqref{eq:appG-balancing}, verified by hand, and predicted by
\Cref{thm:analytic-tropical-balancing}.

\emph{A projection test.} Push forward along $\mathrm{pr}_1$. The ray
$\sigma_2$ is contracted and drops out. The rays $\sigma_1$ and $\sigma_3$
map onto $\mathbf R_{\ge0}$ and $\mathbf R_{\le0}$ with lattice indices
$[\mathbf Z:\mathrm{pr}_1(\mathbf Z(1,0))]=1$ and
$[\mathbf Z:\mathrm{pr}_1(\mathbf Z(1,1))]=1$, so
$\mathrm{pr}_{1,*}\Trop_f(X)=\mathbf R$ with weight $1$. This is
$\Trop(\mathbf G_m)$ with its canonical calibration
(\Cref{ex:trop-torus-explicit}), consistent with the fact that
$s:X\to\mathbf G_m^{\An}$ is generically of degree $1$. Multiplying by a
further factor $\deg(s)=1$ would be harmless here; the next example shows
that in general it is not.
\end{example}

\subsection{Two charts with the same local skeleton}

\begin{example}[Refinement independence, exhibited]
\label{ex:trop-two-charts}
Let $X=\mathbf G_{m,K}^{\An}$, of level $1$, with coordinate $T$.

\emph{(a) A monomial pair.} Consider the moment maps
\[
 f=T:X\to\mathbf G_m,
 \qquad
 g=(T,T^2):X\to\mathbf G_m^2 .
\]
By \Cref{lem:trop-monomial-degree} \cref{itm:trop-mono:preimage} and
\Cref{thm:trop-sigma-characterization},
\[
 \Sigma_f=\Sigma_g=S(\mathbf G_m)=\{\eta_u:u\in\mathbf R\}
\]
--- the two charts see exactly the same subset of $X$, although their targets
have different dimensions and neither refines the other in the naive sense.
The tropical polyhedra are different: $\Pi_f=\mathbf R$ while
$\Pi_g=\{(a,2a):a\in\mathbf R\}\subset\mathbf R^2$.

Compute the calibration of $\Sigma:=S(\mathbf G_m)$ in the chart $g$, using
each of the two admissible projections in turn. With $q=\mathrm{pr}_1$ one
has $q\circ g=T$, which is the identity, of degree $1$; the affine section
of $\mathrm{pr}_1$ with image $\langle\Pi_g\rangle$ is
$\sigma_q(a)=(a,2a)$, so
\[
 \mu=1\cdot|(1,2)| .
\]
With $q=\mathrm{pr}_2$ one has $q\circ g=T^2$, of degree $2$ by
\Cref{lem:trop-monomial-degree} \cref{itm:trop-mono:degree}, and the section
is $\sigma_{q}(a)=(a/2,a)$, so
\[
 \mu=2\cdot\bigl|(\tfrac12,1)\bigr|=1\cdot|(1,2)| .
\]
The two agree, as they must by
\Cref{thm:trop-calibration-canonical}. Note that the agreement is a
cancellation between a degree and a lattice normalization: the analytic
degree $2$ is exactly compensated by the halving of the section.

Finally compare with the chart $f$. Here
$\mu=1\cdot|e_1|$ on $\Pi_f=\mathbf R$, and
$g_{\mathrm{trop}}=(1,2)\circ f_{\mathrm{trop}}$ carries $|e_1|$ to
$|(1,2)|$: the two calibrations of $\Sigma$ coincide, which is
\Cref{thm:trop-calibration-chart-independent} in its smallest nontrivial
instance. The weight of $\Sigma$ is $1$ in both charts, because $|(1,2)|$ is
the primitive volume vector of the tangent lattice
$\mathbf Z(1,2)$ of $\Pi_g$.

\emph{(b) A non-monomial pair.} Let $V=\{x:0<|T(x)|<1\}$ and consider
\[
 f=T|_V,\qquad h=T(1-T)|_V .
\]
On $V$ one has $|1-T|=1$, so $h_{\mathrm{trop}}=f_{\mathrm{trop}}$
\emph{exactly}, not merely up to refinement.
\Cref{cor:trop-sigma-translation} gives $\Sigma_f\cap V=\Sigma_h\cap V$, and
\Cref{thm:trop-degree-translation} with zero constant gives equality of the
local degrees, hence of the calibrations. One can also see this by hand:
if $t,t'\in V$ satisfy $t-t^2=t'-t'^2$, then $(t-t')(1-t-t')=0$, and
$|t+t'|<1$ forces $t=t'$; so $h$ is injective on $V$ and induces an
isomorphism onto its image, of degree $1$ --- the same degree as $f$. The
two charts therefore carry the same canonical calibration on the common
local skeleton $\{\eta_u:u>0\}$, even though $h$ is not a monomial in $T$
and neither chart refines the other.

\emph{(c) The product test.} The product moment map $(f,h)$ on $V$ has
tropical image the diagonal $\{(a,a):a>0\}$, and both projections have
lattice index $1$; the calibration computed in the product chart is
$1\cdot|(1,1)|$, and pushes forward to $1\cdot|e_1|$ in either factor. This
is the mechanism of the proof of
\Cref{thm:trop-calibration-chart-independent} in miniature: one goes up to
the product, where the comparison is a tautology, and comes back down twice.
\end{example}

\subsection{An elliptic curve with multiplicative reduction}

\begin{example}[The Tate curve]
\label{ex:trop-tate-curve}
Let $q\in K^\times$ with $0<|q|<1$ and let
\[
 E=\mathbf G_{m,K}^{\An}\big/q^{\mathbf Z}
\]
be the associated Tate curve, an elliptic curve over $K$ with split
multiplicative reduction, studied as a degeneration in
\Cref{ex:tate-degeneration} and \Cref{sec:worked-tate}, and metrically in
\Cref{ex:ch08-tate-curve}. Write $\lambda=v(q)=-\log|q|>0$ and let
$p:\mathbf G_{m}^{\An}\to E$ be the quotient map, a local isomorphism
because $q^{\mathbf Z}$ acts freely and discontinuously.

\emph{No global moment map.} $E$ is proper, so
$\mathcal O(E)^\times=K^\times$ and every global moment map on $E$ has
constant tropicalization, hence $\Sigma=\varnothing$ for it. Local skeleta
are genuinely local objects; this example is the reason
\Cref{cor:trop-calibration-glue} is stated as a gluing theorem and not as a
single chart.

\emph{Two charts.} Set
\[
 U_1=p\bigl(\{0<v(T)<\lambda\}\bigr),
 \qquad
 U_2=p\bigl(\{\tfrac\lambda2<v(T)<\tfrac{3\lambda}2\}\bigr),
\]
two open subsets covering $E$: the first omits the locus $v(T)\equiv0$ and
the second the locus $v(T)\equiv\lambda/2$ modulo $\lambda$. On each of them
$p$ admits a section $s_j$, and we let $f_j=T\circ s_j:U_j\to\mathbf G_m$ be
the resulting moment maps.

The overlap $U_1\cap U_2$ has two connected components, corresponding to
$v(T)\bmod\lambda\in(0,\lambda/2)$ and to
$v(T)\bmod\lambda\in(\lambda/2,\lambda)$. On the second, $s_1$ and $s_2$
have the same image $\{\lambda/2<v(T)<\lambda\}$ and coincide, so
$f_2=f_1$ there. On the first, $s_1$ lands in
$\{0<v(T)<\lambda/2\}$ while $s_2$ lands in
$\{\lambda<v(T)<3\lambda/2\}$, so $s_2=q\,s_1$ and
\[
 f_2=q\,f_1,
 \qquad
 f_{2,\mathrm{trop}}=f_{1,\mathrm{trop}}+\lambda .
\]
Thus the two tropicalizations differ by a constant on each connected
component of the overlap. By \Cref{cor:trop-sigma-translation} the two local
skeleta agree there, and by \Cref{thm:trop-degree-translation} so do the
local degrees, hence the calibrations. This is a second concrete instance of
refinement independence, and the one that makes the global object exist.

\emph{The skeleton and its calibration.} Since
$\Sigma_{f_j}=p\bigl(S(\mathbf G_m)\cap\{\cdots\}\bigr)$ and $S(\mathbf G_m)$
is the line $\{\eta_u:u\in\mathbf R\}$, the local skeleta glue to
\[
 \Sk(E)=p\bigl(S(\mathbf G_m)\bigr)\ \simeq\ \mathbf R/\lambda\mathbf Z,
\]
a circle. Each $f_j$ is an isomorphism onto its image, so every local degree
is $1$ and $\sigma_{q}$ is the identity for $q=\mathrm{id}$; the canonical
calibration is therefore the constant $1$-vector $|e|$ of the integral
affine structure in which $v(T)$ is an affine coordinate. The total lattice
length of $\Sk(E)$ is $\lambda=v(q)$.

\emph{Consistency with Parts I and II.} The circle of circumference $v(q)$
is exactly the dual complex of the standard semistable model of $E$ computed
in \Cref{ex:tate-degeneration}, with the integral affine structure of
\Cref{def:integral-affine}; the retraction $E\to\Sk(E)$ is the skeletal
retraction of \Cref{prop:skeleton-retraction}; and the total mass computation
of \Cref{ex:ch08-tate-curve} uses the same normalization, the lattice length
of an edge being the valuation of the corresponding node parameter, as
recorded in \Cref{ex:ch08-An-edge}. There is no discrepancy factor between
the model-theoretic and tropical normalizations here, and that is a
consequence of the weight being $1$: for a curve with a node whose local
equation is $xy=\pi^m$ the corresponding edge has lattice length $m$, and
the extra factor is absorbed by the integral affine structure, not by the
calibration. See \Cref{rem:ch08-base-change-lengths} for the base-change
bookkeeping.
\end{example}

\subsection{A multiplicity different from one}

\begin{example}[A weight-two ray, computed from graded reduction]
\label{ex:trop-multiplicity-two}
Assume $\operatorname{char}\widetilde K\ne2$ and let
\[
 Z=\{1+T_1+T_2^2=0\}\subset\mathbf G_{m,K}^2,
 \qquad X=Z^{\An},
\]
of level $1$, with moment map $f=(T_1,T_2)|_Z$. Put $s=T_2|_Z$, so
$T_1|_Z=-1-s^2$ and $Z\simeq\{s:s\ne0,\ 1+s^2\ne0\}$.

\emph{The tropical image.} With $u=(v(-1-s^2),v(s))$ one finds, exactly as
in \Cref{ex:trop-line-in-torus}:
\[
 \Trop_f(X)=\mathbf R_{\ge0}(1,0)\ \cup\ \mathbf R_{\ge0}(0,1)
 \ \cup\ \mathbf R_{\ge0}(-2,-1).
\]
Indeed $v(s)>0$ gives $v(1+s^2)=0$; $v(s)<0$ gives $v(1+s^2)=2v(s)$; and
$v(s)=0$ gives $v(1+s^2)\ge0$ arbitrary.

\emph{Balancing forces a weight two.} The primitive outgoing directions are
$(1,0)$, $(0,1)$, $(-2,-1)$. The only positive solution of
$m_1(1,0)+m_2(0,1)+m_3(-2,-1)=0$ up to scale is
$(m_1,m_2,m_3)=(2,1,1)$. So if the tropicalization is balanced --- which it
is, by \Cref{thm:analytic-tropical-balancing} --- the horizontal ray must
carry weight $2$. We verify this from the definition, twice, over two
different base fields.

\emph{The rays of weight one.} On $\mathbf R_{\ge0}(0,1)$ take
$q=\mathrm{pr}_2$: then $q\circ f=s$ is an open immersion, of degree $1$, and
$\sigma_q(a)=(0,a)$, so the weight is $1$. On
$\mathbf R_{\ge0}(-2,-1)$ take $q=\mathrm{pr}_2$ again: the section with
image the line $\{(2a,a)\}$ is $\sigma_q(a)=(2a,a)$, giving
$\sigma_{q,*}(e)=|(2,1)|$, the primitive volume vector of the tangent
lattice; the degree is again $1$, so the weight is $1$.

\emph{The horizontal ray over an algebraically closed field.} Assume first
$K$ algebraically closed and let $\rho\in(0,1)$. The fibre of
$f_{\mathrm{trop}}$ over $(-\log\rho,0)$ meets $\Sigma_f$ in the points $x$
with $|s(x)|=1$, $|(1+s^2)(x)|=\rho$, and $\widetilde{T_1(x)}$ transcendental
over $\widetilde K$. Factor $1+s^2=(s-\rmi)(s+\rmi)$ with $\rmi^2=-1$. Since
$|2\rmi|=1$ and $|(s+\rmi)-(s-\rmi)|=|2\rmi|$, the two factors cannot both
have absolute value $<1$; so either $|s-\rmi|=\rho$ and $|s+\rmi|=1$, or the
other way round. There are therefore exactly two points of $\Sigma_f$ over
$(-\log\rho,0)$, namely $\zeta_{\rmi,\rho}$ and $\zeta_{-\rmi,\rho}$, and
near each of them $-1-s^2$ is a unit times $(s\mp\rmi)$, so
$q\circ f=T_1|_Z$ has local degree $1$ at each. By
\Cref{thm:trop-multiplicity-formula} the weight of the ray is
\[
 d_D=1+1=2 .
\]

\emph{The horizontal ray over $\mathbf Q_p$.} Now take $K=\mathbf Q_p$ with
$p\equiv3\pmod4$, so that $-1$ is not a square in $\mathbf F_p$, and let
$\eta$ be the point of $\mathbf G_m^{\An}$ with $|T_1(\eta)|=\rho<1$ and
Gauss seminorm, that is $\eta=\eta_{-\log\rho}$. The fibre of
$q\circ f=T_1|_Z$ over $\eta$ is
\[
 \mathcal M\bigl(\mathcal H(\eta)[s]/(s^2+1+T_1)\bigr).
\]
Because $|T_1|<1$ and $p\ne2$, Hensel's lemma provides
$c\in\mathcal H(\eta)$ with $c^2=1+T_1$, so the algebra is
$\mathcal H(\eta)[\sqrt{-1}]$. The degree-zero residue field of
$\mathcal H(\eta)$ is $\mathbf F_p$ when the class of $\rho$ in
$\mathbf R_{>0}/|K^\times|$ has infinite order, and it is
$\mathbf F_p(\bar\theta)$ when $\rho^m\in|K^\times|$ for some $m\ge1$ (after
normalizing a suitable power of the coordinate). In both cases $-1$ is not a
square in it, because $\mathbf F_p$ is relatively algebraically closed in
$\mathbf F_p(\bar\theta)$; and a square
in $\mathcal H(\eta)$ of absolute value $1$ reduces to a square in the
residue field. So $\mathcal H(\eta)[\sqrt{-1}]$ is a field, the fibre is the
single point $x=\mathcal M(\mathcal H(x))$ with
$\mathcal H(x)=\mathcal H(\eta)(\sqrt{-1})$, and
\[
 \deg_x(q\circ f)=\dim_{\mathcal H(\eta)}\mathcal H(x)=2 .
\]
Now read off the three factors of \Cref{thm:trop-degree-graded}. The fibre is
reduced, so the length is $1$. The extension
$\mathcal H(x)/\mathcal H(\eta)$ is unramified with residue extension of
degree $2$ --- namely $\mathbf F_{p^2}/\mathbf F_p$ or
$\mathbf F_{p^2}(\bar\theta)/\mathbf F_p(\bar\theta)$ according to the two
cases above --- so the graded residue degree
$[\widetilde{\mathcal H(x)}:\widetilde{\mathcal H(\eta)}]$ is $2$. The
residue characteristic is $p\ne0$, but the extension is unramified, hence
defectless, so the defect is $1$. Therefore
\[
 \deg_x(q\circ f)=1\cdot2\cdot1=2,
\]
and the weight of the horizontal ray is again $2$, this time supplied
entirely by a graded residue degree rather than by a count of sheets.

\emph{Compatibility of the two computations.} The two answers must agree,
and they do so through
\Cref{prop:trop-calibration-base-change}: base change from $\mathbf Q_p$ to
$\widehat{\overline{\mathbf Q_p}}$ splits the single point $x$ of degree $2$
into two points of degree $1$, and pushforward adds them. This is the
cleanest illustration of the warning in
\Cref{sec:appG-analytic-multiplicity} that a tropical multiplicity is never
a cardinality: over $\mathbf Q_p$ the fibre has one point, over
$\mathbf C_p$ it has two, and the multiplicity is $2$ in both cases.
\end{example}

\begin{example}[A nonreduced source: where the length appears]
\label{ex:trop-length-factor}
Let $K$ be algebraically closed and let $X$ be the analytic space attached to
the nonreduced affine curve $\Spec K[T^{\pm1},\varepsilon]/(\varepsilon^2)$
--- that is, the double of $\mathbf G_m$ in the $\varepsilon$-direction ---
with the moment map $f=T$. Then $X_{\mathrm{red}}=\mathbf G_m^{\An}$, and
$\Sigma_f=\Sigma_{f_{\mathrm{red}}}=S(\mathbf G_m)$, since $\Sigma_f$ depends
only on the reduction (the graded residue field of a point of $X$ is that of
the corresponding point of $X_{\mathrm{red}}$). For $x=\eta_u$ the fibre of
$f$ is $\mathcal M\bigl(\mathcal H(x)[\varepsilon]/(\varepsilon^2)\bigr)$,
a single point whose local ring has length $2$ and residue field
$\mathcal H(x)$; so
\[
 \deg_x(f)=2=\underbrace{2}_{\text{length}}\cdot
 \underbrace{1}_{\text{graded residue degree}}\cdot
 \underbrace{1}_{\text{defect}},
\]
and the weight of $\Pi_f=\mathbf R$ is $2$. The calibrated degree of the
closed immersion $X_{\mathrm{red}}\hookrightarrow X$ is
$\tfrac12$ by \Cref{prop:trop-calibrated-degree} \cref{itm:trop-dc:red},
which is exactly what is needed for
$\Trop_f(X)=2\,\Trop_f(X_{\mathrm{red}})$ to be consistent with
\Cref{def:trop-cycle-calibration}, in which the cycle $[X]$ of the
nonreduced space is $2[X_{\mathrm{red}}]$.
\end{example}

\section{Interfaces with later extensions}

Ordinary CLD charts use tropicalizations by invertible analytic functions.
Two later enlargements begin from the present construction.
\begin{enumerate}
\item Delta-forms retain the same calibrated tropical cycles but allow
  coefficients supported on polyhedra. They contain the curvature of
  piecewise affine model metrics and support tropical intersection products.
\item Harmonic tropicalizations allow suitable harmonic coordinate functions,
  producing a larger sheaf with better Poincaré-lemma and cohomological
  behavior. They do not turn every polyhedral current into a smooth form.
\item Curve-test semipositivity studies restrictions to analytic curves and
  gives a local positivity criterion for PL functions. Its comparison with
  CLD-psh functions is a theorem with hypotheses, not a definition.
\end{enumerate}
These constructions are developed together in \Cref{chap:delta-forms} so that their
overlap and their different purposes remain visible.

\section{Dictionary established in this chapter}
\label{sec:trop-dictionary-closing}

The following correspondences are used without further comment from
\Cref{chap:cld-forms} onwards.
\[
\begin{array}{ccl}
\text{invertible functions }f_1,\ldots,f_r
 &\longleftrightarrow&
 \text{moment map }f:X\to\mathbf G_m^r,\\[1mm]
\text{monomial change of the }f_i
 &\longleftrightarrow&
 \text{integer-affine chart change},\\[1mm]
\text{graded independence of }\widetilde{f_i(x)}
 &\longleftrightarrow&
 x\in\Sigma_f,\\[1mm]
\text{Abhyankar point, }d_K(x)=\dim_xX
 &\longleftrightarrow&
 \text{point of some local skeleton},\\[1mm]
\text{Gauss point }\eta_u\text{ of a torus}
 &\longleftrightarrow&
 \text{point of }S(T),\ \trop(\eta_u)=u,\\[1mm]
\text{local degree of }q\circ f
 &\longleftrightarrow&
 \text{calibration }\mu_C,\\[1mm]
\text{length}\cdot\text{res.\ degree}\cdot\text{defect}
 &\longleftrightarrow&
 \deg_x(f),\\[1mm]
\text{finite analytic morphism}
 &\longleftrightarrow&
 \text{one degree factor in pushforward},\\[1mm]
\text{boundary }\partial X
 &\longleftrightarrow&
 \text{tropically thin, weakly analytic set},\\[1mm]
\text{polyhedral subspace of an analytic space}
 &\longleftrightarrow&
 \text{canonically calibrated tropical charts}.
\end{array}
\]
The third line is \Cref{thm:trop-sigma-characterization}. The fourth and
fifth lines are
\Cref{cor:trop-sigma-abhyankar} and \Cref{prop:trop-torus-skeleton}; the
sixth is \Cref{def:trop-canonical-calibration}; the seventh is
\Cref{thm:trop-degree-graded}; the eighth is
\Cref{thm:trop-finite-pushforward}; the ninth is
\Cref{prop:trop-boundary-thin}; and the last is
\Cref{cor:trop-calibration-glue}.

Four separations are equally part of the dictionary, and each of them is a
place where the literature is easy to misread.

First, \emph{the tropical image is not the local skeleton}. The image
$\Trop_f(X)$ lives in $N_{\mathbf R}$ and may have pieces of every dimension
$\le n$; the local skeleton $\Sigma_f$ lives in $X$, and only its
top-dimensional part $\Sigma_f^{(n)}$ carries a calibration. The two are
related by $\Pi_f=f_{\mathrm{trop}}(\Sigma_f^{(n)})=\Trop_f(X)^{(n)}$ in the
proper case (\Cref{thm:trop-dimension-sigma} \cref{itm:trop-dimsig:c}), and
not otherwise.

Second, \emph{the multiplicity is not a number of points}. It is built from
a local analytic degree, which over a nonalgebraically closed field can
contain a residue degree and over a nonreduced space can contain a length
(\Cref{ex:trop-multiplicity-two,ex:trop-length-factor}). Relative to the
canonical saturated tangent lattice, the scalar weight $m_D$ of the
fundamental cycle is a positive integer
(\Cref{cor:trop-weights-integral}). Fractional coefficients arise only after
changing to a noncanonical lattice normalization or allowing
rational-coefficient cycles; a nonreduced source by itself does not make the
canonical weight fractional.

Third, \emph{refinement is not symmetric}. Charts are compared by passing up
to a product, and a comparison of weights is therefore always a pushforward
statement, with the lattice index \eqref{eq:appG-lattice-index} present
exactly once.

Fourth, \emph{a skeleton of a model and a local skeleton of a moment map are
different constructions with a common intersection}. The former, from
\Cref{chap:skeleta}, is global, defined by a model over the valuation ring,
and comes with a retraction; the latter is local, defined by analytic
functions, and comes with a calibration. They agree in the cases computed in
\Cref{ex:trop-torus-explicit} and \Cref{ex:trop-tate-curve}, and the general
comparison --- model metrics against CLD currents --- is the subject of
\Cref{chap:compare-ma}, not of a definition.

\section*{Source notes}

\paragraph{Two lineages.}
The material of this chapter has two independent ancestries, and it is worth
keeping them apart because they answer different questions.

The \emph{Berkovich-analytic} lineage asks what a point of an analytic space
knows about a family of invertible functions. Its objects are completed
residue fields, graded reductions, and the invariants $d_K(x)$, $s(x)$,
$t(x)$. Berkovich introduced the spectrum and the analytification
\cite{Berk12}, together with the invariants and the local structure theory
that make $d_K(x)$ meaningful \cite{Berk93}. Temkin's graded reduction of a
germ \cite{Temkin2004}, and the associated Riemann--Zariski picture
\cite{Tem11}, are what turn the naive residue field into an object rich
enough to detect radii and rational ranks: \Cref{prop:trop-torus-skeleton}
and \Cref{thm:trop-sigma-characterization} are statements in that language
and have no formulation in terms of the ordinary residue field, as the
warning in \Cref{sec:appG-graded-reduction} records. Ducros supplied the
local structure of tropicalization germs (\Cref{thm:trop-local-germ}) and the
family reduction theory \cite{Duc21} used in the polyhedrality statements.
The theory of defect and Abhyankar extensions used in
\Cref{thm:trop-degree-graded} is classical valuation theory; Ostrowski's
lemma, in the form that the defect of a finite extension of complete
rank-one valued fields is a power of the residue characteristic, is in
\cite[Ch.~VI]{ZariskiSamuelII}, and the analytic packaging is
\cite[\S13.8]{CLD25}.

The \emph{tropical-geometric} lineage asks what a balanced weighted
polyhedral complex is, independently of any analytic space. Its objects are
lattices, primitive normal vectors, weights, balancing, stable intersection,
and pushforward with lattice indices. We follow Mikhalkin--Rau
\cite[\S\S1.3, 2.1, 4.2, 4.6 and 6.2]{MR18Tropical} for terminology and for
the lattice-index convention; in particular the coefficient in
\eqref{eq:analytic-tropical-pushforward} is the index appearing in their
formula~(4.9), with no additional degree factor. Equation
\eqref{eq:analytic-max-min-firewall} records the exact isomorphism between
their max-plus convention and the valuation convention used here. The
foundational results of this lineage --- balancing of tropicalizations of
subvarieties of tori, the structure theorem, the fundamental theorem of
tropical geometry --- were obtained independently of Berkovich theory, and
they remain valid statements about polyhedral complexes; what the
Berkovich-analytic route adds is that the weights have an intrinsic
definition, valid without strictness, without algebraic closure, and for
spaces that are not algebraizable. The warning in
\Cref{sec:appG-initial-tropical} marks the exact point at which the two
lineages must not be conflated: the initial-form calculus is bookkeeping,
whereas the identification of a corner locus with valuations of actual zeros
is a theorem with hypotheses.

\paragraph{Attributions in this chapter.}
Moment maps, tropical charts, characteristic polyhedra, canonical
calibrations, integration, and refinement independence follow the 2025
Chambert-Loir--Ducros manuscript \cite[\S\S12--15]{CLD25}, of which the 2012
version is \cite{CLD12}; the numbered correspondences are as follows.
\Cref{prop:trop-torus-skeleton} and \Cref{lem:trop-monomial-degree} are
\cite[Lemma~12.1.6 and Corollary~12.1.7]{CLD25};
\Cref{thm:trop-local-germ} is \cite[Proposition~12.2.8]{CLD25};
\Cref{thm:analytic-tropical-polyhedrality} is \cite[Theorem~12.2.4]{CLD25},
whose algebraizable case is due to Berkovich \cite[Corollary~6.2.2]{Berk93};
\Cref{thm:trop-sigma-characterization} and
\Cref{thm:trop-sigma-polyhedral} are \cite[\S\S12.6 and 12.9]{CLD25};
\Cref{thm:trop-dimension-sigma} is \cite[Theorem~12.8.2]{CLD25};
\Cref{prop:trop-degree-constant} is \cite[Proposition~14.2.2]{CLD25};
\Cref{thm:trop-degree-translation} is \cite[Theorem~14.2.7]{CLD25};
\Cref{thm:trop-calibration-canonical} is
\cite[Proposition~14.3.1]{CLD25};
\Cref{thm:trop-calibration-chart-independent} and
\Cref{cor:trop-calibration-glue} are \cite[\S14.3 and Lemma~14.3.4]{CLD25};
\Cref{thm:trop-multiplicity-formula} is \cite[\S14.5]{CLD25};
\Cref{thm:trop-finite-pushforward} is \cite[Proposition~14.4.2]{CLD25}; and
\Cref{prop:trop-weakly-analytic-image} together with
\Cref{prop:trop-boundary-thin} is \cite[\S11.10]{CLD25}, whose
affinoid-boundary case goes back to Ducros. Finally,
\Cref{thm:cld-partitions-unity} is
\cite[Proposition~15.2.8]{CLD25}. The proofs given here are our own
arrangement where supplied; those of
\Cref{thm:trop-sigma-characterization},
\Cref{cor:trop-sigma-abhyankar},
\Cref{thm:trop-dimension-formula},
\Cref{lem:trop-skeleton-attained},
\Cref{thm:trop-dimension-sigma},
\Cref{prop:trop-degree-constant},
\Cref{thm:trop-calibration-canonical},
\Cref{thm:trop-calibration-chart-independent},
\Cref{thm:trop-degree-graded} and all examples are complete, whereas
\Cref{thm:analytic-tropical-polyhedrality},
\Cref{thm:trop-local-germ},
\Cref{thm:trop-sigma-polyhedral},
\Cref{prop:trop-analytic-inputs} and
\Cref{thm:trop-degree-translation} are imported with exact references and a
description of the mechanism; \Cref{prop:trop-boundary-thin} and
\Cref{thm:cld-partitions-unity} are direct applications of the cited CLD
results.

\paragraph{Adjacent literature not used as a hypothesis.}
Several treatments of tropicalization on Berkovich spaces run parallel to
this one and are not needed here, but a reader will meet them. Gubler's
tropical analytic geometry and his work with K\"unnemann
\cite{GK17,GK19} develop the same calculus with an eye to Arakelov theory,
and are the source of the delta-form extension of \Cref{chap:delta-forms};
Gubler--Rabinoff--Werner relate skeleta of semistable models to
tropicalizations and prove that faithful tropicalizations exist, which is a
sharper statement than anything used above; Gubler--Jell--Rabinoff
\cite{GJR26} enlarge the supply of charts by allowing harmonic functions;
Jell \cite{Jell16} proves the Poincar\'e lemma in this calculus;
and Burgos Gil--Gubler--Jell--K\"unnemann \cite{BGJKK25} develop the toric
tropical pluripotential theory that \Cref{chap:cld-examples} compares with
ours. On the purely tropical side, the standard textbook treatments of
balancing and of the structure theorem are those of Maclagan--Sturmfels and
of Mikhalkin--Rau, the latter being our reference \cite{MR18Tropical}; the
fundamental theorem of tropical geometry in the form relating corner loci to
valuations of points is due to Kapranov in the hypersurface case and to
Einsiedler--Kapranov--Lind and Speyer--Sturmfels in general. Those three
sources are cited here in prose only, since they are not part of any proof
above.

\Cref{chap:delta-forms} cites Gubler--K\"unnemann for delta-forms,
Gubler--Jell--Rabinoff for harmonic tropicalizations, and
Gubler--Rabinoff for curve-test psh functions.
